\section{Multiplicities of inner forms and normalising constants}\label{sec:pf:multnc}

\providecommand{\multncAf}{\mathbb{A}_{f}}
\providecommand{\multncm}{m_{\mathrm{cusp}}}
\providecommand{\multncT}{\mathcal{T}}
\providecommand{\multncPhi}{\Phi}
\providecommand{\multncB}{\mathcal{B}}
\providecommand{\multncTy}{\mathfrak{T}}
\providecommand{\multncst}[1]{\textup{(#1)}}

This section proves two results of the formal development \cite{LeanSource2026} which are used by several other
parts of the argument. They are logically independent of each other.

\begin{enumerate}[label=(\Alph*)]
\item \emph{Equality of cuspidal multiplicities of inner forms} (\cref{pf:multnc:main-mult}). Let $Y_0$ and $W_i$ be
hermitian spaces of the same odd rank $r$ over a CM field $E$, isometric at all finite places, and let $\sigma$ be an
irreducible admissible representation of the finite-adelic unitary group $G_{Y_0}(\multncAf)$ which contains a
suitable type. Then the multiplicity of $\sigma\otimes\pi_\infty$ in the cuspidal spectrum of $G_{Y_0}$ equals the
multiplicity of $\sigma\circ\Phi^{-1}\otimes\pi'_\infty$ in the cuspidal spectrum of $G_{W_i}$, for every isomorphism
$\Phi$ of the finite-adelic groups induced by a finite-adelic isometry, where $\pi_\infty$, $\pi'_\infty$ are discrete
series with the same Harish-Chandra parameters. Although the result is called ``twisted multiplicity one'' in the
development, it is an \emph{equality} of multiplicities of two inner forms, not a multiplicity-one theorem in the sense of
\cite{AizenbudGourevitchRallisSchiffmann2010} (see \cref{rem:families:mult}).
\item \emph{The explicit normalising constant} (\cref{pf:multnc:main-nc}). For the bundle of counting data attached to a
Hida setting (the \emph{mass datum} of the row counts), a certain normalising constant equals an explicit constant built
from a mass constant of unimodular hermitian lattices, a power of the discriminant of $F$, special values
$L(i,\eta^i)$ of the quadratic character $\eta$ of $E/F$, and the mass of the row lattices. This is a unitary analogue
of Siegel's mass formula \cite{Siegel1935,GanHankeYu2001,Prasad1989}, adapted to the row lattices of the construction;
see the discussion after \cref{prop:families:mass}.
\end{enumerate}

\paragraph{Idea of (A).} The proof compares two simple trace formulas in the style of
\cite{DeligneKazhdanVigneras1984,Arthur1989}. One uses test functions $f=\bigotimes_vf_v$ whose archimedean components
are pseudo-coefficients of discrete series, whose components at a finite set $S$ of split places are built from types of
$\GL_r$ over the residue fields (these come from Green's parametrisation of cuspidal characters, \cref{def:green}), and
whose component at an auxiliary split place is the characteristic function of a regular Iwahori double coset. For such
$f$ the operator $R(f)$ on $L^2(G(F)\backslash G(\A))$ is of trace class. Its trace is a finite sum of multiplicities
times traces of $\sigma_i(f_f)$ (spectral side), and it equals the integral of the elliptic part of the kernel, a finite
sum of products of local orbital integrals (geometric side). At the real places the stable orbital integrals of the
pseudo-coefficients are given by the Weyl character formula; at the finite places the local orbital integrals are
preserved by local isometries. Hence the two trace formulas agree term by term, and a linear independence statement for
traces of irreducible representations containing the types (proved from Bernstein--Zelevinsky theory) separates the
multiplicities. The pseudo-coefficients are not quoted from \cite{ClozelDelorme1984}: they are constructed explicitly,
as polynomials in the matrix entries cut off along a radial coordinate, by a deformation argument
(\cref{pf:multnc:sec-pc}).

\paragraph{Idea of (B).} The statement is an identity between real numbers attached to a record, the mass datum
$\multncB$. After an algebraic reduction which uses the defining relations of the record, it becomes a product formula,
which is proved by combining (i) an archimedean identity (a disintegration constant times three archimedean volumes),
(ii) a splitting of the measure of the adelic orbit of a frame into local orbit measures, from an integral Witt theorem,
(iii) explicit local volumes at all finite places, (iv) a mass formula for the row lattices and (v) the mass ratio of
unimodular hermitian lattices.

\paragraph{Organisation.} \Cref{pf:multnc:sec-setting} fixes notation for unitary groups and multiplicities and states
the first theorem; \cref{pf:multnc:sec-red} reduces it to a comparison statement. \Cref{pf:multnc:sec-inputs} introduces
the test functions and the spectral and geometric inputs, and \cref{pf:multnc:sec-comp} proves the comparison from them.
\Cref{pf:multnc:sec-pc} constructs the pseudo-coefficients and \cref{pf:multnc:sec-sides} proves the spectral,
separation and geometric inputs; \cref{pf:multnc:sec-kpoly} proves the two analytic statements on which the
pseudo-coefficients rest (the $K$-polynomial and the limit formula), and \cref{pf:multnc:sec-classint} the statements on
class integrals and the elliptic expansion used on the geometric side. \Cref{pf:multnc:sec-nc} states the second theorem,
\cref{pf:multnc:sec-ncpf} proves it, \cref{pf:multnc:sec-ncin} proves the row-lattice mass formula and
\cref{pf:multnc:sec-ncproofs} the remaining analytic inputs (local orbit volumes, the archimedean constant, the Jacobian
identity, the real bracket and the splitting identity). \Cref{pf:multnc:sec-remarks} collects remarks on the formal
statements.

\paragraph{Conventions.} All representations are complex. Dimensions are cardinal numbers (in the formal statements,
ranks of modules); they are finite whenever we compare them as integers. Results of the development which are proved
elsewhere, or whose proofs are deferred, are cited by their Lean names with a statement of what we use; they are proved
in \cite{LeanSource2026}.

% ===========================================================================
\subsection{Unitary groups, discrete series models and cuspidal multiplicities}\label{pf:multnc:sec-setting}

In \cref{pf:multnc:sec-setting,pf:multnc:sec-red,pf:multnc:sec-inputs,pf:multnc:sec-comp,pf:multnc:sec-pc,pf:multnc:sec-sides}
we fix a global datum $D$ of the development. We use only the following parts of it: a totally real number field $F$, a
CM extension $E/F$ with non-trivial automorphism $c$, which is unramified at all finite places, a finite set $\multncT$
of archimedean labels, and the hermitian spaces of $D$ described below.

\begin{definition}[Hermitian spaces and unitary groups]\label{pf:multnc:def-groups}
A \emph{hermitian space} $V$ of rank $n$ over $E$ is given by a matrix $J_V\in M_n(E)$ with $c(J_V)^{t}=J_V$ and
$\det J_V\ne0$. For a commutative $F$-algebra $k$ put
\[
G_V(k)=\{g\in\GL_n(k\otimes_FE)\;:\;c(g)^{t}J_Vg=J_V\}.
\]
For a finite place $v$ of $F$ let $G_v=G_V(F_v)$, and let $K_v\subset G_v$ be the stabiliser of the completion at $v$
of the \emph{frame lattice} $L_V\subset E^n$; for an archimedean place $w$ let $G_w=G_V(F_w)$. The adelic group
$G_V(\A)$ is the restricted product of the $G_v$ with respect to the $K_v$, times $\prod_wG_w$, and $G_V(F)$ is
embedded diagonally. Let $G_f\subset G_V(\A)$ and $G_\infty\subset G_V(\A)$ be the subgroups of elements whose
archimedean, respectively finite, components are trivial, and $\mathrm{loc}_v\colon G_v\to G_f$ the inclusion at $v$.
\end{definition}

The set $\multncT$ comes with a bijection $\rho\mapsto w_\rho$ onto the set of archimedean places of $F$ and, for each
$\rho$, an embedding $\tau_\rho\colon E\to\C$ above $w_\rho$. Through $\tau_\rho$ the group $G_{w_\rho}$ is identified
with the unitary group of the hermitian matrix $\tau_\rho(J_V)\in M_n(\C)$. For a sign vector $\eps\in\{\pm1\}^n$ we
write $J_\eps=\mathrm{diag}(\eps)$, $U(\eps)=\{g\in\GL_n(\C):g^*J_\eps g=J_\eps\}$ (a group isomorphic to $\Uni(p,q)$,
where $p$ and $q=n-p$ are the numbers of signs $+1$ and $-1$) and $K(\eps)=U(\eps)\cap\Uni(n)\cong\Uni(p)\times\Uni(q)$.

\begin{definition}[Discrete series models]\label{pf:multnc:def-dsmodel}
Let $V$ be a hermitian space of rank $n$ and $S$ a family of parameter sets. A \emph{discrete series model} $M$ of $V$
with parameters $S$ consists, for each $\rho\in\multncT$, of a Hilbert space $H_\rho$, a sign vector $\eps_\rho$, a
vector $\lambda_\rho\in\R^n$ enumerating $S(\rho)$, a matrix $\mathrm{fr}_\rho\in\GL_n(\C)$ with
$\mathrm{fr}_\rho^*J_{\eps_\rho}\mathrm{fr}_\rho=\tau_\rho(J_V)$, and a map $\sigma_\rho$ from $M_n(\C)$ to bounded
operators on $H_\rho$ whose restriction to $U(\eps_\rho)$ is an irreducible square-integrable unitary representation
with Harish-Chandra parameter $\lambda_\rho$ (in the sense of (H2) in \cref{pf:multnc:hyp-published}). The \emph{discrete
series representation of $G_{w_\rho}$ attached to $M$} is $\pi_\rho(g)=\sigma_\rho(\mathrm{fr}_\rho\,g\,\mathrm{fr}_\rho^{-1})$.
\formal{Formalism.IndefR.Mult.DSModel}
\end{definition}

Fix a Haar measure $\mu$ on $G_V(\A)$ (it is also right invariant). Let $L^2=L^2(G_V(F)\backslash G_V(\A))$ for the
quotient measure, $R$ the right regular representation, and $L^2_{\mathrm{cusp}}\subset L^2$ the closure of the space of
continuous cuspidal functions.

\begin{definition}[Cuspidal multiplicities]\label{pf:multnc:def-mult}
Let $\sigma$ be a representation of $G_f$ on a complex vector space $W$ (no topology), and $M$ a discrete series model of
$V$. Let $\mathrm{Hom}(\sigma,M)$ be the space of $\C$-linear maps $\Psi$ from $W$ to the space of continuous
multilinear maps $\prod_{\rho}H_\rho\to L^2$ such that every $\Psi(w)$ takes values in $L^2_{\mathrm{cusp}}$,
$\Psi(\sigma(g)w)=R(g)\circ\Psi(w)$ for $g\in G_f$, and
$\Psi(w)(\dots,\pi_\rho(h)x_\rho,\dots)=R(h)\,\Psi(w)(\dots,x_\rho,\dots)$ for $h\in G_{w_\rho}$. The \emph{cuspidal
multiplicity} is $\multncm(\mu;\sigma,M)=\dim_\C\mathrm{Hom}(\sigma,M)$.
\formallabel{GlobalTheory.AutMult.mCuspDS}
\end{definition}

The statement of the main theorem uses a second version $m^{\mathrm{sm}}_\theta(\sigma,M)$ of the multiplicity, defined
with smooth vectors: it is the dimension of the space of $G_f$-equivariant linear maps from $W$ to the span of the
multilinear maps $T$ from $\prod_\rho H_\rho^{\infty}$ (smooth vectors) to functions on $G_V(\A)$ such that every $T(x)$
is a smooth automorphic function with respect to a fixed family $\theta$ of archimedean characters, is cuspidal, is
$G_{w_\rho}$-equivariant in the variable $x_\rho$, and is continuous for a Sobolev norm in each variable. The two
versions agree:

\begin{proposition}[{\lref{GlobalTheory.T.cuspidalMultiplicity_eq_discreteSeriesModel_maps}}]\label{pf:multnc:cuspmult}
Assume (H3), (H4) and (H6) of \cref{pf:multnc:hyp-published}. Let $\theta$ be an admissible family of archimedean
characters, $V$ a hermitian space, $M$ a discrete series model of $V$, $\mu$ a Haar measure on $G_V(\A)$ and $\sigma$ a
smooth representation of $G_f$. Then $m^{\mathrm{sm}}_\theta(\sigma,M)=\multncm(\mu;\sigma,M)$.
\end{proposition}

This expresses that cuspidal automorphic forms are rapidly decreasing and that the smooth vectors of an isotypic
component of $L^2_{\mathrm{cusp}}$ are automorphic forms; the formal proof (about $600$ Lean lines) is deferred.

\begin{hypothesis}[Published inputs]\label{pf:multnc:hyp-published}
The first main theorem assumes the following ten statements of the literature, in the formulation of
\cite{LeanSource2026}.
\begin{enumerate}[label=(H\arabic*)]
\item Analyticity: matrix coefficients of $K(\eps)$-finite vectors of admissible representations of $U(\eps)$ are
real-analytic; and a density statement for $K$-finite smooth subspaces of irreducible unitary representations of Levi
subgroups \cite{HarishChandra1953}.
\item Classification of discrete series: for every $n$ there are irreducible square-integrable unitary representations
$\pi(\eps,\lambda)$ of $U(\eps)$, for $\lambda\in\R^n$ strictly decreasing with $\lambda_a-\frac{n-1}2\in\Z$, with
infinitesimal character $\lambda$ and a unique lowest $K(\eps)$-type of highest weight $\lambda+\rho_n-\rho_c$ (all
other $K$-types differ from it by sums of positive roots); and every irreducible square-integrable unitary
representation of $U(\eps)$ is isomorphic, after a permutation of the coordinates, to some $\pi(\eps\circ s,\lambda)$
\cite{HarishChandra1966}.
\item Rapid decay: cuspidal automorphic functions which are smooth at infinity and of uniform moderate growth are rapidly
decreasing.
\item Finiteness: for an open subgroup $U\subset G_V(\A)$, given $K_\infty$-types and eigenvalues of the archimedean
Casimir operators, the space of automorphic forms of level $U$ with these data is finite-dimensional
\cite{HarishChandra1968}.
\item Trace class: for $f\in C_c^\infty(G_V(\A))$ the operator $R(f)$ on $L^2_{\mathrm{cusp}}$ is of trace class.
\item Dixmier--Malliavin factorisation: every $f\in C_c^\infty(G)$, $G$ a closed subgroup of some $\GL_n(\C)$, is a finite
sum of convolutions $g_i*h_i$ with $g_i,h_i\in C_c^\infty(G)$, $\supp g_i$ in a prescribed neighbourhood of $1$ and
$\supp h_i\subset\supp f$ \cite{DixmierMalliavin1978}.
\item Flath factorisation: irreducible unitary representations of $G_V(\A)$ are restricted tensor products of local ones
\cite{Flath1979}.
\item Zelevinsky's theorem: over a non-archimedean local field the Grothendieck ring of the groups $\GL_n$ is the
polynomial ring on the segments \cite{Zelevinsky1980}.
\item Geometric lemma: the filtration of Jacquet modules of parabolically induced representations of Bernstein and
Zelevinsky and the corresponding cell semisimplification \cite{BernsteinZelevinsky1977}.
\item Green's parametrisation: the statement $\sfname{Green}$ of \cref{def:green}, which holds by \cref{lit:top:green}.
\end{enumerate}
\formal{GlobalTheory.Published.HarishChandraAnalyticity}
\formal{GlobalTheory.Published.DiscreteSeriesClassification}
\formal{GlobalTheory.Published.CuspRapidDecay}
\formal{GlobalTheory.Published.HarishChandraFinite}
\formal{GlobalTheory.Published.CuspTraceClass}
\formal{GlobalTheory.Published.DixmierMalliavinFactorisation}
\formal{GlobalTheory.Published.FlathFactorisation}
\formal{GlobalTheory.Published.ZelevinskyPolynomialRing}
\formal{GlobalTheory.Published.GeometricLemma}
\formal{GlobalTheory.Published.GreenCuspidal}
\end{hypothesis}

\begin{remark}\label{pf:multnc:rem-unused}
Hypotheses (H1) and (H7) are hypotheses of the main theorem but are not used in its proof (they are discarded in
\cref{pf:multnc:thm-assembly}). Hypothesis (H4) is used only through \cref{pf:multnc:cuspmult,pf:multnc:specfin}.
\end{remark}

\paragraph{The realisation data.} The main theorem concerns a realisation $R$ of the datum $D$ with a face $F$; these are
data of the row construction, defined elsewhere in the development. We need only the following parts of $R$: the base
hermitian space $Y_0$ and the spaces $W_i$ for $i<m$, all of rank $r$; the family $\theta$ of archimedean characters;
for each $i$ a discrete series model $M_{W,i}$ of $W_i$; matrices $\beta_{Y}\in\GL_r(E\otimes\multncAf)$ and
$\beta_{W,i}$ (level frames) and a finite set $P_g$ of places of $E$; a finite group $G$ with an isomorphism
$G\cong G_0=\prod_{\mathfrak P\in P_g}\GL_r(k_{\mathfrak P})$, where $k_{\mathfrak P}$ is the residue field; a set
$\mathcal X$ of characters of $G$; and the function $c_R\colon G\to\C$, $c_R(g)=$ the coefficient of $g$ in the central
idempotent of $\mathcal X$.

\begin{definition}[Good representations; finite-adelic isometries]\label{pf:multnc:def-good}
A representation $\sigma$ of $G_f(Y_0)$ on $W$ is \emph{good} if there is $x\in W\smallsetminus\{0\}$ which is fixed by
the level subgroup of $G_f(Y_0)$ attached to $\beta_Y$ and the trivial level datum, and a family $\kappa\colon G\to
G_f(Y_0)$ of lifts of the elements of $G_0$ to the parahoric subgroup defined by $\beta_Y$ with
$\sum_{g\in G}c_R(g)\,\sigma(\kappa(g))x=x$. A topological group isomorphism $\Phi\colon G_f(Y_0)\to G_f(W_i)$ is
\emph{induced by the finite-adelic isometry} $j\in M_{r}(E\otimes\multncAf)$ if $c(j)^{t}J_{W_i}j=J_{Y_0}$,
$\Phi(g)j=jg$ for all $g$, and $j\beta_Y=\beta_{W,i}u$ for some $u\in\GL_r(\OO_E\otimes\widehat\Z)$ whose image in
$\GL_r(k_{\mathfrak P})$ is trivial for every $\mathfrak P\in P_g$.
\formal{GlobalTheory.CuspMult.IsGoodRep}
\formal{GlobalTheory.CuspMult.IsFinIso}
\end{definition}

\begin{definition}[Equality of multiplicities]\label{pf:multnc:def-multeq}
$\sfname{MultEq}(R)$ is the statement: for every $i<m$, every discrete series model $M_0$ of $Y_0$ with the parameters
prescribed by $D$ at the base, every $j$ and $\Phi$ such that $\Phi$ is induced by the finite-adelic isometry $j$, and
every irreducible admissible good representation $\sigma$ of $G_f(Y_0)$,
\begin{equation}\label{pf:multnc:eq-multeq}
m^{\mathrm{sm}}_\theta(\sigma,M_0)=m^{\mathrm{sm}}_\theta(\sigma\circ\Phi^{-1},M_{W,i}).
\end{equation}
\formal{GlobalTheory.CuspMult.MultEq}
\end{definition}

\begin{definition}[The standing hypotheses]\label{pf:multnc:def-given}
$\sfname{Given}(R,F)$ consists of: the core clauses of $R$ along $F$ (in particular $\theta$ is an admissible family of
archimedean characters); the existence of size data for $D$ with a parameter choice; absolute irreducibility of the
characters in $\mathcal X$; the good-types condition and the guard condition on $R$; and the gap condition
\begin{equation}\label{pf:multnc:eq-gapR}
100\,r^4\le\lambda_\rho(n)-\lambda_\rho(n+1)\qquad(\rho\in\multncT,\ n+1<r)
\end{equation}
for the Harish-Chandra parameters $\lambda_\rho(0)>\lambda_\rho(1)>\dots$ of $D$.
\formal{GlobalTheory.LiftRows.Given}
\end{definition}

\begin{theorem}[Twisted multiplicity one]\label{pf:multnc:main-mult}
Assume (H1)--(H10). Let $R$ be a realisation of $D$ with face $F$, and assume the following conditions of the row
construction: the admissibility of the family at some level $s$; the weak characterisation of $R$ along $F$; the
existence condition at the base; the host hypothesis for $D$ (interlacing and gaps of the Harish-Chandra parameters);
the existence of size data with a parameter choice; the triviality of the host's class field condition
($\mathrm{cc}_K=\top$); the guard condition; and the good-types condition. Then $\sfname{MultEq}(R)$ holds.
\formallabel{GlobalTheory.T.twisted_multiplicity_one}
\formallabel{GlobalTheory.MultEqFinal.twisted_multiplicity_eq}
\end{theorem}

These conditions enter only through the following shared lemma.

\begin{lemma}[{\lref{GlobalTheory.GivenOfTop.given_of_top}}]\label{pf:multnc:given-of-top}
Under the hypotheses of \cref{pf:multnc:main-mult} (other than (H1)--(H10)), $\sfname{Given}(R,F)$ holds for the
underlying realisation and the face $F$.
\end{lemma}

Its proof only assembles fields: the core clauses come from the weak characterisation and the existence condition
(\lref{GlobalTheory.CoreOfWeak.core_of_weak}, which extracts them from the characterisation of $R$), the
irreducibility of the characters from the admissibility of the family, and the gap \eqref{pf:multnc:eq-gapR} is part of
the host hypothesis.

% ===========================================================================
\subsection{Reduction to a comparison statement}\label{pf:multnc:sec-red}

\begin{definition}[Types]\label{pf:multnc:def-types}
Let $V$ be a hermitian space of rank $n$. A \emph{type datum} $\multncTy$ for $V$ consists of a finite set $S$ of finite
places of $F$, for each $v\in S$ a splitting datum at $v$ (an $F$-embedding $E\to F_v$ and a basis identifying
$G_v\cong\GL_n(F_v)$ and $K_v\cong\GL_n(\OO_{F_v})$, so that reduction gives $\mathrm{red}_v\colon K_v\to\GL_n(k_v)$),
and an irreducible finite-dimensional representation $\tau_v$ of $\GL_n(k_v)$ on $W_v$. It is \emph{cuspidal} if
$S\ne\varnothing$ and for every $0<a<n$ there is $v\in S$ such that $\tau_v$ has no non-zero vector fixed by the
unipotent radical of the standard parabolic subgroup of type $(a,n-a)$. A representation $\sigma$ of $G_f$ on $W$
\emph{contains} $\multncTy$ if there is $x\ne0$ in $W$ which is fixed by $\mathrm{loc}_v(K_v)$ for $v\notin S$ and which,
for every $v\in S$, lies in the span of the images of the $K_v$-equivariant maps
$(\tau_v\circ\mathrm{red}_v,W_v)\to(\sigma\circ\mathrm{loc}_v,W)$.
\formallabel{GlobalTheory.TrMult.HasTypeVec}
\end{definition}

\begin{definition}[Lattice conjugations]\label{pf:multnc:def-latconj}
Let $V,V'$ be hermitian spaces of the same rank $n$. A topological isomorphism $j_v\colon G_v(V)\to G_v(V')$ is a
\emph{lattice conjugation} if there is an $F_v\otimes E$-linear isometry $a\colon(F_v\otimes E)^n\to(F_v\otimes E)^n$ from
$J_V$ to $J_{V'}$ which maps the completed frame lattice of $V$ onto that of $V'$ and satisfies $j_v(g)\,a(y)=a(gy)$.
\formallabel{GlobalTheory.IndefiniteRowSideFullHecke.PlusPolynomialRoots.IsLatticeConj}
\end{definition}

\begin{definition}[The comparison statement]\label{pf:multnc:def-trcusp}
For $\lambda\in\Z^r$ the \emph{strong gap condition} $\multncst{G$_{100}$}$ is: $100\,r^4\le\lambda_i-\lambda_{i+1}$ for all
$i$. The \emph{comparison statement} $\multncst{C}$ asserts: let $V$ be a hermitian space of odd rank $r$,
$\lambda_\rho\in\Z^r$ ($\rho\in\multncT$) satisfying $\multncst{G$_{100}$}$, $M$ a discrete series model of $V$ with
parameters $\lambda$, and $\multncTy$ a cuspidal type datum for $V$; let $V'$ be a hermitian space of rank $r$ and $M'$ a
discrete series model of $V'$ with parameters $\lambda$; let $j_v\colon G_v(V)\to G_v(V')$ be lattice conjugations for
all finite $v$ and $\Phi\colon G_f(V)\to G_f(V')$ a topological isomorphism with
$\Phi\circ\mathrm{loc}_v=\mathrm{loc}'_v\circ j_v$ for all $v$; and let $\mu,\mu'$ be Haar measures on $G_V(\A)$,
$G_{V'}(\A)$. Then for every irreducible admissible representation $\sigma$ of $G_f(V)$ containing $\multncTy$,
\begin{equation}\label{pf:multnc:eq-trcusp}
\multncm(\mu;\sigma,M)=\multncm(\mu';\sigma\circ\Phi^{-1},M').
\end{equation}
\formallabel{GlobalTheory.PsCo.Pc.GapG}
\formallabel{GlobalTheory.PsCo.Pc.TRCusp'}
\end{definition}

The following three statements are proved in the development; their proofs are deferred.

\begin{lemma}[{\lref{GlobalTheory.T.exists_typeVector_of_good_level_representation}}]\label{pf:multnc:typevec}
Assume (H10) and $\sfname{Given}(R,F)$. If $\sigma$ is an admissible good representation of $G_f(Y_0)$, then there is a
cuspidal type datum $\multncTy$ for $Y_0$ contained in $\sigma$.
\end{lemma}

\begin{lemma}[{\lref{GlobalTheory.T.exists_localIsometries_innerAut_adelicFrameIsometry}}]\label{pf:multnc:localiso}
Assume $\sfname{Given}(R,F)$. If $\Phi\colon G_f(Y_0)\to G_f(W_i)$ is induced by a finite-adelic isometry, then there are
lattice conjugations $j_v\colon G_v(Y_0)\to G_v(W_i)$ for all finite $v$, a topological isomorphism
$\Phi_0\colon G_f(Y_0)\to G_f(W_i)$ with $\Phi_0\circ\mathrm{loc}_v=\mathrm{loc}_v\circ j_v$ for all $v$, and $b\in
G_f(Y_0)$ with $\Phi(g)=\Phi_0(bgb^{-1})$ for all $g$.
\end{lemma}

\begin{lemma}\label{pf:multnc:params}
Assume $\sfname{Given}(R,F)$, let $i<m$ and let $M_0$ be a discrete series model of $Y_0$ with the base parameters of $D$.
Then $W_i$ and $Y_0$ both have rank $r$, $r$ is odd, and there are $\lambda_\rho\in\Z^r$ satisfying
$\multncst{G$_{100}$}$ such that the parameters of $M_0$ and of $M_{W,i}$ are both $\lambda_\rho$.
\formallabel{GlobalTheory.PsCo.Pc.params_match'}
\end{lemma}

\begin{proof}
The ranks of $Y_0$ and $W_i$ are $r$ by the rank conditions on the spaces of $D$, and $r$ is odd by the definition of the
ranks of $D$. The Harish-Chandra parameters $\lambda_\rho(n)$ of $D$ are integers; put
$\lambda_\rho=(\lambda_\rho(0),\dots,\lambda_\rho(r-1))$. Both models are enumerated by their parameter sets, which by
construction are the sets $\{\lambda_\rho(n):n<r\}$; since $n\mapsto\lambda_\rho(n)$ is strictly decreasing on $[0,r)$,
the enumerations are $\lambda_\rho$. Finally $\multncst{G$_{100}$}$ is \eqref{pf:multnc:eq-gapR}.
\end{proof}

\begin{proposition}\label{pf:multnc:multeq-of-trcusp}
Assume (H3), (H4), (H6), (H10) and the comparison statement $\multncst{C}$. Then $\sfname{Given}(R,F)$ implies
$\sfname{MultEq}(R)$.
\formallabel{GlobalTheory.PsCo.Pc.multEq_of_trCusp'}
\end{proposition}

\begin{proof}
Let $i,M_0,j,\Phi$ and $\sigma$ be as in \cref{pf:multnc:def-multeq}. By \cref{pf:multnc:typevec} there is a cuspidal
type datum $\multncTy$ for $Y_0$ contained in $\sigma$. By \cref{pf:multnc:localiso} there are lattice conjugations $j_v$,
an isomorphism $\Phi_0$ compatible with them and $b\in G_f(Y_0)$ with $\Phi(g)=\Phi_0(bgb^{-1})$. By \cref{pf:multnc:params}
there are parameters $\lambda_\rho$ satisfying $\multncst{G$_{100}$}$, common to $M_0$ and $M_{W,i}$, and $r$ is odd. Let
$\mu,\mu'$ be Haar measures on $G_{Y_0}(\A)$, $G_{W_i}(\A)$; they are right invariant, the adelic unitary groups being
unimodular (\lref{GlobalTheory.T.isMulRightInvariant_haar_adelicUnitaryGroup_hermitianSpace}, which states the right
invariance of Haar measures on $G_V(\A)$). The representation $\sigma\circ\Phi^{-1}$ is admissible
(\lref{GlobalTheory.CuspMult.isAdmissible_comp_equiv}: composing with a topological isomorphism preserves
admissibility), and admissible representations are smooth. The family $\theta$ is admissible by the core clauses. By
\cref{pf:multnc:cuspmult} and by $\multncst{C}$ applied to $(Y_0,W_i,M_0,M_{W,i},\multncTy,(j_v),\Phi_0)$,
\[
m^{\mathrm{sm}}_\theta(\sigma,M_0)=\multncm(\mu;\sigma,M_0)=\multncm(\mu';\sigma\circ\Phi_0^{-1},M_{W,i}).
\]
The representations $\sigma\circ\Phi^{-1}$ and $\sigma\circ\Phi_0^{-1}$ are isomorphic: since
$\Phi^{-1}(h)=b^{-1}\Phi_0^{-1}(h)b$ for $h\in G_f(W_i)$, the operator $\sigma(b)$ intertwines them. Isomorphic
representations have the same cuspidal multiplicity, since composing with the isomorphism identifies the spaces
$\mathrm{Hom}$ of \cref{pf:multnc:def-mult}. Hence $\multncm(\mu';\sigma\circ\Phi_0^{-1},M_{W,i})=
\multncm(\mu';\sigma\circ\Phi^{-1},M_{W,i})=m^{\mathrm{sm}}_\theta(\sigma\circ\Phi^{-1},M_{W,i})$, the last equality by
\cref{pf:multnc:cuspmult}.
\end{proof}

\begin{proof}[Proof of \cref{pf:multnc:main-mult}]
By \cref{pf:multnc:given-of-top}, $\sfname{Given}(R,F)$ holds. By \cref{pf:multnc:thm-assembly} below, (H1)--(H10) imply
that $\sfname{Given}(R,F)$ implies $\sfname{MultEq}(R)$.
\end{proof}


% ===========================================================================
\subsection{Test functions and the inputs of the comparison}\label{pf:multnc:sec-inputs}

Throughout this subsection $V$ is a hermitian space of rank $n$.

\paragraph{Measures and traces.} For finite places $v$ let $\mu_v$ be Haar measures on $G_v$ with $\mu_v(K_v)=1$ for
almost all $v$, for archimedean $w$ let $\mu_w$ be Haar measures on $G_w$, and let $\mu$ be a Haar measure on
$G_V(\A)$. We say that $\mu$ is the \emph{restricted product} of the $\mu_w,\mu_v$ if
$\mu\bigl(\prod_wU_w\times\prod_{v\in S}U_v\times\prod_{v\notin S}K_v\bigr)=\prod_w\mu_w(U_w)\prod_{v\in S}\mu_v(U_v)
\prod_{v\notin S}\mu_v(K_v)$ for all finite $S$ and Borel sets $U_w,U_v$. A family $(f_v)_v$ is a \emph{finite family} if
$f_v\in C_c^\infty(G_v)$ for all $v$ and $f_v=\one_{K_v}$ for almost all $v$; we write $f_f(g)=\prod_vf_v(g_v)$ on $G_f$.
A Haar measure $\nu$ on $G_f$ is \emph{adapted} to $(\mu_v)$ if $\int_{G_f}f_f\,d\nu=\prod_v\int_{G_v}f_v\,d\mu_v$ for
every finite family. A bounded operator on a Hilbert space \emph{has trace $c$} if it is of trace class with trace $c$.
For a representation $\sigma$ of $G_f$ on a vector space $W$ (no topology), a function $\phi$ on $G_f$ and $c\in\C$, an
operator $A$ \emph{represents} $\sigma(\phi)$ if $\ell(Ax)=\int_{G_f}\phi(g)\,\ell(\sigma(g)x)\,d\nu(g)$ (the integrand
being integrable) for all $x\in W$ and $\ell\in W^*$; we say that \emph{$\sigma(\phi)$ has trace $c$} if it is represented
by an operator of finite rank with trace $c$, and that $\sigma(\phi)\ne0$ if the zero operator does not represent
$\sigma(\phi)$. For $f\in C_c^\infty(G_V(\A))$ let $R(f)=\int f(g)R(g)\,d\mu(g)$ on $L^2(G_V(F)\backslash G_V(\A))$.

\paragraph{Orbital integrals.} Let $H$ be a locally compact group with Haar measure $dh$, $\gamma\in H$ with centraliser
$H_\gamma$, $dt$ a Haar measure on $H_\gamma$ and $f\colon H\to\C$. The \emph{orbital integral} of $f$ at $\gamma$ is
$c$ if the function $x\mapsto f(x^{-1}\gamma x)$ on $H_\gamma\backslash H$ is integrable for the quotient measure
$dh/dt$ with integral $c$; we then write $O_\gamma(f)=c$. An element of a matrix group is \emph{regular semisimple} if
its characteristic polynomial is separable.

\begin{definition}[Pseudo-coefficients]\label{pf:multnc:def-pc}
Let $H$ be a locally compact group with a continuous injective homomorphism $\iota\colon H\to M_N(\C)$, a Haar measure
$dh$ and an irreducible unitary representation $\pi_0$. A function $f$ on $H$ is a \emph{pseudo-coefficient} of $\pi_0$ if
$f$ is smooth and compactly supported in the coordinates given by $\iota$, $\pi_0(f)$ has trace $1$, $\pi(f)$ has trace
$0$ for every irreducible unitary $\pi\not\cong\pi_0$, and $O_s(f)=0$ for every regular semisimple $s$ whose centraliser is
not compact (and every Haar measure on the centraliser). For $H=U(\eps)$, $\lambda\in\Z^r$ and $\pi_0=\pi(\eps,\lambda)$,
$f$ is a pseudo-coefficient \emph{with elliptic orbital integrals} if moreover $f$ is bi-$K(\eps)$-finite and, for every
regular semisimple $s$ with compact centraliser and every $y\in U(\eps)$, $\vartheta\in\R^r$ with
$y^{-1}sy=t(\vartheta):=\mathrm{diag}(e^{i\vartheta_1},\dots,e^{i\vartheta_r})$,
\begin{equation}\label{pf:multnc:eq-orbell}
O_s(f)=\overline{\Theta_{\eps,\lambda}(\vartheta)}\,/\,\mathrm{vol}(H_s),\qquad
\Theta_{\eps,\lambda}(\vartheta)=(-1)^{pq}\,\frac{\sum_{w\in W_K(\eps)}\mathrm{sgn}(w)\,
e^{i\langle\lambda\circ w^{-1},\vartheta\rangle}}{\prod_{a<b}\bigl(e^{i(\vartheta_a-\vartheta_b)/2}-
e^{-i(\vartheta_a-\vartheta_b)/2}\bigr)},
\end{equation}
where $W_K(\eps)=\{w\in\mathfrak S_r:\eps\circ w=\eps\}$ is the compact Weyl group ($\Theta_{\eps,\lambda}$ is
Harish-Chandra's formula for the character of $\pi(\eps,\lambda)$ on the compact Cartan subgroup \cite{HarishChandra1966}).
\formallabel{GlobalTheory.PsCo.Pc.IsPCG}
\formallabel{GlobalTheory.PsCo.Pc.IsPC'}
\end{definition}

\begin{definition}[Existence of pseudo-coefficients]\label{pf:multnc:def-localpc}
The statement $\multncst{PC}$ is: for every odd $r$, every $\eps\in\{\pm1\}^r$, every $\lambda\in\Z^r$ satisfying
$\multncst{G$_{100}$}$, every irreducible square-integrable unitary representation $\pi$ of $U(\eps)$ with Harish-Chandra
parameter $\lambda$ and every Haar measure on $U(\eps)$, $\pi$ has a pseudo-coefficient with elliptic orbital integrals.
\formallabel{GlobalTheory.PsCo.Pc.LocalPC}
\end{definition}

\begin{definition}[Test functions]\label{pf:multnc:def-testfn}
Let $\multncTy=(S,(\tau_v))$ be a type datum for $V$. An \emph{auxiliary place} $U$ for $V$ consists of a finite place
$u$ of $F$, a splitting datum at $u$ (so $G_u\cong\GL_n(F_u)$ and $K_u\cong\GL_n(\OO_{F_u})$; let $I_u$ be the Iwahori
subgroup) and a uniformiser $\varpi$ of $F_u$. A family $(f_v)_v$ is a \emph{finite test family} for $(\multncTy,U)$ and
$(\mu_v)$ if $u\notin S$ and
\begin{enumerate}[label=(\roman*)]
\item for $v\in S$, $f_v=e_v*h_v*e_v$ with $h_v\in C_c^\infty(G_v)$, where $e_v$ is the type idempotent of
$\tau_v\circ\mathrm{red}_v$ (a function supported on $K_v$) and $*$ is convolution for $\mu_v$;
\item $f_u=\one_{I_utI_u}$ with $t=\mathrm{diag}(\varpi^{a_1},\dots,\varpi^{a_n})$ for some strictly decreasing $a\in\Z^n$;
\item $f_v\in C_c^\infty(G_v)$ for $v\notin S\cup\{u\}$, and $f_v=\one_{K_v}$ for almost all $v$.
\end{enumerate}
It is \emph{spherical off the types} if $f_v$ is bi-$K_v$-invariant for $v\notin S\cup\{u\}$. Let $M$ be a discrete series
model of $V$ and $\mu_w$ Haar measures on the $G_w$. An \emph{archimedean test family} for $M$ is a family $(f_w)_w$ such
that $f_{w_\rho}$ is a pseudo-coefficient of $\pi_\rho$ on $G_{w_\rho}$ (with $\iota=\tau_\rho$) for every $\rho$. We
write $f=f_\infty\otimes f_f$, $f(g)=\prod_wf_w(g_w)\prod_vf_v(g_v)$.
\formallabel{GlobalTheory.Orb.IsTestFin}
\end{definition}

\begin{definition}[The archimedean orbital condition]\label{pf:multnc:def-archo}
The condition $\multncst{AO}$ for $(M,(\mu_w),(f_w))$ is: for every $\rho$, every regular semisimple $s\in G_{w_\rho}$ with
compact centraliser and every $g\in G_{w_\rho}$, $\vartheta\in\R^n$ with
$\mathrm{fr}_\rho\,g^{-1}sg\,\mathrm{fr}_\rho^{-1}=t(\vartheta)$, one has
$O_s(f_{w_\rho})=\overline{\Theta_{\eps_\rho,\lambda_\rho}(\vartheta)}/\mathrm{vol}(G_{w_\rho,s})$.
\formallabel{GlobalTheory.PsCo.Pc.ArchO}
\end{definition}

\begin{definition}[Test data]\label{pf:multnc:def-testdata}
The tuple $(V,\lambda,M,\multncTy,U,(\mu_w),(\mu_v),\mu,(f_w),(f_v))$ is a set of \emph{test data} if $n$ is odd,
$\lambda_\rho\in\Z^n$ satisfies the \emph{weak gap condition} $\lambda_\rho(i)-\lambda_\rho(i+1)\ge n$ for all $\rho,i$, the
parameters of $M$ are $\lambda$, $\multncTy$ is cuspidal, $\mu$ is the restricted product of the $\mu_w,\mu_v$, $(f_w)$ is
an archimedean test family for $M$ and $(f_v)$ a finite test family for $(\multncTy,U)$ and $(\mu_v)$.
\formallabel{GlobalTheory.IndefiniteRowSideFullHecke.PlusPolynomialRoots.IsTestDataX}
\end{definition}

\paragraph{Elliptic terms.} The geometric side is indexed by the characteristic polynomials of elliptic elements.

\begin{definition}[Plus polynomials, tori and classes]\label{pf:multnc:def-ppol}
A \emph{plus polynomial} of degree $n$ is a monic irreducible $P\in E[X]$ of degree $n$ with $P(0)\ne0$ and
$P=P^\sharp$, where $P^\sharp(X)=c(P(0))^{-1}X^n\,P^c(X^{-1})$. Let $K_P=E[X]/(P)$, $\alpha$ the class of $X$, and $*$
the involution of $K_P$ extending $c$ with $\alpha^*=\alpha^{-1}$. The torus $T_P$ is the unitary group of the hermitian
form $1$ on $K_P$: $T_P(k)=\{t\in(k\otimes_FK_P)^\times:tt^*=1\}$. For an $F$-algebra $k$ let $\mathrm{Cls}_P(k)$ be the
group of $*$-invariant units of $k\otimes_FK_P$ modulo the norms $yy^*$, and let $X_P(k,J_V)\subset\mathrm{Cls}_P(k)$ be the
set of classes of those $x$ for which the hermitian form $(y,y')\mapsto\mathrm{Tr}_{K_P/E}(x\,y\,y'^*)$ on
$k\otimes_FK_P$ is isometric to $(k\otimes_FE^n,J_V)$. An isometry $\phi$ identifies multiplication on $k\otimes_FK_P$
with an embedding $T_P(k)\to G_V(k)$; for each class $q$ we fix such a $\phi$ and put $\gamma_q=\phi(\alpha)$.
\formallabel{GlobalTheory.IndefiniteRowSideFullHecke.PlusPolynomialRoots.PPol}
\formallabel{GlobalTheory.IndefiniteRowSideFullHecke.PlusPolynomialRoots.XSet}
\end{definition}

\begin{definition}[Elliptic terms]\label{pf:multnc:def-terms}
For a place $v$ (finite or archimedean), Haar measures $\mu_v$ on $G_v$ and $dt_v$ on $T_P(F_v)$ (transported to the
centralisers of the $\gamma_q$) and a function $f_v$ on $G_v$, put
\begin{equation}\label{pf:multnc:eq-Phi}
\multncPhi_v(P;f_v)=\sum_{q\in X_P(F_v,J_V)}O_{\gamma_q}(f_v),
\end{equation}
and say that $\multncPhi_v(P;f_v)$ is \emph{well defined} if $X_P(F_v,J_V)$ is finite and all these orbital integrals
exist. For $\rho\in\multncT$ and $\lambda\in\Z^n$ put
\begin{equation}\label{pf:multnc:eq-PhiReal}
\multncPhi^{\R}_\rho(P;\lambda)=\begin{cases}\overline{\chi_{\lambda^\flat}(\alpha_1,\dots,\alpha_n)}\,/\,
\mathrm{vol}(T_P(F_{w_\rho}))&\text{if $T_P(F_{w_\rho})$ is compact,}\\ 0&\text{otherwise,}\end{cases}
\end{equation}
where $\alpha_1,\dots,\alpha_n$ are the roots of $\tau_\rho(P)$ in $\C$, $\lambda^\flat_j=\lambda_j-(\frac{n-1}2-j)$
($0\le j<n$), and $\chi_\nu(a)=\det(a_i^{\nu_j+n-1-j})/\det(a_i^{n-1-j})$ is the Weyl character of $\GL_n$ of highest
weight $\nu$ \cite{Weyl1946}. Given Haar measures $dt_P$ on $T_P(\A)$ which are restricted products of measures
$dt_{P,w},dt_{P,v}$, put
\begin{align}
\mathrm{t}(P)&=\mathrm{vol}(T_P(F)\backslash T_P(\A))\prod_v\multncPhi_v(P;f_v)\prod_{\rho\in\multncT}
\multncPhi^{\R}_\rho(P;\lambda_\rho),\label{pf:multnc:eq-term}\\
\mathrm{t}^{\mathrm{orb}}(P)&=\mathrm{vol}(T_P(F)\backslash T_P(\A))\prod_v\multncPhi_v(P;f_v)
\prod_{w}\multncPhi_w(P;f_w),\label{pf:multnc:eq-termorb}
\end{align}
the products over the finite places $v$ having almost all factors equal to $1$. We say that $\mathrm{t}(P)$ is
\emph{finitely defined} if the volume is finite, every $\multncPhi_v(P;f_v)$ is well defined and $\multncPhi_v(P;f_v)=1$ for
almost all $v$. Finally the \emph{elliptic kernel} of $f$ is the function on $G_V(F)\backslash G_V(\A)$
\begin{equation}\label{pf:multnc:eq-ellkernel}
K^{\mathrm{ell}}_f(x)=\sum_{\delta\in G_V(F),\ \mathrm{charpoly}(\delta)\ \text{irreducible over }E}f(x^{-1}\delta x).
\end{equation}
\formallabel{GlobalTheory.IndefiniteRowSideFullHecke.PlusPolynomialRoots.term}
\end{definition}

\paragraph{The inputs.} We now state the inputs of the comparison.

\begin{proposition}[{\lref{GlobalTheory.T.cuspidal_spectrum_finiteness_of_decay_harishChandra}}]\label{pf:multnc:specfin}
Assume (H3), (H4) and (H6). Let $V$ be a hermitian space, $M$ a discrete series model of $V$ and $\mu$ a Haar measure on
$G_V(\A)$. Then: (a) $\multncm(\mu;\sigma,M)$ is finite for every irreducible admissible $\sigma$; (b) for every open
subgroup $K\subset G_f$ there are finitely many pairwise non-isomorphic irreducible admissible representations
$\sigma_1,\dots,\sigma_N$ of $G_f$ such that every irreducible admissible $\sigma$ with $\multncm(\mu;\sigma,M)\ne0$ and a
non-zero $K$-fixed vector is isomorphic to some $\sigma_i$. We call (a) and (b) \emph{finiteness of the cuspidal
spectrum} $\multncst{Fin}$.
\end{proposition}

\begin{definition}[The spectral expansion]\label{pf:multnc:def-spectrace}
The statement $\multncst{Spec}$ is: for all test data, every Haar measure $\nu$ on $G_f$ adapted to $(\mu_v)$, and every
family $\sigma_1,\dots,\sigma_N$ of pairwise non-isomorphic irreducible admissible representations of $G_f$ such that every
irreducible admissible $\sigma$ with $\multncm(\mu;\sigma,M)\ne0$ and $\sigma(f_f)\ne0$ is isomorphic to some $\sigma_i$: if
$\sigma_i(f_f)$ has trace $c_i$ for all $i$, then $R(f)$ has trace $\sum_i\multncm(\mu;\sigma_i,M)\,c_i$.
\formallabel{GlobalTheory.TrMult.SpecTraceX}
\end{definition}

\begin{definition}[Separation]\label{pf:multnc:def-separation}
The statement $\multncst{Sep}$ is: for every type datum $\multncTy=(S,(\tau_v))$ for $V$, every auxiliary place $U$ with
$u\notin S$, all Haar measures $\mu_v$ with $\mu_v(K_v)=1$ for almost all $v$, every adapted $\nu$, every family
$\sigma_1,\dots,\sigma_N$ of pairwise non-isomorphic irreducible admissible representations of $G_f$ and every $a\in\C^N$:
if $\sum_ia_ic_i=0$ whenever $(f_v)$ is a finite test family for $(\multncTy,U)$ which is spherical off the types and
$\sigma_i(f_f)$ has trace $c_i$ for all $i$, then $a_i=0$ for every $i$ such that $\sigma_i$ contains $\multncTy$.
\formallabel{GlobalTheory.TrMult.Separation}
\end{definition}

\begin{definition}[The geometric side]\label{pf:multnc:def-ell}
The \emph{elliptic trace identity} $\multncst{Ell}$ is: for all test data, $R(f)$ has trace
$\int_{G_V(F)\backslash G_V(\A)}K^{\mathrm{ell}}_f(x)\,dx$. The \emph{elliptic expansion} $\multncst{EE}$ is: for all
test data and all choices of restricted product measures $dt_P$ on the tori, for every plus polynomial $P$ of degree $n$,
$\mathrm{t}(P)$ is finitely defined and $\multncPhi_w(P;f_w)$ is well defined for every archimedean $w$;
$\mathrm{t}^{\mathrm{orb}}(P)\ne0$ for only finitely many $P$; $K^{\mathrm{ell}}_f$ is integrable; and
\begin{equation}\label{pf:multnc:eq-ellpart}
\int_{G_V(F)\backslash G_V(\A)}K^{\mathrm{ell}}_f(x)\,dx=\sum_P\mathrm{t}^{\mathrm{orb}}(P).
\end{equation}
\formallabel{GlobalTheory.TrMult.EllTraceX}
\formallabel{GlobalTheory.IndefiniteRowSideFullHecke.PlusPolynomialRoots.EllipticPartX}
\end{definition}

\begin{definition}[Derived statements]\label{pf:multnc:def-derived}
The \emph{real stable identity} $\multncst{RS}$ is: for $V$ of odd rank $n$, strictly decreasing $\lambda_\rho\in\Z^n$, a
discrete series model $M$ with parameters $\lambda$, Haar measures $\mu_w$, an archimedean test family $(f_w)$ for $M$
satisfying $\multncst{AO}$, every plus polynomial $P$ of degree $n$, every $\rho$ and every Haar measure on
$T_P(F_{w_\rho})$, one has $\multncPhi_{w_\rho}(P;f_{w_\rho})=\multncPhi^{\R}_\rho(P;\lambda_\rho)$. The \emph{geometric
expansion} $\multncst{Geo}$ is: for all test data satisfying $\multncst{AO}$ and all restricted product measures on the
tori, every $\mathrm{t}(P)$ is finitely defined, $\mathrm{t}(P)\ne0$ for only finitely many $P$, and $R(f)$ has trace
$\sum_P\mathrm{t}(P)$. The \emph{trace equality} $\multncst{TE}$ is: given test data for $V$ satisfying $\multncst{AO}$, a
hermitian space $V'$ of rank $n$, a discrete series model $M'$ of $V'$ with parameters $\lambda$, lattice conjugations
$j_v\colon G_v(V)\to G_v(V')$, Haar measures $\mu'_w$ on the $G_w(V')$ with the restricted product $\mu'$ of the $\mu'_w$
and the $j_{v*}\mu_v$, and an archimedean test family $(f'_w)$ for $M'$ satisfying $\multncst{AO}$, there is $c\in\C$
such that $R(f)$ and $R(f')$ both have trace $c$, where $f'=f'_\infty\otimes(f_v\circ j_v^{-1})_v$.
\formallabel{GlobalTheory.PsCo.Pc.RealStableX}
\formallabel{GlobalTheory.PsCo.Pc.GeomExpansionX}
\formallabel{GlobalTheory.PsCo.Pc.TraceEqX}
\end{definition}

The inputs are proved in \cref{pf:multnc:sec-pc,pf:multnc:sec-sides}:

\begin{theorem}\label{pf:multnc:inputs}
\begin{enumerate}[label=(\alph*)]
\item $\multncst{PC}$ holds (\cref{pf:multnc:localpc}).
\item (H2) and (H5) imply $\multncst{Spec}$ (\cref{pf:multnc:spectrace}).
\item (H8) and (H9) imply $\multncst{Sep}$ (\cref{pf:multnc:separation}).
\item (H5) implies $\multncst{Ell}$ (\cref{pf:multnc:elltrace}).
\item $\multncst{EE}$ holds (\cref{pf:multnc:ellpart}).
\end{enumerate}
\end{theorem}


% ===========================================================================
\subsection{Proof of the comparison statement}\label{pf:multnc:sec-comp}

\subsubsection*{Real places: stable orbital integrals of pseudo-coefficients}

We use three facts about the real tori, proved in the development (proofs deferred). Let $V$ have rank $n$, $P$ a plus
polynomial of degree $n$, $\rho\in\multncT$ and $w=w_\rho$.

\begin{lemma}\label{pf:multnc:realtori}
\begin{enumerate}[label=(\alph*)]
\item For every class $q\in X_P(F_w,J_V)$ the element $\gamma_q$ is regular semisimple
(\lref{GlobalTheory.T.isRegularSemisimple_torusEllipticElement_realPlace}); the measure on its centraliser
transported from a Haar measure $dt$ on $T_P(F_w)$ is a right invariant Haar measure of total mass $dt(T_P(F_w))$, and the
centraliser is compact if and only if $T_P(F_w)$ is compact
(\lref{GlobalTheory.T.centralizer_torusEllipticElement_realPlace}).
\item If $T_P(F_w)$ is compact and $M$ is a discrete series model of $V$, then $X_P(F_w,J_V)$ is finite and there are
$y_q\in\mathfrak S_n$ and $\vartheta_q\in\R^n$ ($q\in X_P(F_w,J_V)$) such that $(s,q)\mapsto y_q^{-1}s$ is a bijection
$W_K(\eps_\rho)\times X_P(F_w,J_V)\to\mathfrak S_n$, $e^{i\vartheta_q(a)}=\alpha_{y_q^{-1}(a)}$ for all $a$ (where
$\alpha_1,\dots,\alpha_n$ are the roots of $\tau_\rho(P)$), and $\mathrm{fr}_\rho\,g^{-1}\gamma_qg\,\mathrm{fr}_\rho^{-1}=
t(\vartheta_q)$ for some $g\in G_w$ (\lref{GlobalTheory.T.compactRealPlace_torusClasses_signVectors}).
\item (Character identity) Let $n$ be odd, $\lambda\in\Z^n$, $\alpha\in\C^n$ with pairwise distinct entries, $I$ a finite
set and $y_q\in\mathfrak S_n$, $\vartheta_q\in\R^n$ ($q\in I$) such that $(s,q)\mapsto y_q^{-1}s$ is a bijection
$W_K(\eps)\times I\to\mathfrak S_n$ and $e^{i\vartheta_q(a)}=\alpha_{y_q^{-1}(a)}$. Then
$\sum_{q\in I}\Theta_{\eps,\lambda}(\vartheta_q)=\chi_{\lambda^\flat}(\alpha)$
(\lref{GlobalTheory.T.finsum_discreteSeriesCharacter_eq_weylCharacter}).
\end{enumerate}
\end{lemma}

Part (c) is the Weyl character formula for $\GL_n$ written as a sum over $W_K(\eps)\backslash\mathfrak S_n$: the numerators
of the $\Theta_{\eps,\lambda}(\vartheta_q)$ together run over all permutations, and the Weyl denominators differ by the
signs of the permutations.

\begin{proposition}\label{pf:multnc:realstable}
The real stable identity $\multncst{RS}$ holds.
\formallabel{GlobalTheory.PsCo.Pc.realStableX}
\end{proposition}

\begin{proof}
Let $V,\lambda,M,(\mu_w),(f_w),P,\rho$ and $dt$ be as in \cref{pf:multnc:def-derived}, $w=w_\rho$, $\lambda=\lambda_\rho$
and $\eps=\eps_\rho$. Since $n$ is odd and $\lambda$ is integral and strictly decreasing, $\lambda$ is a Harish-Chandra
parameter ($\lambda_a-\frac{n-1}2\in\Z$).

Suppose first that $T_P(F_w)$ is compact, and take $y_q,\vartheta_q$ as in \cref{pf:multnc:realtori}(b). For
$q\in X_P(F_w,J_V)$, $\gamma_q$ is regular semisimple with compact centraliser of volume $dt(T_P(F_w))$ (part (a)) and is
conjugate to $t(\vartheta_q)$ through $\mathrm{fr}_\rho$; by $\multncst{AO}$,
\[
O_{\gamma_q}(f_w)=\overline{\Theta_{\eps,\lambda}(\vartheta_q)}\,/\,dt(T_P(F_w)).
\]
The roots of $\tau_\rho(P)$ are pairwise distinct since $P$ is irreducible, hence separable. Summing over the finite set
$X_P(F_w,J_V)$ and using \cref{pf:multnc:realtori}(c),
\[
\multncPhi_w(P;f_w)=\overline{\textstyle\sum_q\Theta_{\eps,\lambda}(\vartheta_q)}\,/\,dt(T_P(F_w))
=\overline{\chi_{\lambda^\flat}(\alpha)}\,/\,dt(T_P(F_w))=\multncPhi^{\R}_\rho(P;\lambda).
\]
If $T_P(F_w)$ is not compact, then for every $q$ the centraliser of the regular semisimple element $\gamma_q$ is not
compact (part (a)), so $O_{\gamma_q}(f_w)=0$ because $f_w$ is a pseudo-coefficient; hence
$\multncPhi_w(P;f_w)=0=\multncPhi^{\R}_\rho(P;\lambda)$.
\end{proof}

\begin{proposition}\label{pf:multnc:geomexp}
$\multncst{Ell}$, $\multncst{EE}$ and $\multncst{RS}$ imply $\multncst{Geo}$.
\formallabel{GlobalTheory.PsCo.Pc.geomExpansionX_of_parts}
\end{proposition}

\begin{proof}
Fix test data satisfying $\multncst{AO}$ and restricted product measures on the tori. By $\multncst{EE}$ every
$\mathrm{t}(P)$ is finitely defined, the support of $P\mapsto\mathrm{t}^{\mathrm{orb}}(P)$ is finite and
\eqref{pf:multnc:eq-ellpart} holds. The weak gap condition gives $\lambda_\rho(i)-\lambda_\rho(i+1)\ge n\ge1$, so the
parameters are strictly decreasing and $\multncst{RS}$ applies at every $\rho$:
$\multncPhi_{w_\rho}(P;f_{w_\rho})=\multncPhi^\R_\rho(P;\lambda_\rho)$. Reindexing the product over the archimedean places by
the bijection $\rho\mapsto w_\rho$ we get $\mathrm{t}^{\mathrm{orb}}(P)=\mathrm{t}(P)$ for all $P$. Hence $\mathrm{t}$ has
finite support and, by $\multncst{Ell}$, $R(f)$ has trace $\int K^{\mathrm{ell}}_f=\sum_P\mathrm{t}^{\mathrm{orb}}(P)
=\sum_P\mathrm{t}(P)$.
\end{proof}

\subsubsection*{Transport along local isometries}

\begin{lemma}\label{pf:multnc:transport}
Let $V,V'$ be hermitian spaces of the same rank and $j_v\colon G_v(V)\to G_v(V')$ lattice conjugations for all finite $v$.
\begin{enumerate}[label=(\alph*)]
\item $j_v(K_v)=K'_v$ for all $v$; and for every type datum $\multncTy$ and auxiliary place $U$ for $V$ there are a type
datum $\multncTy'$ and an auxiliary place $U'$ for $V'$ such that $\multncTy'$ is cuspidal if $\multncTy$ is, and
$(f_v\circ j_v^{-1})_v$ is a finite test family for $(\multncTy',U')$ and $(j_{v*}\mu_v)$ whenever $(f_v)$ is one for
$(\multncTy,U)$ and $(\mu_v)$ (\lref{GlobalTheory.T.testFunctions_transport_localLatticeIsometries}).
\item For a plus polynomial $P$, a Haar measure $\mu_v$ on $G_v(V)$, a Haar measure on $T_P(F_v)$ and a function $f_v$ such
that $\multncPhi_v(P;f_v)$ and $\multncPhi'_v(P;f_v\circ j_v^{-1})$ (computed in $G_v(V')$ with $j_{v*}\mu_v$) are well
defined, these two numbers are equal (\lref{GlobalTheory.T.stableOrbitalIntegral_preserved_latticeIsometry}).
\item If $\Phi\colon G_f(V)\to G_f(V')$ is a topological isomorphism with $\Phi\circ\mathrm{loc}_v=\mathrm{loc}'_v\circ j_v$,
then $(f_v\circ j_v^{-1})_f(\Phi(g))=f_f(g)$ for every family $(f_v)$, and $\Phi_*\nu$ is adapted to $(j_{v*}\mu_v)$ whenever
$\nu$ is adapted to $(\mu_v)$ (\lref{GlobalTheory.T.transport_pureTensors_productMeasure_finiteAdelic}).
\end{enumerate}
\end{lemma}

\begin{proposition}\label{pf:multnc:traceeq}
$\multncst{Geo}$ implies $\multncst{TE}$.
\formallabel{GlobalTheory.PsCo.Pc.traceEqX_of_geomExpansionX}
\end{proposition}

\begin{proof}
Let the data be as in the definition of $\multncst{TE}$. By \cref{pf:multnc:transport}(a) there are $\multncTy',U'$ such
that $\multncTy'$ is cuspidal and $f'_f:=(f_v\circ j_v^{-1})_v$ is a finite test family for $(\multncTy',U')$ and
$(j_{v*}\mu_v)$. The ranks are equal and the parameters $\lambda$ of $M'$ satisfy the same gap condition, so the primed
data are test data for $V'$, and they satisfy $\multncst{AO}$ by hypothesis. Choose, for every plus polynomial $P$,
restricted product measures on $T_P$ (they exist,
\lref{GlobalTheory.IndefiniteRowSideFullHecke.PlusPolynomialRoots.exists_isTorusMeas}, which provides a Haar
measure on $T_P(\A)$ that is the restricted product of local Haar measures); the torus $T_P$ depends only on $P$, so the same
measures serve for $V$ and $V'$. By $\multncst{Geo}$ for $V$ and for $V'$, $R(f)$ has trace $\sum_P\mathrm{t}_V(P)$ and
$R(f')$ has trace $\sum_P\mathrm{t}_{V'}(P)$, all terms being finitely defined. In \eqref{pf:multnc:eq-term} the volume
factor and the archimedean factors depend only on $P$, $\lambda$ and the torus measures, hence agree for $V$ and $V'$. For
every finite $v$ both $\multncPhi_v(P;f_v)$ and $\multncPhi'_v(P;f_v\circ j_v^{-1})$ are well defined, so they are equal by
\cref{pf:multnc:transport}(b). Hence $\mathrm{t}_V(P)=\mathrm{t}_{V'}(P)$ for all $P$, and $c=\sum_P\mathrm{t}_V(P)$ is the
common trace.
\end{proof}

\subsubsection*{Archimedean test functions}

\begin{lemma}\label{pf:multnc:testarch}
Assume $\multncst{PC}$. Let $V$ have odd rank $n$, let $M$ be a discrete series model of $V$ with integral parameters
$\lambda_\rho$ satisfying $\multncst{G$_{100}$}$, and let $\mu_w$ be Haar measures on the $G_w$. Then there is an
archimedean test family $(f_w)$ for $M$ satisfying $\multncst{AO}$.
\formallabel{GlobalTheory.PsCo.Pc.exists_testArch}
\formallabel{GlobalTheory.PsCo.Pc.ArchT.exists_testArch}
\end{lemma}

\begin{proof}
Fix $\rho$ and let $e_\rho\colon G_{w_\rho}\to U(\eps_\rho)$, $g\mapsto\mathrm{fr}_\rho g\,\mathrm{fr}_\rho^{-1}$; it is an
isomorphism of topological groups, and $\pi_\rho=\sigma_\rho\circ e_\rho$. By $\multncst{PC}$, applied to $\sigma_\rho$
and the Haar measure $e_{\rho*}\mu_{w_\rho}$, there is a pseudo-coefficient $\phi_\rho$ of $\sigma_\rho$ on $U(\eps_\rho)$
with elliptic orbital integrals. Put $f_{w_\rho}=\phi_\rho\circ e_\rho$. Transport of structure along $e_\rho$ shows that
$f_{w_\rho}$ is a pseudo-coefficient of $\pi_\rho$: the two matrix embeddings are conjugate by the invertible matrix
$\mathrm{fr}_\rho$, so smoothness with compact support, regular semisimplicity and compactness of centralisers are
preserved; the irreducible unitary representations of $G_{w_\rho}$ are exactly the $\pi\circ e_\rho$; and traces and orbital
integrals are unchanged. If $s\in G_{w_\rho}$ is regular semisimple with compact centraliser and
$\mathrm{fr}_\rho g^{-1}sg\,\mathrm{fr}_\rho^{-1}=t(\vartheta)$, then $e_\rho(g)^{-1}e_\rho(s)e_\rho(g)=t(\vartheta)$, and
\eqref{pf:multnc:eq-orbell} for $\phi_\rho$ gives $\multncst{AO}$ at $\rho$. Since $\rho\mapsto w_\rho$ is a bijection,
this defines $f_w$ for every archimedean $w$.
\end{proof}

\subsubsection*{The comparison}

We use the following auxiliary facts (proofs deferred).

\begin{lemma}\label{pf:multnc:aux}
Let $V$ be a hermitian space.
\begin{enumerate}[label=(\alph*)]
\item For every finite set $S$ of places there is an auxiliary place $U$ for $V$ with $u\notin S$
(\lref{GlobalTheory.T.exists_splitPlaceHermitianSpace_notMem_finset}).
\item If $\mu$ is a Haar measure on $G_V(\A)$ and $\mu_v$ are Haar measures with $\mu_v(K_v)=1$ for almost all $v$, there
are Haar measures $\mu_w$ such that $\mu$ is the restricted product of the $\mu_w,\mu_v$, and a Haar measure $\nu$ on $G_f$
adapted to $(\mu_v)$ (\lref{GlobalTheory.T.exists_restrictedProduct_haar_given_finiteFactors}).
\item For a type datum $\multncTy$ and an auxiliary place $U$ there is an open compact subgroup $K\subset G_f$ such that,
for all Haar measures $\mu_v$, every finite test family $(f_v)$ for $(\multncTy,U)$ which is spherical off the types is a
finite family and satisfies $f_f(kg)=f_f(g)$ for $k\in K$
(\lref{GlobalTheory.T.exists_openCompact_level_sphericalOffTypes}).
\item Let $\nu$ be a Haar measure on $G_f$, $(f_v)$ a finite family and $\sigma$ an admissible representation of $G_f$.
Then $\sigma(f_f)$ has a trace (it is represented by an operator of finite rank); and if $f_f$ is left invariant under an
open compact subgroup $K$ and $\sigma(f_f)\ne0$, then $\sigma$ has a non-zero $K$-fixed vector
(\lref{GlobalTheory.T.admissibleRep_hasTrace_finiteAdelicPureTensorFunction}).
\item If $\Phi\colon G\to G'$ is an isomorphism of topological groups and $\sigma$ a representation of $G$, then
$\sigma\circ\Phi^{-1}$ is irreducible (admissible) if $\sigma$ is; and if $f'\circ\Phi=f$, then an operator represents
$\sigma(f)$ for $\nu$ if and only if it represents $(\sigma\circ\Phi^{-1})(f')$ for $\Phi_*\nu$
(\lref{GlobalTheory.T.representation_transport_topGroupEquiv}).
\item Given finite families $(A_i)$, $(B_i)$ of representations of a group and a representation $\sigma$, there is a family
$\tau_1,\dots,\tau_m$ of pairwise non-isomorphic representations, each of which is one of the $A_i$, the $B_i$ or $\sigma$,
such that every $A_i$ and every $B_i$ is isomorphic to some $\tau_k$, and $\tau_{k_0}$ has every property of $\sigma$ for
some $k_0$ (\lref{GlobalTheory.TrMultEnd.exists_family}).
\end{enumerate}
\end{lemma}

\begin{theorem}\label{pf:multnc:trcusp-steps}
$\multncst{PC}$, $\multncst{Fin}$, $\multncst{Spec}$, $\multncst{Sep}$ and $\multncst{TE}$ imply the comparison
statement $\multncst{C}$.
\formallabel{GlobalTheory.PsCo.Pc.trCusp'_of_steps}
\end{theorem}

\begin{proof}
Let $V,\lambda,M,\multncTy,V',M',(j_v),\Phi,\mu,\mu'$ and $\sigma$ be as in \cref{pf:multnc:def-trcusp}; write
$m(\cdot)=\multncm(\mu;\cdot,M)$ and $m'(\cdot)=\multncm(\mu';\cdot,M')$, and let $S$ be the set of places of $\multncTy$.

\emph{Step 1: measures.} By \cref{pf:multnc:aux}(a) choose an auxiliary place $U$ for $V$ with $u\notin S$. Let $\mu_v$
be the Haar measure on $G_v$ with $\mu_v(K_v)=1$ (all $v$). By \cref{pf:multnc:aux}(b) there are $\mu_w$ such that $\mu$
is the restricted product, and an adapted $\nu$. By \cref{pf:multnc:transport}(a), $j_v(K_v)=K'_v$, so
$\mu'_v:=j_{v*}\mu_v$ satisfies $\mu'_v(K'_v)=1$, and \cref{pf:multnc:aux}(b) gives $\mu'_w$ such that $\mu'$ is the
restricted product of the $\mu'_w,\mu'_v$; by \cref{pf:multnc:transport}(c), $\nu':=\Phi_*\nu$ is adapted to $(\mu'_v)$ and
$(f_v\circ j_v^{-1})_f\circ\Phi=f_f$ for all families. Let $\multncTy',U'$ be as in \cref{pf:multnc:transport}(a). The
parameters of $M'$ are $\lambda$, they satisfy $\multncst{G$_{100}$}$, and $r$ is odd. By \cref{pf:multnc:testarch} there
are archimedean test families $(f_w)$ for $M$ and $(f'_w)$ for $M'$ satisfying $\multncst{AO}$.

\emph{Step 2: a finite set of representations.} Let $K\subset G_f(V)$ be as in \cref{pf:multnc:aux}(c), and $K'=\Phi(K)$,
an open compact subgroup of $G_f(V')$. By $\multncst{Fin}$(b) for $(V,M,\mu,K)$ and $(V',M',\mu',K')$ there are finite
families $(\sigma^{(1)}_i)$ of representations of $G_f(V)$ and $(\sigma^{(2)}_i)$ of $G_f(V')$, irreducible, admissible and
pairwise non-isomorphic, such that every irreducible admissible representation with non-zero multiplicity and a non-zero
$K$-fixed (respectively $K'$-fixed) vector is isomorphic to a member of the family. Put $B_i=\sigma^{(2)}_i\circ\Phi$; it is
irreducible and admissible (\cref{pf:multnc:aux}(e)). Apply \cref{pf:multnc:aux}(f) to $(\sigma^{(1)}_i)$, $(B_i)$ and
$\sigma$: we get pairwise non-isomorphic $\tau_1,\dots,\tau_m$, each irreducible and admissible, and $k_0$ such that
$m(\tau_{k_0})=m(\sigma)$, $m'(\tau_{k_0}\circ\Phi^{-1})=m'(\sigma\circ\Phi^{-1})$ and $\tau_{k_0}$ contains $\multncTy$.
Put $\tau'_k=\tau_k\circ\Phi^{-1}$; these are irreducible, admissible and pairwise non-isomorphic (an isomorphism
$\tau'_k\cong\tau'_l$ composed with $\Phi$ gives $\tau_k\cong\tau_l$).

\emph{Step 3: covering.} Let $(f_v)$ be a finite test family for $(\multncTy,U)$ which is spherical off the types. Then
$f_f$ is left $K$-invariant (\cref{pf:multnc:aux}(c)), and $f'_f:=(f_v\circ j_v^{-1})_f=f_f\circ\Phi^{-1}$ is left
$K'$-invariant. If $\rho$ is an irreducible admissible representation of $G_f(V)$ with $m(\rho)\ne0$ and
$\rho(f_f)\ne0$, then $\rho$ has a non-zero $K$-fixed vector (\cref{pf:multnc:aux}(d)), so $\rho\cong\sigma^{(1)}_i\cong
\tau_k$ for some $i,k$. Likewise, if $\rho'$ is irreducible admissible for $G_f(V')$ with $m'(\rho')\ne0$ and
$\rho'(f'_f)\ne0$, then $\rho'\cong\sigma^{(2)}_i=B_i\circ\Phi^{-1}\cong\tau_k\circ\Phi^{-1}=\tau'_k$ for some $i,k$.

\emph{Step 4: the trace identity.} By $\multncst{Fin}$(a) the multiplicities $m(\tau_k)$ and $m'(\tau'_k)$ are finite; put
$a_k=m(\tau_k)-m'(\tau'_k)\in\Z$. Let $(f_v)$ be a finite test family for $(\multncTy,U)$, spherical off the types, and
suppose that $\tau_k(f_f)$ has trace $c_k$ for every $k$ (for $\nu$). Then
$(V,\lambda,M,\multncTy,U,(\mu_w),(\mu_v),\mu,(f_w),(f_v))$ are test data (the strong gap condition implies the weak one,
since $r\le100r^4$), and so are the primed data with $f'_f$, by \cref{pf:multnc:transport}(a). By \cref{pf:multnc:aux}(e),
$\tau'_k(f'_f)$ has trace $c_k$ for $\nu'$. By Step~3 the covering hypotheses of $\multncst{Spec}$ hold on both sides, so
$R(f)$ has trace $\sum_km(\tau_k)c_k$ and $R(f')$ has trace $\sum_km'(\tau'_k)c_k$, where $f=f_\infty\otimes f_f$ and
$f'=f'_\infty\otimes f'_f$. By $\multncst{TE}$, $R(f)$ and $R(f')$ have a common trace. Since the trace of a trace class
operator is unique, $\sum_km(\tau_k)c_k=\sum_km'(\tau'_k)c_k$, that is
\begin{equation}\label{pf:multnc:eq-key}
\sum_ka_kc_k=0.
\end{equation}

\emph{Step 5: separation.} By \eqref{pf:multnc:eq-key} the hypothesis of $\multncst{Sep}$ holds for
$(\multncTy,U,(\mu_v),\nu,(\tau_k),(a_k))$. Since $\tau_{k_0}$ contains $\multncTy$, $a_{k_0}=0$, so
$m(\sigma)=m(\tau_{k_0})=m'(\tau'_{k_0})=m'(\sigma\circ\Phi^{-1})$ as integers. Both sides are finite cardinals
($\multncst{Fin}$(a), applied to $\sigma$ and to the irreducible admissible $\sigma\circ\Phi^{-1}$), which gives
\eqref{pf:multnc:eq-trcusp}.
\end{proof}

\begin{theorem}\label{pf:multnc:multeq-of-sides}
Assume $\multncst{PC}$, (H3), (H4), (H6), (H10), $\multncst{Spec}$, $\multncst{Sep}$, $\multncst{Ell}$ and
$\multncst{EE}$. Then for every realisation $R$ and face $F$, $\sfname{Given}(R,F)$ implies $\sfname{MultEq}(R)$.
\formallabel{GlobalTheory.PsCo.Pc.multEq_of_sides}
\end{theorem}

\begin{proof}
$\multncst{Fin}$ holds by \cref{pf:multnc:specfin}. $\multncst{Geo}$ holds by \cref{pf:multnc:geomexp,pf:multnc:realstable},
hence $\multncst{TE}$ by \cref{pf:multnc:traceeq}, hence $\multncst{C}$ by \cref{pf:multnc:trcusp-steps}. Conclude by
\cref{pf:multnc:multeq-of-trcusp}.
\end{proof}

\begin{theorem}\label{pf:multnc:thm-assembly}
Assume (H1)--(H10). Then for every realisation $R$ and face $F$, $\sfname{Given}(R,F)$ implies $\sfname{MultEq}(R)$.
\formallabel{UpperTheory.Lit.I5.unitary_theory_of_harishChandra_langlands_alt}
\end{theorem}

\begin{proof}
By \cref{pf:multnc:inputs} the inputs of \cref{pf:multnc:multeq-of-sides} hold: $\multncst{PC}$ and $\multncst{EE}$
unconditionally, $\multncst{Spec}$ from (H2) and (H5), $\multncst{Sep}$ from (H8) and (H9), $\multncst{Ell}$ from (H5).
The remaining hypotheses of \cref{pf:multnc:multeq-of-sides} are (H3), (H4), (H6) and (H10). Hypotheses (H1) and (H7) are
not used.
\end{proof}

\begin{corollary}\label{pf:multnc:trcusp-holds}
Under (H2)--(H6), (H8) and (H9), the comparison statement $\multncst{C}$ holds.
\end{corollary}

\begin{proof}
Combine \cref{pf:multnc:inputs}, \cref{pf:multnc:specfin}, \cref{pf:multnc:realstable,pf:multnc:geomexp,pf:multnc:traceeq}
and \cref{pf:multnc:trcusp-steps}.
\end{proof}


% ===========================================================================
\subsection{Pseudo-coefficients of discrete series}\label{pf:multnc:sec-pc}

In this subsection $r$ is odd, $\eps\in\{\pm1\}^r$ has $p$ signs $+1$ and $q$ signs $-1$, $K=K(\eps)$ carries its
normalised Haar measure $dk$, and $\mu$ is a Haar measure on $U(\eps)$. The pseudo-coefficients are constructed explicitly
as polynomials in the matrix entries, cut off along a radial coordinate, and their properties are obtained by a deformation
argument.

\begin{definition}[Shapes, profiles and Koszul data]\label{pf:multnc:def-shapes}
Let $\mathrm{Pol}_r=\C[x_{ab},\bar x_{ab}:a,b\le r]$ and, for $g\in M_r(\C)$, $\mathrm{ev}_g(P)=P(g,\bar g)$. Let
$\mathrm{Mo}(\eps)=\{(a,b):\eps_a=1,\ \eps_b=-1\}$, so $N:=\#\mathrm{Mo}(\eps)=pq$, and
\[
r_M(g)=\sum_{(a,b)\in\mathrm{Mo}(\eps)}|g_{ab}|^2,\qquad
d(\vartheta)=\prod_{(a,b)\in\mathrm{Mo}(\eps)}|e^{i\vartheta_a}-e^{i\vartheta_b}|^2 .
\]
A \emph{cut function} is a smooth $H\colon\R\to\R$ vanishing on some $[R,\infty)$. For a cut function $H$ and
$P\in\mathrm{Pol}_r$ let $b_{H,P}(g)=H(r_M(g))\,\mathrm{ev}_g(P)$. The \emph{shape space} $\mathrm{Sh}(\eps,d)$ is the span of
the $b_{H,P}$ with $\deg P\le d$. For $Z\in M_r(\C)$ let $\bar Z=-J_\eps Z^{t}J_\eps$ (so $\bar Z$ is the complex conjugate of
$Z$ if $Z$ lies in the Lie algebra $\mathfrak u(\eps)$ of $U(\eps)$), and let $\partial_Z$ be the $\C$-linear derivation of
$\mathrm{Pol}_r$ given on the matrices of variables $x=(x_{ab})$, $\bar x=(\bar x_{ab})$ by $\partial_Zx=xZ-Zx$ and
$\partial_Z\bar x=\bar x\bar Z-\bar Z\bar x$. Put $b^Z_{H,P}=b_{H',\,\partial_Z(r_P)P}+b_{H,\,\partial_ZP}$, where
$r_P=\sum_{(a,b)\in\mathrm{Mo}(\eps)}x_{ab}\bar x_{ab}$ (so $\mathrm{ev}_g(r_P)=r_M(g)$). The \emph{commutator space}
$\mathrm{Ad}(\eps)$ is the span of the $b^Z_{H,P}$ with $Z\in\mathfrak u(\eps)$; for such $Z$,
$b^Z_{H,P}(g)=\frac{d}{du}\big|_{u=0}b_{H,P}(e^{-uZ}ge^{uZ})$, and $b^Z_{H,P}\in\mathrm{Sh}(\eps,\deg P+2)$
(\lref{GlobalTheory.PsCo.Pc1.B.cutoffPolynomial_conjugation_deriv_mem_shapeSpan}, which verifies both facts by a direct
differentiation). A
\emph{profile} of order $N$ is a smooth $G\colon\R\to\R$ with $G=0$ on $[1,\infty)$, $G^{(N)}\ge0$ and $G^{(N)}(0)>0$. For a
profile $G$, a sequence $p=(p_k)$ in $\mathrm{Pol}_r$ and $t>0$ put
\[
\phi_{G,p,t}(g)=\sum_{k=0}^{N}\Bigl(\frac{d}{dy}\Bigr)^k\bigl[G(ty)\bigr]_{y=r_M(g)}\,\mathrm{ev}_g(p_k).
\]
For $x=(a,b)\in\mathrm{Mo}(\eps)$ let $\partial_x=\partial_{E_{ba}}$, where $E_{ba}$ is the elementary matrix; let
$\Omega_P=(\partial_x\bar x_y)_{x,y\in\mathrm{Mo}(\eps)}$ and, for $M\in\N$ and $P\in\mathrm{Pol}_r$,
$\mathrm{top}(M,P)=\det(\bar x)^M\,P\,\det\Omega_P$. A \emph{Koszul datum} is a pair of sequences $p_k$, $q_{x,k}$
($x\in\mathrm{Mo}(\eps)$) with $(k+1)p_{k+1}+r_Pp_k=\sum_x\bigl(\partial_xq_{x,k}+\partial_x(r_P)\,q_{x,k-1}\bigr)$ for all
$k\ge0$ ($q_{x,-1}=0$), $p_k=0$ for $k>N$ and $q_{x,k}=0$ for $k\ge N$. Finally $\mathrm{Blk}(\eps)$ is the subalgebra
generated by the $x_{ab}$ with $\eps_a=\eps_b$.
\formallabel{GlobalTheory.PsCo.Pc1.ShapeM}
\formallabel{GlobalTheory.PsCo.Pc1.AdSpanM}
\formallabel{GlobalTheory.PsCo.Pc1.famF}
\formallabel{GlobalTheory.PsCo.Pc1.IsKoszul}
\end{definition}

The construction uses twelve statements, all proved in the development. We prove (S5) below
(\cref{pf:multnc:homot}), and (S1) and (S6) in \cref{pf:multnc:sec-kpoly} (\cref{pf:multnc:kdata,pf:multnc:stlim}); the
proofs of the others are deferred. In each, $\pi$ ranges over irreducible unitary representations of $U(\eps)$.
\begin{enumerate}[label=(S\arabic*)]
\item ($K$-data; proved in \cref{pf:multnc:kdata}) If $\lambda\in\Z^r$ satisfies
$\multncst{G$_{100}$}$ and $\pi_0=\pi(\eps,\lambda)$, there are $M\in\N$ and $P\in\mathrm{Blk}(\eps)$ such that, with
$\mathrm{top}=\mathrm{top}(M,P)$: $k\mapsto\mathrm{ev}_k(\mathrm{top})$ is a class function on $K$;
$\mathrm{ev}_{t(\vartheta)}(\mathrm{top})=\overline{\Theta_{\eps,\lambda}(\vartheta)}\,d(\vartheta)$ whenever the
$e^{i\vartheta_a}$ are distinct; and the operator $B_\pi=\int_K\mathrm{ev}_k(\mathrm{top})\,\pi(k)\,dk$, whenever it has a
trace, has trace $1$ if $\pi=\pi_0$ and trace $0$ if $\pi\not\cong\pi_0$.
\item (Koszul resolution; \lref{GlobalTheory.PsCo.Pc1.A.stAlg}) For $P\in\mathrm{Blk}(\eps)$ and $M\in\N$ there is a
Koszul datum $(p,q)$ with $p_N=\mathrm{top}(M,P)$.
\item (Profiles; \lref{GlobalTheory.PsCo.Pc1.stProfile_holds}) Profiles of every order exist.
\item (Family; \lref{GlobalTheory.PsCo.Pc1.stFam}) For a profile $G$ of order $N$ and any $p$ there is $d$ with
$\deg p_k\le d$ for $k\le N$, $\phi_{G,p,t}\in\mathrm{Sh}(\eps,d)$ for $t>0$, and $\phi_{G,p,t}(g)=0$ whenever
$t\,r_M(g)>1$.
\item (Homotopy; proved in \cref{pf:multnc:homot}) If $(p,q)$ is a Koszul datum and $G$ a profile of order $N$, then
$\phi_{G,p,1}-\phi_{G,p,t}\in\mathrm{Ad}(\eps)$ for $t\ge1$.
\item (Limit; proved in \cref{pf:multnc:stlim}) For a profile $G$ of order $N$ there is $c\ne0$ such that for every
$p$: (i) $\int\phi_{G,p,t}\,w\,d\mu\to c\int_K\mathrm{ev}_k(p_N)\,w(k)\,dk$ as $t\to\infty$, for every continuous $w$; and
(ii) if $s=t(\vartheta)$ with distinct $e^{i\vartheta_a}$ and $\mathrm{ev}_{k^{-1}sk}(p_N)=\mathrm{ev}_s(p_N)$ for $k\in K$,
then $\int\phi_{G,p,t}(x^{-1}sx)\,d\mu(x)\to c\,\mathrm{ev}_s(p_N)/d(\vartheta)$.
\item (Regularity; \lref{GlobalTheory.ShapeSmoothFinite.stShape}) Elements of $\mathrm{Sh}(\eps,d)$ restrict to smooth,
compactly supported, bi-$K$-finite functions on $U(\eps)$.
\item (Orbital integrals of commutators; \lref{GlobalTheory.PsCo.Pc1.orbInt_adSpan_eq_zero}) Elements of
$\mathrm{Ad}(\eps)$ have orbital integral $0$ at every regular semisimple element.
\item (Elliptic orbital integrals; \lref{GlobalTheory.PsCo.Pc1.stOrbCpt}) If $s$ is regular semisimple with compact
centraliser, $y\in U(\eps)$ and $f$ is continuous with compact support, then
$O_s(f)=\int f(x^{-1}y^{-1}syx)\,d\mu(x)/\mathrm{vol}(U(\eps)_s)$.
\item (Non-elliptic elements; \lref{GlobalTheory.PsCo.NonEll.stNonEll}) If $s$ is regular semisimple with non-compact
centraliser, there is $R>0$ with $r_M(x^{-1}sx)\ge R$ for all $x$.
\item (Traces of shapes; \lref{GlobalTheory.PsCo.Pc3.stTr}) For every $\pi$ and $d$ there is a continuous $w$ on
$U(\eps)$ such that $\pi(\phi)$ has trace $\int\phi\,w\,d\mu$ for $\phi\in\mathrm{Sh}(\eps,d)$, and
$\int_K\mathrm{ev}_k(P)\pi(k)\,dk$ has trace $\int_K\mathrm{ev}_k(P)w(k)\,dk$ for $\deg P\le d$.
\item (Traces of commutators; \lref{GlobalTheory.PsCo.Pc3.stTrAd}) $\pi(\phi)$ has trace $0$ for every $\pi$ and
$\phi\in\mathrm{Ad}(\eps)$.
\end{enumerate}

Statements (S1) (the existence of a $K$-polynomial with the trace properties, about $12\,000$ Lean lines, using the Casimir
operator, the lowest $K$-type and the strong gap condition) and (S6) (a Laplace-type asymptotic expansion of
$\phi_{G,p,t}$ near $K$, about $1\,600$ lines) carry most of the analytic weight of the construction.

\begin{theorem}[Existence of pseudo-coefficients]\label{pf:multnc:localpc}
Statements (S1)--(S12) imply $\multncst{PC}$. Consequently $\multncst{PC}$ holds.
\formallabel{GlobalTheory.PsCo.Pc1.localPC_of_steps}
\formallabel{GlobalTheory.PsCo.Pc.localPC_holds}
\formallabel{GlobalTheory.T.localPseudoCoefficientExistence_unitaryGroup}
\end{theorem}

\begin{proof}
Let $\lambda\in\Z^r$ satisfy $\multncst{G$_{100}$}$ and $\pi_0=\pi(\eps,\lambda)$. Take $M,P$ from (S1), a Koszul datum
$(p,q)$ with $p_N=\mathrm{top}(M,P)$ from (S2), a profile $G$ of order $N$ from (S3), $d$ from (S4) and $c\ne0$ from (S6).
Put $\phi_t=c^{-1}\phi_{G,p,t}$ for $t\ge1$; then $\phi_t\in\mathrm{Sh}(\eps,d)$ and, by (S5), $\phi_1-\phi_t\in
\mathrm{Ad}(\eps)$. We show that $f=\phi_1|_{U(\eps)}$ is a pseudo-coefficient of $\pi_0$ with elliptic orbital integrals.
By (S7) it is smooth, compactly supported and bi-$K$-finite.

\emph{Traces.} Let $\pi$ be irreducible unitary, $B_\pi=\int_K\mathrm{ev}_k(p_N)\pi(k)\,dk$ and $w$ as in (S11) for $\pi$
and $d$; then $B_\pi$ has trace $a:=\int_K\mathrm{ev}_k(p_N)w(k)\,dk$. For $t\ge1$, $\pi(\phi_1-\phi_t)$ has trace
$\int(\phi_1-\phi_t)w\,d\mu$ by (S11) and trace $0$ by (S12), so $u(t):=\int\phi_tw\,d\mu$ is constant on $[1,\infty)$.
By (S6)(i), $u(t)\to a$. A function which is constant on $[1,\infty)$ and has a limit at $\infty$ equals that limit, so
$u(1)=a$, and $\pi(f)$ has trace $u(1)=a$ by (S11). By (S1), $a=1$ for $\pi=\pi_0$ and $a=0$ for $\pi\not\cong\pi_0$.

\emph{Non-elliptic orbital integrals.} Let $s$ be regular semisimple with non-compact centraliser, and $R$ as in (S10).
For $t=1+2/R$ and every $x$, $t\,r_M(x^{-1}sx)\ge(1+2/R)R=R+2>1$, so $\phi_t(x^{-1}sx)=0$ by (S4) and $O_s(\phi_t)=0$.
Since $O_s(\phi_1-\phi_t)=0$ by (S8), $O_s(f)=0$.

\emph{Elliptic orbital integrals.} Let $s$ be regular semisimple with compact centraliser and $y^{-1}sy=t(\vartheta)$.
Then $t(\vartheta)$ is regular semisimple, so the $e^{i\vartheta_a}$ are distinct, and $t(\vartheta)\in K$. By (S9) and
(S8), $O_s(f)=O_s(\phi_t)=\int\phi_t(x^{-1}t(\vartheta)x)\,d\mu(x)/\mathrm{vol}(U(\eps)_s)$ for every $t\ge1$, so the
right-hand side is constant in $t$. Since $\mathrm{ev}_k(p_N)$ is a class function on $K$, (S6)(ii) applies, and by (S1)
the limit of the right-hand side is $c^{-1}\cdot c\,\mathrm{ev}_{t(\vartheta)}(\mathrm{top})/(d(\vartheta)\,
\mathrm{vol}(U(\eps)_s))=\overline{\Theta_{\eps,\lambda}(\vartheta)}/\mathrm{vol}(U(\eps)_s)$. Hence
\eqref{pf:multnc:eq-orbell} holds.
\end{proof}

\subsubsection*{The homotopy}

The deformation $t\mapsto\phi_{G,p,t}$ stays in a fixed class modulo commutators because of the Koszul relations.

\begin{lemma}\label{pf:multnc:adbas}
For every $Z\in M_r(\C)$, every cut function $H$ and every $P\in\mathrm{Pol}_r$, $b^Z_{H,P}\in\mathrm{Ad}(\eps)$.
\formallabel{GlobalTheory.T.conjugationDerivCutoffPolynomial_mem_unitarySpan}
\formallabel{GlobalTheory.PsCo.Pc1.adBas_mem_adSpanM}
\end{lemma}

\begin{proof}
The matrix $\bar Z=-J_\eps Z^tJ_\eps$ depends $\C$-linearly on $Z$, hence so do $\partial_Z$ and $b^Z_{H,P}$. Every
$Z\in M_r(\C)$ is $A+iB$ with $A=\frac12(Z-J_\eps Z^*J_\eps)$ and $B=\frac1{2i}(Z+J_\eps Z^*J_\eps)$, and $A,B\in\mathfrak
u(\eps)$: using $J_\eps^2=1$, $A^*J_\eps+J_\eps A=\frac12(Z^*J_\eps-J_\eps Z+J_\eps Z-Z^*J_\eps)=0$, and similarly for $B$.
Hence $b^Z_{H,P}=b^A_{H,P}+i\,b^B_{H,P}\in\mathrm{Ad}(\eps)$.
\end{proof}

\begin{lemma}\label{pf:multnc:famderiv}
Let $G$ be smooth, $(p,q)$ a Koszul datum, $t>0$ and $g\in M_r(\C)$. Put $H_{k,t}(y)=\bigl(\frac{d}{dy}\bigr)^{k+1}[G(ty)]$.
Then
\[
\frac{d}{dt}\,\phi_{G,p,t}(g)=t^{-1}\sum_{x=(a,b)\in\mathrm{Mo}(\eps)}\ \sum_{k=0}^{N}b^{E_{ba}}_{H_{k,t},\,q_{x,k}}(g).
\]
\formallabel{GlobalTheory.T.scaledKoszulTestFamily_hasDerivAt_scale}
\formallabel{GlobalTheory.PsCo.HomotDeriv.famF_hasDerivAt}
\end{lemma}

\begin{proof}
Let $\varrho=r_M(g)$, $a_j=\bigl(\frac{d}{dy}\bigr)^j[G(ty)]_{y=\varrho}=t^jG^{(j)}(t\varrho)$ and $e_k=\mathrm{ev}_g(p_k)$.
Differentiating $s^kG^{(k)}(s\varrho)$ at $s=t$ gives $kt^{k-1}G^{(k)}(t\varrho)+t^k\varrho\,G^{(k+1)}(t\varrho)=
t^{-1}(ka_k+\varrho a_{k+1})$, so
\[
\frac{d}{dt}\phi_{G,p,t}(g)=t^{-1}\sum_{k=0}^N(ka_k+\varrho a_{k+1})e_k=t^{-1}\sum_{k=0}^Na_{k+1}\bigl((k+1)e_{k+1}+\varrho e_k\bigr),
\]
the second equality by summation by parts, using $e_{N+1}=0$ ($p_{N+1}=0$). Evaluating the Koszul relation at $g$ (note
$\mathrm{ev}_g(r_P)=\varrho$) gives $(k+1)e_{k+1}+\varrho e_k=\sum_x\bigl(\mathrm{ev}_g(\partial_xq_{x,k})+\mathrm{ev}_g(\partial_xr_P)\,
\mathrm{ev}_g(q_{x,k-1})\bigr)$. In the terms with $q_{x,k-1}$ we shift the summation index, using $q_{x,-1}=0$ and
$q_{x,N}=0$:
\[
\frac{d}{dt}\phi_{G,p,t}(g)=t^{-1}\sum_x\sum_{k=0}^N\Bigl(a_{k+1}\,\mathrm{ev}_g(\partial_xq_{x,k})+a_{k+2}\,
\mathrm{ev}_g(\partial_xr_P)\,\mathrm{ev}_g(q_{x,k})\Bigr).
\]
Since $H_{k,t}(\varrho)=a_{k+1}$ and $H'_{k,t}(\varrho)=a_{k+2}$, the summand is $b^{E_{ba}}_{H_{k,t},q_{x,k}}(g)$ by the
definition of $b^Z_{H,P}$.
\end{proof}

\begin{proposition}\label{pf:multnc:homot}
Statement (S5) holds: if $(p,q)$ is a Koszul datum and $G$ a profile of order $N$, then $\phi_{G,p,1}-\phi_{G,p,t}\in
\mathrm{Ad}(\eps)$ for every $t\ge1$.
\formallabel{GlobalTheory.PsCo.Pc1.Homot.stHomot}
\end{proposition}

\begin{proof}
For $k\ge0$ and $s,y\in\R$ put $S_k(s,y)=s^kG^{(k+1)}(sy)$, a smooth function with $\partial_yS_k=S_{k+1}$. Since
$H_{k,s}=s\,S_k(s,\cdot)$ and $H'_{k,s}=s\,S_{k+1}(s,\cdot)$, \cref{pf:multnc:famderiv} gives, for $s>0$,
\[
\frac{d}{ds}\phi_{G,p,s}(g)=\sum_x\sum_{k=0}^N\Bigl(S_{k+1}(s,\varrho)\,\mathrm{ev}_g(\partial_x(r_P)q_{x,k})+S_k(s,\varrho)\,
\mathrm{ev}_g(\partial_xq_{x,k})\Bigr),\qquad\varrho=r_M(g),
\]
a continuous function of $s$. Let $t\ge1$ and $H_k(y)=\int_1^tS_k(s,y)\,ds$. By differentiation under the integral sign,
$H_k$ is smooth with $H'_k(y)=\int_1^tS_{k+1}(s,y)\,ds$; and for $y\ge2$ and $s\in[1,t]$ we have $sy>1$, so all derivatives
of $G$ vanish at $sy$ and $H_k(y)=0$; thus $H_k$ is a cut function. Integrating the derivative from $1$ to $t$ (fundamental
theorem of calculus) gives, for every $g$,
\begin{align*}
\phi_{G,p,t}(g)-\phi_{G,p,1}(g)&=\sum_x\sum_{k=0}^N\Bigl(H'_k(\varrho)\,\mathrm{ev}_g(\partial_x(r_P)q_{x,k})+H_k(\varrho)\,
\mathrm{ev}_g(\partial_xq_{x,k})\Bigr)\\
&=\sum_x\sum_{k=0}^Nb^{E_{ba}}_{H_k,q_{x,k}}(g).
\end{align*}
By \cref{pf:multnc:adbas} each $b^{E_{ba}}_{H_k,q_{x,k}}$ lies in $\mathrm{Ad}(\eps)$, hence so does
$\phi_{G,p,1}-\phi_{G,p,t}$.
\end{proof}

% ===========================================================================
\subsection{The spectral side, separation, and the geometric side}\label{pf:multnc:sec-sides}

\subsubsection*{The spectral side}

\begin{theorem}\label{pf:multnc:spectrace}
(H2) and (H5) imply the spectral expansion $\multncst{Spec}$.
\formallabel{GlobalTheory.PsCo.SpecX.main}
\formallabel{GlobalTheory.PsCo.SpecX.specTrace}
\end{theorem}

\begin{proof}
Fix test data, an adapted $\nu$, and a family $(\sigma_i)$ with the covering property of \cref{pf:multnc:def-spectrace}
and traces $c_i$ of $\sigma_i(f_f)$. We use the following facts of the development (proofs deferred).
\begin{enumerate}[label=(\alph*)]
\item All constant terms of $f$ and of $f^*$ along proper parabolic subgroups vanish
(\lref{GlobalTheory.PsCo.Plug.parabolicT}: this holds for test functions with cuspidal types and the auxiliary place).
Hence, by (H5) and \lref{GlobalTheory.T.cuspidalOperators_vanishingRadicalIntegrals_unitaryGroup} (operators of
functions with vanishing constant terms are of trace class with image in $L^2_{\mathrm{cusp}}$), $R(f)$ is of trace class
and its trace is $\sum_\pi m(\pi)\,\mathrm{tr}\,\pi(f)$, a summable family indexed by the unitary dual of $G_V(\A)$, where
$m(\pi)$ is the multiplicity of $\pi$ in $L^2_{\mathrm{cusp}}$ (\lref{GlobalTheory.PsCo.Plug.cuspOpT}).
\item If $\pi$ is irreducible unitary, $m(\pi)\ne0$ and $\mathrm{tr}\,\pi(f)\ne0$, then $\pi$ is the completion of
$\sigma\otimes\bigotimes_\rho\pi_\rho$ for an irreducible admissible $\sigma$ of $G_f$ with $\sigma(f_f)\ne0$, and
$m(\pi)=\multncm(\mu;\sigma,M)$ (\lref{GlobalTheory.PsCo.SpecX.copyFinDS}, which uses (H2) and the vanishing of the
traces of pseudo-coefficients on all other unitary representations, and
\lref{GlobalTheory.T.cuspMultiplicity_completion_eq_discreteSeriesModel}, which compares the two multiplicities).
\item Conversely, if $\multncm(\mu;\sigma_i,M)\ne0$ there is an irreducible unitary $\pi_i$ which is the completion of
$\sigma_i\otimes\bigotimes_\rho\pi_\rho$ (\lref{GlobalTheory.T.exists_irreducibleUnitary_completion_finiteTimesDS}),
and $\pi_i(f)$ has trace $c_i\prod_\rho\mathrm{tr}\,\pi_\rho(f_{w_\rho})=c_i$ (\lref{GlobalTheory.PsCo.Plug.copyTraceT},
which factors the trace of a completed tensor product).
\item Two such completions are unitarily equivalent if and only if their finite parts are isomorphic
(\lref{GlobalTheory.T.unitaryEquiv_iff_finitePart_equiv_discreteSeries}).
\end{enumerate}
By (b) and the covering property every $\pi$ contributing to the sum in (a) is the completion attached to some $\sigma_i$; by
(c) and (d) the map $i\mapsto[\pi_i]$ (on the $i$ with non-zero multiplicity) is a bijection onto the contributing classes,
with $m(\pi_i)=\multncm(\mu;\sigma_i,M)$ and $\mathrm{tr}\,\pi_i(f)=c_i$. Rearranging the summable family along this
bijection (\lref{GlobalTheory.T.eq_sum_of_hasSum_of_bijective_matching}, an elementary statement on sums over matched
index sets) gives the trace $\sum_i\multncm(\mu;\sigma_i,M)\,c_i$.
\end{proof}

\subsubsection*{Separation}

\begin{theorem}\label{pf:multnc:separation}
(H8) and (H9) imply the separation statement $\multncst{Sep}$.
\formallabel{GlobalTheory.T.trace_separation_of_iwahori_trace}
\formallabel{GlobalTheory.TrMultSep.separation_of_iwahoriTrace}
\end{theorem}

\begin{proof}
Two local statements about $\GL_r$ over a non-archimedean local field are used. The \emph{Iwahori trace statement} (Iwahori
double cosets of dominant elements contract, and a trace formula for Iwahori--Hecke operators) holds
(\lref{GlobalTheory.IwahoriTraceProof.iwahoriTrace}). The \emph{separation of exponents} (for pairwise non-isomorphic
irreducible admissible representations with Iwahori-fixed vectors, the character sums of their exponents are linearly
independent on the strict cone of the torus) follows from (H8) and (H9) (\lref{GlobalTheory.SmGL.Exponents.exponentsSeparate}).
From these two statements one obtains:
(i) if $\pi_1,\dots,\pi_N$ are irreducible admissible representations of $G_u\cong\GL_r(F_u)$ and $\sum_id_ib_i=0$ whenever
$b_i$ is the trace of $\pi_i(\one_{I_utI_u})$, for all $t=\mathrm{diag}(\varpi^{a})$ with $a$ strictly decreasing, then
$\sum_id_ib_i=0$ whenever $b_i$ is the trace of $\pi_i(h)$ for a left $I_u$-invariant $h\in C_c^\infty(G_u)$
(\lref{GlobalTheory.T.admissible_irreducible_torus_characters}).

Now let $\multncTy,U,(\mu_v),\nu,(\sigma_i),a$ be as in \cref{pf:multnc:def-separation}. We use (proofs deferred):
(ii) there is a finite family $e$ such that, for every smooth $\sigma$ and every operator $P$ representing $\sigma(e_f)$,
$P$ is idempotent and $P\ne0$ if $\sigma$ contains $\multncTy$; and for every $g\in G_f$ there is a finite family $f^{g}$,
left $I_u$-invariant at $u$, such that the families obtained by replacing $f^{g}_u$ by $\one_{I_utI_u}$ ($t$ as above) are
finite test families spherical off the types, and $P\circ\sigma(g)\circ P$ represents $\sigma(f^{g}_f)$
(\lref{GlobalTheory.T.exists_testFamilies_projectedTranslationOperators});
(iii) every irreducible admissible $\sigma_i$ has an irreducible admissible local component $\pi_i$ at $u$
(\lref{GlobalTheory.T.exists_irreducible_admissible_localComponent_finiteAdelic}), and for every finite family $h$
there is $c'_i$ such that, for $h'_u\in C_c^\infty(G_u)$, if $\pi_i(h'_u)$ has trace $b$ then $\sigma_i$ of the family $h$
with $h_u$ replaced by $h'_u$ has trace $c'_ib$ (\lref{GlobalTheory.T.pure_tensor_trace_eq_const_mul_local_trace});
(iv) if $P_i$ are idempotents of finite rank on the spaces of pairwise non-isomorphic irreducible representations
$\sigma_i$ of a group $G$, and $\sum_ia_ic_i=0$ whenever $c_i$ is the trace of $P_i\sigma_i(g)P_i$ ($g\in G$), then $a_i=0$
for every $i$ with $P_i\ne0$ (\lref{GlobalTheory.T.coeff_eq_zero_of_sum_trace_projections_inequivalent_irreps}).

Let $P_i$ represent $\sigma_i(e_f)$; it exists and has finite rank by \cref{pf:multnc:aux}(d). By (iv) it suffices to show
$\sum_ia_ic_i=0$ when $c_i$ is the trace of $P_i\sigma_i(g)P_i$. Fix $g$, let $h=f^{g}$ and $c'_i$ as in (iii). For $t$ as
above let $b_i(t)$ be the trace of $\pi_i(\one_{I_utI_u})$. The family $h$ with $h_u$ replaced by $\one_{I_utI_u}$ is a
test family spherical off the types, and $\sigma_i$ of it has trace $c'_ib_i(t)$, so the hypothesis of $\multncst{Sep}$
gives $\sum_i(a_ic'_i)b_i(t)=0$. By (i), $\sum_ia_ic'_ib_i=0$, where $b_i$ is the trace of $\pi_i(h_u)$. By (iii) and (ii),
$c'_ib_i$ and $c_i$ are both traces of $\sigma_i(h_f)$, hence equal. So $\sum_ia_ic_i=0$, and $a_i=0$ whenever $\sigma_i$
contains $\multncTy$.
\end{proof}

\subsubsection*{The geometric side}

\begin{theorem}\label{pf:multnc:elltrace}
(H5) implies the elliptic trace identity $\multncst{Ell}$.
\formallabel{GlobalTheory.PsCo.Plug.ellipticSideDirect}
\formallabel{GlobalTheory.PsCo.Plug.ellTrace_of_middleX}
\end{theorem}

\begin{proof}
Fix test data and let $f=f_\infty\otimes f_f$; it is smooth with compact support. We use the following statements of the
development (proofs deferred).
\begin{enumerate}[label=(\alph*)]
\item All unipotent integrals of $f$ and of $f^*(g)=\overline{f(g^{-1})}$ along proper isotropic flags vanish
(\lref{GlobalTheory.PsCo.Plug.parabolicT}; this uses the cuspidality of the type datum).
\item For such $f$, $R(f)$ is of trace class and its trace is $\int_{G(F)\backslash G(\A)}K_f(x,x)\,dx$, where
$K_f(x,x)=\sum_{\gamma\in G(F)}f(x^{-1}\gamma x)$ (\lref{GlobalTheory.T.cuspidalOperators_vanishingRadicalIntegrals_unitaryGroup},
from (H5)).
\item Every element of $I_utI_u$, $t=\mathrm{diag}(\varpi^a)$ with $a$ strictly decreasing, has separable characteristic
polynomial (\lref{GlobalTheory.IwahoriEigenObj.iwahoriEigenvalues}); hence $f(x^{-1}\delta x)\ne0$ implies that
$\delta$ is regular semisimple (\lref{GlobalTheory.T.isRegularSemisimple_of_testFunction_ne_zero}).
\item Only finitely many $G(F)$-conjugacy classes of regular semisimple elements meet the support of
$x\mapsto f(x^{-1}\cdot x)$ (\cref{pf:multnc:regfin}, with the reduction-theory statement of
\lref{PositivityTheory.Prw.siegelCore_holds}).
\item For each regular semisimple $\gamma$ the class kernel $x\mapsto\sum_{\delta\sim\gamma}f(x^{-1}\delta x)$ is integrable
on $G(F)\backslash G(\A)$ (\cref{pf:multnc:regint}; $f$ is a cusp function by (a)).
\item The integral of the class kernel vanishes for every regular semisimple class whose characteristic polynomial is
reducible over $E$: for the classes of non-field type by \cref{pf:multnc:nonfield}, which applies since the type datum of
test data satisfies the radical hypothesis (\lref{GlobalTheory.T.typeHypothesis_forms_agree}), and for the remaining
classes by \cref{pf:multnc:hyperzero}.
\end{enumerate}
By (a) and (b), $R(f)$ has trace $\int K_f(x,x)\,dx$. By (c) the sum $K_f(x,x)$ runs over regular semisimple $\gamma$
only; group it by conjugacy classes. By (d) only finitely many classes contribute and by (e) each class kernel is
integrable, so the integral is the sum of the integrals of the class kernels
(\lref{GlobalTheory.T.integral_hasSum_eq_integral_indicator_hasSum'}, an interchange of a finite sum and an integral).
By (f) the classes with reducible characteristic polynomial contribute $0$; the remaining ones (irreducibility of the
characteristic polynomial is a property of the class) add up to $\int K^{\mathrm{ell}}_f$.
\end{proof}

\begin{theorem}\label{pf:multnc:ellpart}
The elliptic expansion $\multncst{EE}$ holds.
\formallabel{GlobalTheory.PsCo.Plug.ellipticT}
\formallabel{GlobalTheory.TrEllSum.ellipticPartX}
\end{theorem}

\begin{proof}
Of the test data only three properties are used: $\mu$ is the restricted product of the $\mu_w,\mu_v$, the $f_w$ are
smooth with compact support, and $(f_v)$ is a finite family. Under these assumptions, and since $E/F$ is a CM extension
unramified at all finite places and $J_V$ is hermitian non-degenerate, the conclusion is the elliptic expansion
\cref{pf:multnc:ellexp}, which states exactly the finiteness, integrability and expansion properties of $\multncst{EE}$
for such data.
\end{proof}


% ===========================================================================
\subsection{The $K$-polynomial and the limit formula}\label{pf:multnc:sec-kpoly}

We prove the statements (S1) and (S6) of \cref{pf:multnc:sec-pc}, keeping the notation of that subsection: $r$,
$\eps$ with $p$ signs $+1$ and $q$ signs $-1$, $K=K(\eps)$ with normalised Haar measure $dk$, $\mathrm{Mo}(\eps)$,
$N=pq$, $r_M$, $d(\vartheta)$, $\mathrm{Pol}_r$, $\mathrm{ev}$, $\mathrm{Blk}(\eps)$ and $\mathrm{top}(M,P)$.

\subsubsection*{The classical background}

Statement (S1) replaces a group of classical facts on discrete series. We recall them in the form in which they are
usually stated.

\begin{littheorem}[{\cite[Thm.~9.20 and Ch.~XII]{Knapp1986}, \cite{HarishChandra1966,Schmid1975,ClozelDelorme1984}}]\label{pf:multnc:lit-ds}
Let $G$ be a connected linear reductive Lie group, $K$ a maximal compact subgroup and $T\subset K$ a maximal torus which
is a Cartan subgroup of $G$; let $q_G=\frac12\dim G/K$.
\begin{enumerate}[label=(\alph*)]
\item (Harish-Chandra) For every nonsingular $\lambda\in i\mathfrak t^*$ such that $\lambda+\rho$ is analytically
integral there is a discrete series representation $\pi_\lambda$ with infinitesimal character $\lambda$, whose character
on the regular elements of $T$ is $(-1)^{q_G}\sum_{w\in W_K}\mathrm{sgn}(w)e^{w\lambda}/\prod_{\alpha>0}
(e^{\alpha/2}-e^{-\alpha/2})$; $\pi_\lambda\cong\pi_{\lambda'}$ if and only if $\lambda'\in W_K\lambda$; and every discrete
series representation is of this form.
\item (Schmid) $\pi_\lambda$ contains the $K$-type of highest weight $\lambda+\rho_n-\rho_c$ with multiplicity one and no
$K$-type which is lower in the sense of \cite{Schmid1975}; and $\pi_\lambda$ is characterised by its lowest $K$-type, up to
infinitesimal equivalence, among the irreducible representations of $G$.
\item (Clozel--Delorme) For every discrete series representation $\pi$ there is a $K$-finite $f\in C_c^\infty(G)$ with
$\mathrm{tr}\,\pi(f)=1$ and $\mathrm{tr}\,\pi'(f)=0$ for every irreducible tempered representation $\pi'\not\cong\pi$.
\end{enumerate}
\end{littheorem}

For $G=U(\eps)$ we have $q_G=pq$, and the character in (a) is the function $\Theta_{\eps,\lambda}$ of
\eqref{pf:multnc:eq-orbell}. Statement (S1) differs from these results in three respects. It does not assert the existence
of a pseudo-coefficient, but of a polynomial $\mathrm{top}(M,P)$ in the matrix entries whose restriction to $K$ is
$\pm\overline{\chi_T}\,D_\eps$, where $\chi_T$ is the character of a $K$-type and $D_\eps$ is the alternating sum of the
characters of $K$ on the exterior powers of $\C^{\mathrm{Mo}(\eps)}$, on which $k$ acts by the matrix
$(k_{aa'}\overline{k_{bb'}})$ (see \eqref{pf:multnc:eq-Deps} below); from it the
pseudo-coefficient is built in \cref{pf:multnc:localpc}. The operator $\int_K\mathrm{ev}_k(\mathrm{top})\pi(k)\,dk$ has
trace $0$ for \emph{every} irreducible unitary $\pi\not\cong\pi(\eps,\lambda)$, not only for the tempered ones; this is
where the strong gap condition is used. And all the facts used are proved in the development, (b) in the form (K10)(x)
below, which uses only the value of the Casimir operator and requires a gap condition. The vanishing of the trace
is obtained by an index argument for a Dirac-type operator on the tensor product of the smooth vectors with a fermionic
Fock space, in the spirit of \cite{AtiyahSchmid1977}.

\subsubsection*{The ingredients of (S1)}

For a finite-dimensional unitary representation $T$ of $K$ let $\chi_T(k)=\mathrm{tr}\,T(k)$; a \emph{$K$-type of weight
$\Lambda$} is an irreducible one with highest weight $\Lambda$. A function $\phi$ on $K$ is \emph{polynomial} if there is
$P\in\mathrm{Pol}_r$ with $\phi(k)=\mathrm{ev}_k(P)$ for all $k\in K$. Put
\begin{equation}\label{pf:multnc:eq-Deps}
D_\eps(k)=\det\bigl(\delta_{xy}-k_{aa'}\,\overline{k_{bb'}}\bigr)_{x=(a,b),\,y=(a',b')\in\mathrm{Mo}(\eps)}\qquad(k\in K),
\end{equation}
and $s(\eps)=\#\{(a,b)\in\mathrm{Mo}(\eps):b<a\}$. For $a\le r$ let $\rho_a=\frac{r-1}2-a$, let
$\rho_{c,a}=\frac12(\#\{b>a:\eps_b=\eps_a\}-\#\{b<a:\eps_b=\eps_a\})$, $\rho_{n,a}$ the same with $\eps_b\ne\eps_a$, and
$\rho_{h,a}=q/2$ if $\eps_a=1$, $\rho_{h,a}=-p/2$ if $\eps_a=-1$. For $\lambda\in\R^r$ put
$\Lambda_\lambda=\lambda-\rho_c+\rho_h$, let $\lambda+\rho_n-\rho_c$ be the weight of the lowest $K$-type, and
$c(\lambda)=|\lambda|^2-|\rho|^2$ the Casimir value. We use the following statements of the development.
\begin{enumerate}[label=(K\arabic*)]
\item (Gap conditions; \lref{GlobalTheory.PsCo.PcEF.gapConditions_of_largeGaps}) If $r$ is odd and $\lambda\in\Z^r$
satisfies $\multncst{G$_{100}$}$, then $\lambda$ is a Harish-Chandra parameter, it satisfies the \emph{standard gap
condition} ($\lambda$ is strictly decreasing and $|\lambda_a-\lambda_b|\ge h$ whenever $\eps_a\ne\eps_b$, for some
$h>rpq+1$) and a second gap condition used in (K10), and $|\lambda_a-\lambda_b|\ge h>\theta(\eps)$ whenever
$\eps_a\ne\eps_b$, for an explicit threshold $\theta(\eps)\ge rpq+1$ depending only on $\eps$.
\item (Compact case; \lref{GlobalTheory.PsCo.PcEF.compactDiscreteSeries_ownType}) If $p=0$ or $q=0$, then $D_\eps=1$,
$s(\eps)=0$, and every discrete series representation with parameter $\lambda$ is finite-dimensional and is a $K$-type of
weight $\Lambda_\lambda$.
\item (Schur orthogonality; \lref{GlobalTheory.PsCo.PcEF.norm_irreducibleCharacter_unitaryMaxCompact_eq_one})
$\int_K|\chi_T(k)|^2\,dk=1$ for every $K$-type $T$.
\item (Supertrace; \lref{ReciprocityTheory.ChartRestriction.DS.TRGap.supTrace}) If $p,q\ge1$, $\lambda$ is a
Harish-Chandra parameter with $|\lambda_a-\lambda_b|\ge h>\theta(\eps)$ whenever $\eps_a\ne\eps_b$, and $\pi$ is a
discrete series representation with parameter $\lambda$, then there is a $K$-type $T$ of weight $\Lambda_\lambda$ such
that for every $P\in\mathrm{Pol}_r$ the operator $\int_K\mathrm{ev}_k(P)D_\eps(k)\pi(k)\,dk$ has finite rank and trace
$(-1)^{s(\eps)}\int_K\mathrm{ev}_k(P)\chi_T(k)\,dk$.
\item (Polynomiality; \lref{GlobalTheory.PsCo.PcEF.coefPoly}, \lref{GlobalTheory.PsCo.PcEF.trPoly}) The matrix
coefficients of finite-dimensional unitary representations of $K$ are polynomial on $K$; hence so is
$k\mapsto\overline{\chi_T(k)}$.
\item (Block form; \lref{modeMatrixDet_eq_adjointDet},
\lref{GlobalTheory.PsCo.PcEF.maxCompactPolynomial_blockTimesConjDetPower}) There is $Q\in\mathrm{Pol}_r$ with
$\mathrm{ev}_k(Q)\,\mathrm{ev}_k(\det\Omega_P)=D_\eps(k)$ for $k\in K$; and for every $P\in\mathrm{Pol}_r$ there are
$M\in\N$ and $R\in\mathrm{Blk}(\eps)$ with $\mathrm{ev}_k(P)=\mathrm{ev}_k(\det(\bar x)^MR)$ for $k\in K$.
\item (Class functions; \lref{GlobalTheory.PsCo.PcEF.adjointDetTrace_classFunction}) $D_\eps$ and the characters
$\chi_T$ are class functions on $K$.
\item (The compact torus; \lref{GlobalTheory.PsCo.PcEF.harishChandraWeight_integral_dominant},
\lref{GlobalTheory.PsCo.PcEF.polyCoefIrrep_hasDerivedAction},
\lref{ReciprocityTheory.ChartRestriction.TR.Weyl.weyl}, \lref{GlobalTheory.PsCo.PcEF.torusId_full},
\lref{adjointDet_onCompactTorus}) For a Harish-Chandra parameter $\lambda$ the weight $\Lambda_\lambda$ is integral
and dominant for $K$; a $K$-type $T$ of weight $\Lambda_\lambda$ with polynomial coefficients is the irreducible $K$-module
of highest weight $\Lambda_\lambda$ of highest weight theory, and its character satisfies the Weyl character formula
\cite{Weyl1946}; an identity between the Weyl numerator and Harish-Chandra's numerator then gives
$(-1)^{s(\eps)}\chi_T(t(\vartheta))=\Theta_{\eps,\lambda}(\vartheta)\,D_\eps(t(\vartheta))$ whenever the $e^{i\vartheta_a}$
are distinct; and $|D_\eps(t(\vartheta))|^2=d(\vartheta)$.
\item (Matrix realisation; \lref{GlobalTheory.PsCo.PcEF.unitaryRep_upq_matrixFunction}) Every irreducible unitary
representation $\pi$ of $U(\eps)$ is the restriction of an irreducible unitary representation $\sigma$ on matrices, in
the sense used for discrete series models (\cref{pf:multnc:def-dsmodel}).
\item (Index and equivalence) The eleven statements used in \cref{pf:multnc:equiv} below.
\end{enumerate}

\begin{lemma}[The index of the discrete series]\label{pf:multnc:index}
Let $\lambda$ be a Harish-Chandra parameter with $|\lambda_a-\lambda_b|\ge h>\theta(\eps)$ whenever $\eps_a\ne\eps_b$, and
$\pi$ a discrete series representation of $U(\eps)$ with parameter $\lambda$. There is a $K$-type $T$ of weight
$\Lambda_\lambda$ such that, whenever $P\in\mathrm{Pol}_r$ satisfies $\mathrm{ev}_k(P)=\overline{\chi_T(k)}$ on $K$, every
operator $B=\int_K\mathrm{ev}_k(P)D_\eps(k)\pi(k)\,dk$ which has a trace has trace $(-1)^{s(\eps)}$.
\formallabel{GlobalTheory.PsCo.PcEF.exists_type}
\end{lemma}

\begin{proof}
If $p=0$ or $q=0$, then by (K2) $\pi$ itself is finite-dimensional and is a $K$-type of weight $\Lambda_\lambda$; take
$T=\pi$. Since $D_\eps=1$, $B=\int_K\overline{\chi_T(k)}\,T(k)\,dk$ and, the space being finite-dimensional, the trace of
the integral is the integral of the traces: $\mathrm{tr}\,B=\int_K|\chi_T|^2\,dk=1=(-1)^{s(\eps)}$ by (K3) and
$s(\eps)=0$. If $p,q\ge1$, take $T$ from (K4); then $B$ has finite rank with trace
$(-1)^{s(\eps)}\int_K\overline{\chi_T}\chi_T\,dk=(-1)^{s(\eps)}$ by (K3). In both cases a trace of $B$ is unique, which
gives the claim.
\end{proof}

The second lemma is the converse direction: a non-zero index forces the representation to be the discrete series. Let
$\sigma$ be an irreducible unitary representation on a Hilbert space $H$, $H^\infty$ the space of smooth vectors with the
derived action $E_{ab}$ of the complexified Lie algebra, $\Phi$ the fermionic Fock space on generators indexed by pairs
$(i,a)$, with creation and annihilation operators $c^\pm_{ia}$, and $\mathcal P\subset H^\infty\otimes\Phi$ the
\emph{physical subspace}, the common kernel of the $c^-_{ia}$ with $(i,a)\notin\mathrm{Mo}(\eps)$. The development defines
on $H^\infty\otimes\Phi$ a Dirac-type operator $\mathsf d$, built from the $E_{ia}\otimes c^\pm_{ia}$ with
$(i,a)\in\mathrm{Mo}(\eps)$, a grading operator $\mathsf g$, and an action of $K$; and the supertrace
$\mathrm{str}_Y(k)=\mathrm{tr}(\mathsf g\circ k\mid Y)$ on finite-dimensional stable subspaces $Y$.

\begin{lemma}[Equivalence from a non-zero index]\label{pf:multnc:equiv}
Let $\lambda$ satisfy the two gap conditions of (K1), $T$ a $K$-type of weight $\Lambda_\lambda$, $P\in\mathrm{Pol}_r$ with
$\mathrm{ev}_k(P)=\overline{\chi_T(k)}$ on $K$, $\sigma$ an irreducible unitary representation of $U(\eps)$ and $\pi$ a
discrete series representation with parameter $\lambda$. If $B=\int_K\mathrm{ev}_k(P)D_\eps(k)\sigma(k)\,dk$ has trace
$c\ne0$, then $\sigma$ is unitarily equivalent to $\pi$.
\formallabel{GlobalTheory.PsCo.PcEF.equiv_of_index}
\end{lemma}

\begin{proof}
We use eleven statements of the development, (K10)(i)--(x) (item (viii) consists of two of them).
\begin{enumerate}[label=(\roman*)]
\item The Casimir operator acts on $H^\infty$ by a scalar
(\lref{GlobalTheory.HcPort.casimirActsByScalar_indefUnitary}).
\item For every $d$ there is a finite-dimensional $K$-stable subspace $L\subset H^\infty$ containing
$\int_K\mathrm{ev}_k(P')\sigma(k)u\,dk$ for all $u\in H$ and $\deg P'\le d$
(\lref{GlobalTheory.PsCo.PcEF.rangeFD_closed}).
\item Given (i) and such an $L$, there is a finite-dimensional subspace $Y$ with $(L\otimes\Phi)\cap\mathcal P\subset
Y\subset\mathcal P$ which is stable under $\mathsf d$, $\mathsf g$ and $K$
(\lref{GlobalTheory.PsCo.PcEF.casimirScalar_stableSumSpace}).
\item For every $P'$ there is $d$ such that, for $L$ as in (ii) and $Y$ as in (iii), the operator
$\int_K\mathrm{ev}_k(P')D_\eps(k)\sigma(k)\,dk$ maps $H$ into $L$ and its trace on $L$ is
$\int_K\mathrm{ev}_k(P')\,\mathrm{str}_Y(k)\,dk$ (\lref{GlobalTheory.PsCo.PcEF.StrIdx_proof2}).
\item For such $Y$ there are a finite-dimensional $K$-stable subspace $Z\subset Y\cap\ker\mathsf d$ and an operator
$\Gamma$ on $Z$ with $\mathrm{str}_Y(k)=\mathrm{tr}(\Gamma\circ k\mid Z)$ for $k\in K$: the supertrace is carried by the
kernel of $\mathsf d$ (\lref{GlobalTheory.PsCo.PcEF.supertrace_carriedBy_diracKernel}).
\item If $T,S$ are finite-dimensional unitary representations of $K$, $\Gamma$ an operator on the space of $S$, and
$\int_K\overline{\chi_T(k)}\,\mathrm{tr}(\Gamma\circ S(k))\,dk\ne0$, there is a non-zero $K$-intertwiner from $T$ to $S$
(\lref{GlobalTheory.PsCo.PcEF.tracePairing_nonzero_intertwiner}).
\item A non-zero $K$-intertwiner from a $K$-type of weight $\Lambda$ into $Z$ produces a non-zero $K$-highest weight vector
$z\in Z$ of weight $\Lambda$ (\lref{GlobalTheory.PsCo.PcEF.intertwiner_preserves_highestWeight}).
\item The coordinates of the elements of a finite-dimensional $K$-stable subspace of $H^\infty\otimes\Phi$ span a
finite-dimensional $K$-stable subspace of $H^\infty$ (\lref{GlobalTheory.PsCo.PcEF.fockCoordinates_span_finiteDimensional});
and if $z$ is a $K$-highest weight vector of weight $\Lambda_\lambda$ with $\mathsf dz=0$, the Casimir operator acts by
$c(\lambda)$ on all coordinates of $z$ (\lref{GlobalTheory.PsCo.PcEF.casimirValue_highestWeightCoordinates}).
\item Under the second gap condition, and using that the derived action is infinitesimally unitary (as $\sigma$ is
unitary): if $z\ne0$ is physical, a $K$-highest weight vector of weight $\Lambda_\lambda$ whose
coordinates lie in a finite-dimensional $K$-stable subspace on which the Casimir operator acts by $c(\lambda)$, then
$H^\infty$ contains a $K$-irreducible highest weight submodule $V\ne0$ of weight $\lambda+\rho_n-\rho_c$ on which the
Casimir operator acts by $c(\lambda)$ (\lref{GlobalTheory.PsCo.PcEF.highestWeightKernel_lowestType}).
\item Under the standard gap condition, an irreducible unitary representation whose smooth vectors contain such a $V$ is
unitarily equivalent to every discrete series representation with parameter $\lambda$
(\lref{GlobalTheory.PsCo.PcEF.eq_discreteSeries_of_lowestKType_casimir}); compare \cref{pf:multnc:lit-ds}(b).
\end{enumerate}
Take $d$ from (iv) for $P$, then $L$ from (ii) and $Y$ from (iii). By (iv), $c=\mathrm{tr}(B\mid L)=\int_K\mathrm{ev}_k(P)
\,\mathrm{str}_Y(k)\,dk$, and by (v) this is $\int_K\overline{\chi_T(k)}\,\mathrm{tr}(\Gamma\circ k\mid Z)\,dk$. Since
$c\ne0$, (vi) gives a non-zero $K$-intertwiner $T\to Z$, and (vii) a non-zero highest weight vector $z\in Z$ of weight
$\Lambda_\lambda$; it is physical ($Z\subset Y\subset\mathcal P$) and $\mathsf dz=0$. By (viii) the coordinates of the
elements of $Z$ span a finite-dimensional $K$-stable $S\subset H^\infty$ and the Casimir operator acts by $c(\lambda)$ on
the coordinates of $z$; by (ix) we obtain $V$; and by (x), $\sigma\cong\pi$.
\end{proof}

\begin{theorem}[The $K$-polynomial]\label{pf:multnc:kdata}
Statements (K1)--(K10) imply (S1). Consequently (S1) holds.
\formallabel{GlobalTheory.PsCo.PcEF.stKData_of}
\formallabel{GlobalTheory.PsCo.Pc1.maximalCompact_pseudoCoefficientData'}
\formallabel{GlobalTheory.T.maximalCompact_pseudoCoefficientData}
\end{theorem}

\begin{proof}
Let $r$ be odd, $\lambda\in\Z^r$ with $\multncst{G$_{100}$}$ and $\pi_0=\pi(\eps,\lambda)$. By (K1), $\lambda$ satisfies the
hypotheses of \cref{pf:multnc:index,pf:multnc:equiv}. Let $T$ be the $K$-type of \cref{pf:multnc:index}; by (K5) there is
$P\in\mathrm{Pol}_r$ with $\mathrm{ev}_k(P)=\overline{\chi_T(k)}$ on $K$. Write $s=s(\eps)$. By (K6) choose $Q$ with
$\mathrm{ev}_k(Q)\mathrm{ev}_k(\det\Omega_P)=D_\eps(k)$ on $K$, and then $M$ and $R\in\mathrm{Blk}(\eps)$ with
$\mathrm{ev}_k(\det(\bar x)^MR)=\mathrm{ev}_k((-1)^sPQ)$ on $K$. Since $\mathrm{top}(M,R)=\det(\bar x)^MR\det\Omega_P$, for
$k\in K$
\begin{equation}\label{pf:multnc:eq-topK}
\mathrm{ev}_k(\mathrm{top}(M,R))=(-1)^s\,\mathrm{ev}_k(P)\,\mathrm{ev}_k(Q)\,\mathrm{ev}_k(\det\Omega_P)=
(-1)^s\,\overline{\chi_T(k)}\,D_\eps(k).
\end{equation}
We verify the four properties of (S1) for $M$ and $R$.

\emph{Class function.} By \eqref{pf:multnc:eq-topK} and (K7).

\emph{Values on the torus.} Let the $e^{i\vartheta_a}$ be distinct. Since $(-1)^s$ is real, \eqref{pf:multnc:eq-topK} and
(K8) give $\mathrm{ev}_{t(\vartheta)}(\mathrm{top})=\overline{(-1)^s\chi_T(t(\vartheta))}\,D_\eps(t(\vartheta))=
\overline{\Theta_{\eps,\lambda}(\vartheta)}\,|D_\eps(t(\vartheta))|^2=\overline{\Theta_{\eps,\lambda}(\vartheta)}\,
d(\vartheta)$.

\emph{Trace on $\pi_0$.} Let $B=\int_K\mathrm{ev}_k(\mathrm{top})\pi_0(k)\,dk$ have trace $c$. By
\eqref{pf:multnc:eq-topK}, $(-1)^sB=\int_K\mathrm{ev}_k(P)D_\eps(k)\pi_0(k)\,dk$, which has trace $(-1)^sc$; by
\cref{pf:multnc:index}, $(-1)^sc=(-1)^s$, so $c=1$.

\emph{Trace on other representations.} Let $\pi\not\cong\pi_0$ be irreducible unitary and suppose that
$B_\pi=\int_K\mathrm{ev}_k(\mathrm{top})\pi(k)\,dk$ has trace $c\ne0$. Realise $\pi$ as $\sigma$ by (K9). Then
$\int_K\mathrm{ev}_k(P)D_\eps(k)\sigma(k)\,dk=(-1)^sB_\pi$ has trace $(-1)^sc\ne0$, and \cref{pf:multnc:equiv} gives
$\sigma\cong\pi_0$, a contradiction. Hence $c=0$.
\end{proof}

\subsubsection*{The limit formula (S6)}

The limit (S6) is a Laplace-type computation in polar coordinates $U(\eps)=K\exp(\mathfrak p)$. For
$B\in\C^{\mathrm{Mo}(\eps)}$ let $p(B)\in M_r(\C)$ be the hermitian matrix with $p(B)_{ab}=B_{(a,b)}$ and
$p(B)_{ba}=\overline{B_{(a,b)}}$ for $(a,b)\in\mathrm{Mo}(\eps)$ and all other entries $0$, and let $\mathrm{blk}(g)=
(g_{ab})_{(a,b)\in\mathrm{Mo}(\eps)}$, so that $r_M(g)=|\mathrm{blk}(g)|^2$. We use the following statements of the
development.
\begin{enumerate}[label=(B\arabic*)]
\item (\lref{GlobalTheory.PsCo.Pc3.stEB}) There is a continuous map $B\mapsto e(B)$ from $\C^{\mathrm{Mo}(\eps)}$ to
$U(\eps)$ with $e(0)=1$, $e(B)=\exp p(B)$ and $e(B)^{-1}=\exp(-p(B))$ as matrices.
\item (\lref{GlobalTheory.PsCo.Pc3.stR}) $g\mapsto r_M(g)$ is continuous and bi-$K$-invariant on $U(\eps)$, vanishes
exactly on $K$, and its sublevel sets are compact.
\item (\lref{GlobalTheory.PsCo.Pc3.stPolar}) For every Haar measure $\mu$ there is a continuous $J$ on
$\C^{\mathrm{Mo}(\eps)}$ with $J(0)>0$ such that for continuous compactly supported $f$, $\int f\,d\mu=\int J(B)\int_K
f(k\,e(B))\,dk\,dB=\int J(B)\int_Kf(e(B)\,k)\,dk\,dB$, with $dB$ the Lebesgue measure.
\item (\lref{GlobalTheory.PsCo.Pc3.stExpLow}) $|B|^2\le r_M(\exp p(B))$.
\item (\lref{GlobalTheory.PsCo.Pc3.stExpDeriv}) At $B=0$ the map $B\mapsto\mathrm{blk}(\exp p(B))$ has derivative the
identity; and for $s\in\C^r$ the map $B\mapsto\mathrm{blk}(\exp(-p(B))\,\mathrm{diag}(s)\exp p(B))$ has derivative
$w\mapsto((s_a-s_b)w_{(a,b)})_{(a,b)}$.
\item (\lref{GlobalTheory.PsCo.Pcd.stConjLow}) If $s\in U(\eps)$ is diagonal with distinct entries, then for every
$\delta>0$ there is $\eta>0$ such that $r_M(x^{-1}sx)\le\eta$ implies $r_M(x)\le\delta$.
\item (\lref{GlobalTheory.PsCo.Pcd.stFamBlow}) Let $\iota$ be a finite set with $\#\iota=N$,
$\beta\colon\C^\iota\to\C^\iota$ continuous with $\beta(0)=0$, differentiable at $0$ with derivative $w\mapsto(d_xw_x)$ for
some $d\in(\C^\times)^\iota$, and such that for every $\delta>0$ there is $\eta>0$ with $|\beta(B)|^2\le\eta\Rightarrow|B|\le
\delta$. Let $G$ be a profile of order $N$ and $\psi_0,\dots,\psi_N$ continuous on $\C^\iota$. Then, as $t\to\infty$,
\[
\int\sum_{k=0}^N\Bigl(\frac{d}{dy}\Bigr)^k\bigl[G(ty)\bigr]_{y=|\beta(B)|^2}\psi_k(B)\,dB\longrightarrow
\frac{\psi_N(0)\,I_G}{\prod_x|d_x|^2},\qquad I_G=\int_{\C^\iota}G^{(N)}(|w|^2)\,dw.
\]
\item (\lref{GlobalTheory.PsCo.Pc3.stConst_holds}) $I_G\ne0$ for every profile $G$ of order $N=\#\iota$.
\end{enumerate}

\begin{proposition}[The limit formula]\label{pf:multnc:stlim}
Statements (B1)--(B8) imply (S6), with $c=J(0)\,I_G$. Consequently (S6) holds.
\formallabel{GlobalTheory.PsCo.Pcd.lim1}
\formallabel{GlobalTheory.PsCo.Pcd.lim2}
\formallabel{GlobalTheory.PsCo.Pcd.stLim_of}
\formallabel{GlobalTheory.PsCo.Pcd.stLim}
\end{proposition}

\begin{proof}
Fix $\mu$, a profile $G$ of order $N$ and $J$ as in (B3); put $c=J(0)I_G$, which is non-zero by (B3) and (B8). Let $p$ be
a sequence in $\mathrm{Pol}_r$ and $\phi_t=\phi_{G,p,t}$; for $t>0$, $\phi_t$ is continuous and vanishes where
$t\,r_M>1$ (S4), so $\phi_t|_{U(\eps)}$ has compact support by (B2).

(i) Let $w$ be continuous and $\beta(B)=\mathrm{blk}(e(B))$. By (B2), $r_M(k\,e(B))=r_M(e(B))=|\beta(B)|^2$ for $k\in K$,
so $\phi_t(k\,e(B))=\sum_j(\frac d{dy})^j[G(ty)]_{y=|\beta(B)|^2}\mathrm{ev}_{k e(B)}(p_j)$ and, by (B3),
\begin{gather*}
\int\phi_tw\,d\mu=\int\sum_{j=0}^N\Bigl(\frac d{dy}\Bigr)^j\bigl[G(ty)\bigr]_{y=|\beta(B)|^2}\psi_j(B)\,dB,\\
\psi_j(B)=J(B)\int_K\mathrm{ev}_{k e(B)}(p_j)\,w(k\,e(B))\,dk .
\end{gather*}
The $\psi_j$ are continuous (the integrand is jointly continuous and $K$ is compact). The map $\beta$ is continuous with
$\beta(0)=\mathrm{blk}(1)=0$ and derivative the identity at $0$ (B1, B5), and $|B|^2\le|\beta(B)|^2$ by (B4), so the
condition of (B7) holds with $\eta=\delta^2$. By (B7) with $d=1$ the integral tends to $\psi_N(0)I_G=c\int_K
\mathrm{ev}_k(p_N)w(k)\,dk$.

(ii) Let $s=t(\vartheta)$ with distinct $d_a=e^{i\vartheta_a}$, and assume $\mathrm{ev}_{k^{-1}sk}(p_N)=\mathrm{ev}_s(p_N)$
for $k\in K$. Take $\eta_1$ from (B6) for $\delta=1$. For $t>1/\eta_1$ the function $x\mapsto\phi_t(x^{-1}sx)$ is continuous
with compact support: where it is non-zero, $r_M(x^{-1}sx)\le1/t<\eta_1$, hence $r_M(x)\le1$, a compact set by (B2). Put
$\beta(B)=\mathrm{blk}(e(B)^{-1}s\,e(B))$. By the second form of (B3) and the invariance $r_M(k^{-1}yk)=r_M(y)$,
\begin{gather*}
\int\phi_t(x^{-1}sx)\,d\mu(x)=\int\sum_{j=0}^N\Bigl(\frac d{dy}\Bigr)^j\bigl[G(ty)\bigr]_{y=|\beta(B)|^2}\psi_j(B)\,dB,\\
\psi_j(B)=J(B)\int_K\mathrm{ev}_{k^{-1}e(B)^{-1}se(B)k}(p_j)\,dk .
\end{gather*}
Here $\beta$ is continuous, $\beta(0)=\mathrm{blk}(s)=0$ as $s$ is diagonal, and its derivative at $0$ is
$w\mapsto((d_a-d_b)w_{(a,b)})$ by (B5), with $d_a\ne d_b$ for $(a,b)\in\mathrm{Mo}(\eps)$ because $a\ne b$. Given
$\delta>0$, (B6) gives $\eta$ with $|\beta(B)|^2=r_M(e(B)^{-1}se(B))\le\eta\Rightarrow r_M(e(B))\le\delta^2$, and then
$|B|\le\delta$ by (B4). By (B7) the integral tends to $\psi_N(0)I_G/\prod_{(a,b)}|d_a-d_b|^2$. Finally
$\psi_N(0)=J(0)\int_K\mathrm{ev}_{k^{-1}sk}(p_N)\,dk=J(0)\,\mathrm{ev}_s(p_N)$ and $\prod_{(a,b)\in\mathrm{Mo}(\eps)}
|d_a-d_b|^2=d(\vartheta)$, so the limit is $c\,\mathrm{ev}_s(p_N)/d(\vartheta)$.
\end{proof}

% ===========================================================================
\subsection{Class integrals and the elliptic expansion}\label{pf:multnc:sec-classint}

We prove the inputs (d)--(f) of the proof of \cref{pf:multnc:elltrace} and the elliptic expansion used in
\cref{pf:multnc:ellpart}. Let $E/F$ be a CM extension, $J\in M_n(E)$ hermitian and non-degenerate and $G=\Uni(J)$; for a
hermitian space $V$ of the datum, $J=J_V$ and $G=G_V$. Let $\mathcal C(G)$ be the span of the products of a smooth compactly
supported function on $G(F\otimes\R)$ and a locally constant compactly supported function on $G(\multncAf)$; its elements
are continuous with compact support. For $f\in\mathcal C(G)$ and
$\gamma\in G(F)$ the \emph{class kernel} is $k^f_\gamma(x)=\sum_{\delta\in[\gamma]}f(x^{-1}\delta x)$, where $[\gamma]$ is the
$G(F)$-conjugacy class of $\gamma$; it is left $G(F)$-invariant, since $\delta\mapsto g^{-1}\delta g$ permutes $[\gamma]$ for
$g\in G(F)$, and it only depends on the class $[\gamma]$. A \emph{cusp function} is an $f\in\mathcal C(G)$ whose unipotent integrals along all proper
isotropic flags vanish. We say that the \emph{class integral of $f$ at $\gamma$ is $c$} if, for every fundamental domain
$\mathcal F$ of $G(F)$ in $G(\A)$, the series defining $k^f_\gamma$ converges absolutely at every point, $k^f_\gamma$ is
integrable on $\mathcal F$ and $\int_{\mathcal F}k^f_\gamma=c$.

\subsubsection*{Reduction theory}

A \emph{Witt frame} consists of $r\ge0$ and $X,Y\in M_{n\times r}(E)$ with $X^*JX=Y^*JY=0$, $X^*JY=1$ and anisotropic
orthogonal complement of the columns; Witt frames exist (\lref{GlobalTheory.SiegelSet.exists_isWittFrame}, Witt's
decomposition). Let $P_{\min}$ be the stabiliser of the flag spanned by the columns of $X$ and $P^1_{\min}(\A)$ the
elements of $P_{\min}(\A)$ whose diagonal entries on that flag have idelic norm $1$. For $t\in\R^r_{>0}$ let $a_t$ denote
the archimedean elements which act on the $i$th column of $X$ by $t_i$, on that of $Y$ by $t_i^{-1}$ and trivially on the
complement. For $c>0$ let $\mathrm{cone}(c)=\{t:c\,t_{i+1}\le t_i\ (i<r),\ c\le t_r\}$, and for compact $\Omega\subset
P^1_{\min}(\A)$ and a compact subgroup $K\subset G(\A)$ the \emph{Siegel set} $\mathfrak S=\{p\,a_tk:p\in\Omega,\ t\in
\mathrm{cone}(c),\ k\in K\}$. Let $\delta(t)=\prod_{i<r}t_i^{2[F:\Q](n-2i-1)}$.

\begin{definition}[The reduction-theory statement]\label{pf:multnc:def-siegel}
The statement $\multncst{Red}$ is: for every CM extension $E/F$, hermitian non-degenerate $J$ and Witt frame $(X,Y)$ there
are a compact subgroup $K$, a compact $\Omega\subset P^1_{\min}(\A)$ and $c>0$ such that $G(\A)=P_{\min}(\A)K$; the
elements of $P_{\min}(\A)\cap K$ are in good position (they factor through the Levi and unipotent parts inside $K$);
$G(\A)=G(F)\,\mathfrak S$; the set $\{a_t^{-1}pa_t:p\in\Omega,t\in\mathrm{cone}(c)\}$ is relatively compact; the heights of
$g$ and $\gamma g$ are polynomially comparable for $g\in\mathfrak S$, $\gamma\in G(F)$; and smooth cuspidal functions of
moderate growth decay rapidly on $\mathfrak S$.
\formal{PositivityTheory.ThKR.SiegelCore}
\end{definition}

This is reduction theory for unitary groups \cite{Borel1969} together with the rapid decay of cusp forms
\cite{HarishChandra1968}. It is proved in the development (\lref{PositivityTheory.Prw.siegelCore_holds}), and the
published form \lref{PositivityTheory.ThKR.siegel_of_core} derives from it, among other things, the closedness of
$\mathfrak S$ and the \emph{Siegel integration bound}: for every Haar measure $\mu$ on $G(\A)$ there is $C$ such that
$\int_{\mathfrak S}\phi\,d\mu\le C\int_{\mathrm{cone}(c)}h(t)\,(\delta(t)\prod_it_i)^{-1}\,dt$ whenever $\phi\ge0$,
$h$ is measurable and $\phi(p\,a_tk)\le h(t)$ on $\mathfrak S$.

\begin{lemma}[Finiteness of regular classes]\label{pf:multnc:regfin}
Assume $\multncst{Red}$. Let $C\subset G(\A)$ be compact. Then only finitely many $G(F)$-conjugacy classes of regular
semisimple elements $\gamma\in G(F)$ contain an element $\delta$ with $x^{-1}\delta x\in C$ for some $x\in G(\A)$. In
particular, for $f\in\mathcal C(G)$ only finitely many regular semisimple classes meet the support of $x\mapsto
f(x^{-1}\cdot x)$.
\formallabel{GlobalTheory.CompanionTheory.CFin.classes_finite}
\formallabel{GlobalTheory.CompanionTheory.CFin.regClassFin}
\formallabel{GlobalTheory.T.siegel_implies_regular_class_finite}
\end{lemma}

\begin{proof}
Fix a Witt frame $(X,Y)$ and complete it (\lref{GlobalTheory.T.exists_witt_complement_frame}) by a basis $Z$ of the
orthogonal complement of its columns, so that $Q=(X\,|\,Z\,|\,Y)$ is invertible, $X^*JZ=Y^*JZ=0$; index the columns of $Q$
by $\mathrm{Idx}=[r]\sqcup[m]\sqcup[r]$. For $\delta\in G(F)$ put $\tilde\delta=Q^{-1}\delta Q$; $\delta\mapsto\tilde\delta$
is injective. Take $K,\Omega,c$ from $\multncst{Red}$ and put $C_1=\{kyk^{-1}:k\in K,y\in C\}$ and
$\mathcal C=\mathrm{pr}_f(\{pyp^{-1}:p\in\Omega,y\in C_1\})$, both compact. For a finite set $S\subset[r]$ let
$\mathrm{lev}_S\colon\mathrm{Idx}\to\Z$ be the block index attached to $S$ ($\#\{s\in S:s<i\}-\#S$ on the $X$-part, $0$ on
the $Z$-part and $\#S-\#\{s\in S:s<i\}$ on the $Y$-part), and $\mathrm{bd}_S(N)$ the matrix which agrees with $N$ on the
entries $(u,v)$ with $\mathrm{lev}_S(u)=\mathrm{lev}_S(v)$ and is $0$ elsewhere. Let $\mathrm{wt}(t)$ be $t_i$ on the $X$-part,
$1$ on the $Z$-part and $t_i^{-1}$ on the $Y$-part. We use:
\begin{enumerate}[label=(\alph*)]
\item (\lref{GlobalTheory.T.siegel_compact_translate_bound}) there is $t_1$ such that if $p\in\Omega$,
$t\in\mathrm{cone}(c)$, $k\in K$ and $(pa_tk)^{-1}\gamma(pa_tk)\in C$, then $\gamma$ maps $X_l$ into the span of
$X_1,\dots,X_j$ for every $l\le j$, for every \emph{deep} index $j$, i.e.\ with $t_j>t_1$ and $t_1t_{j+1}<t_j$;
\item (\lref{GlobalTheory.T.deep_cone_weight_ratio_bound}) for $t'_1\ge1$ there is $R$ such that for $t\in\mathrm{cone}
(c)$ and the set $S$ of indices that are deep for $t'_1$, $\mathrm{wt}(t)_u\le R\,\mathrm{wt}(t)_v$ whenever
$\mathrm{lev}_S(u)=\mathrm{lev}_S(v)$;
\item (\lref{GlobalTheory.CompanionTheory.CFin.arch_bound}) for every $R$ there is $B$ such that, if $p\in\Omega$,
$t\in\R^r_{>0}$, $y\in C_1$, $S$ satisfies the weight condition of (b), $\delta\in G(F)$ equals $pa_tya_t^{-1}p^{-1}$ and
stabilises the flags of (a) for all $j\in S$, then all archimedean absolute values of the entries of
$\mathrm{bd}_S(\tilde\delta)$ are at most $B$;
\item (\lref{GlobalTheory.T.compact_finite_rational_denominator}) there is $d\ge1$ such that $d\tilde\delta$ has
integral entries whenever the finite component of $\delta\in G(F)$ lies in $\mathcal C$;
\item (\lref{GlobalTheory.T.bounded_integral_matrices_finite}) the set of matrices $N$ over $E$ with $dN$ integral and
all archimedean absolute values of the entries at most $B$ is finite;
\item (\lref{GlobalTheory.CompanionTheory.CFin.exists_conj_levi}) if $\delta$ is regular semisimple and stabilises the
flags of (a) for all $j\in S$, then there is $\delta'\in G(F)$ conjugate to $\delta$ with $\tilde\delta'=
\mathrm{bd}_S(\tilde\delta)$.
\end{enumerate}
Take $t_1$ from (a), $R$ from (b) for $t'_1=\max(t_1,1)$, $B$ from (c), $d$ from (d) and the finite set $\mathcal N$ of (e).
Let $[\gamma]$ be a regular semisimple class and $\delta\in[\gamma]$, $x\in G(\A)$ with $x^{-1}\delta x\in C$. By
$\multncst{Red}$ write $x=\eta\,pa_tk$ with $\eta\in G(F)$ and $pa_tk\in\mathfrak S$, and put $\delta_1=\eta^{-1}\delta\eta
\in[\gamma]$; then $(pa_tk)^{-1}\delta_1(pa_tk)\in C$, and $\delta_1$ is regular semisimple. Let $S$ be the set of indices
that are deep for $t'_1$; by (a), $\delta_1$ stabilises the corresponding flags, and (f) gives $\delta'\in[\gamma]$ with
$\tilde\delta'=\mathrm{bd}_S(\tilde\delta_1)$. With $y=k\,(pa_tk)^{-1}\delta_1(pa_tk)\,k^{-1}\in C_1$ we have
$\delta_1=pa_tya_t^{-1}p^{-1}$; since $a_t$ has trivial finite component, the finite component of $\delta_1$ is that of
$pyp^{-1}$, which lies in $\mathcal C$. By (d), $d\tilde\delta_1$ and hence $d\,\mathrm{bd}_S(\tilde\delta_1)$ are integral,
and by (b) and (c) the entries of $\mathrm{bd}_S(\tilde\delta_1)$ are bounded by $B$ at all archimedean places. So
$\tilde\delta'\in\mathcal N$ and $[\gamma]=[\delta']$ lies in the finite set $\{[\delta']:\tilde\delta'\in\mathcal N\}$. The
last assertion follows with $C=\supp f$.
\end{proof}

\begin{lemma}[Integrability of class kernels]\label{pf:multnc:regint}
Let $V$ be a hermitian space of the datum, $\mu$ a Haar measure on $G(\A)$, $f\in\mathcal C(G)$ a cusp function and
$\gamma\in G(F)$ regular semisimple. Then $k^f_\gamma$ descends to an integrable function on $G(F)\backslash G(\A)$.
\formallabel{MainStatements.WeightKernel.RegInt.regClassInt}
\formallabel{MainStatements.WeightKernel.regClassInt_holds}
\end{lemma}

\begin{proof}
By \lref{PositivityTheory.Prw.siegelCore_holds}, $\multncst{Red}$ holds; fix a Witt frame and $K,\Omega,c$ as there,
with Siegel set $\mathfrak S$. The function $k=k^f_\gamma$ is continuous
(\lref{MainStatements.WeightKernel.RegInt.continuous_classKernel}: near each point only finitely many terms are
non-zero, since $f$ has compact support and $G(F)$ is discrete) and left $G(F)$-invariant. Since $\mathfrak S$ is
measurable and $G(\A)=G(F)\mathfrak S$, a left $G(F)$-invariant measurable function which is integrable on $\mathfrak S$
descends to an integrable function on the quotient
(\lref{GlobalTheory.CompanionTheory.AutCover.integrable_descend_of_cover}). Choose $m_0$ with
$\int_{\mathrm{cone}(c)}\prod_it_i^{-m_0}\,(\delta(t)\prod_it_i)^{-1}dt<\infty$
(\lref{GlobalTheory.T.cone_integral_finite_large_exponent}). The \emph{class bound}
(\lref{MainStatements.WeightKernel.RegInt.class_bound}) gives $B\ge0$ with
\[
|k(p\,a_tk')|\le B\prod_it_i^{-m_0}\qquad(p\in\Omega,\ t\in\mathrm{cone}(c),\ k'\in K);
\]
its proof combines a polynomial bound in $\prod_i\max(t_i,1)$ for the number of $\delta\in[\gamma]$ with
$a_t^{-1}\delta a_t$ in a fixed compact set with the decay of the unipotent sums of the cusp function $f$ along the deep
coordinates (\lref{MainStatements.WeightKernel.Dec.decay_holds}). By the Siegel integration bound with
$h(t)=B\prod_it_i^{-m_0}$, $\int_{\mathfrak S}|k|\,d\mu\le CB\int_{\mathrm{cone}(c)}\prod_it_i^{-m_0}(\delta(t)\prod_i
t_i)^{-1}dt<\infty$.
\end{proof}

\subsubsection*{Classes with reducible characteristic polynomial}

Let $\gamma\in G(F)$ be regular semisimple, $E[\gamma]\subset M_n(E)$ the $E$-algebra generated by $\gamma$, and
$E[\gamma]^+$ the $F$-algebra of its elements which are self-adjoint for $J$. We say that $\gamma$ is of \emph{non-field
type} if $E[\gamma]^+\cong\prod_{i<k}K_i$ with fields $K_i$, $k\ge2$ and every $K_i\otimes_FE$ a field. For a typed place $v$
and $0<a<n$ let $U_a(k_v)\subset\GL_n(k_v)$ be the unipotent radical of the standard parabolic subgroup with blocks
$(a,n-a)$. The type datum $\multncTy=(S,(\tau_v))$ satisfies the \emph{radical hypothesis} $\multncst{Rad}$ if for every
$0<a<n$ there is $v\in S$ such that $\tau_v$ has no non-zero vector fixed by $U_a(k_v)$.

\begin{lemma}[Non-field type classes]\label{pf:multnc:nonfield}
Let $V$ have odd rank, $M$ a discrete series model, $\multncTy$ a type datum satisfying $\multncst{Rad}$, $U$ an auxiliary
place, $\mu$ the restricted product of Haar measures $\mu_w,\mu_v$, $(f_w)$ an archimedean test family for $M$, $(f_v)$ a
finite test family for $(\multncTy,U)$, $f=f_\infty\otimes f_f$, and $\gamma\in G(F)$ of non-field type. Then:
\begin{enumerate}[label=(\alph*)]
\item (local clause) for every isomorphism $e\colon E[\gamma]^+\cong\prod_{i<k}K_i$ and every $i$, the idempotent
$\epsilon_i=e^{-1}(1_{K_i})$ satisfies $\epsilon_i^2=\epsilon_i$, is self-adjoint, commutes with $E[\gamma]$, has rank
$\dim_E(E\otimes_FK_i)$ with $0<\mathrm{rank}\,\epsilon_i<n$; and $O_\gamma(f_v)=0$, for every Haar measure on the
centraliser, at every $v\in S$ at which $\tau_v$ has no non-zero vector fixed by $U_{\mathrm{rank}\,\epsilon_i}(k_v)$;
\item (product clause) for every restricted product $dt$ of Haar measures on the centralisers of $\gamma$ and every
fundamental domain $\mathcal T$ of $G_\gamma(F)$ in $G_\gamma(\A)$: $\mathrm{vol}(\mathcal T)<\infty$; the local orbital
integrals $o_w=O_\gamma(f_w)$, $o_v=O_\gamma(f_v)$ exist, $o_v=1$ for almost all $v$, $O_\gamma(f)=\prod_wo_w\prod_vo_v$,
and the class integral of $f$ at $\gamma$ is $\mathrm{vol}(\mathcal T)\prod_wo_w\prod_vo_v$;
\item (zero clause) the class integral of $f$ at $\gamma$ is $0$.
\end{enumerate}
\formallabel{GlobalTheory.NonFieldType.localClause}
\formallabel{GlobalTheory.NonFieldType.productClauseX}
\formallabel{GlobalTheory.NonFieldType.zeroClause}
\formallabel{GlobalTheory.NonFieldType.nonFieldTypeVanishingX}
\formallabel{GlobalTheory.PsCo.Plug.nonFieldT}
\end{lemma}

\begin{proof}
(a) The properties of $\epsilon_i$ are \lref{GlobalTheory.T.selfAdjointFactorIdempotent_orthogonal_splitting_rank}
(applied with the structure statement \lref{GlobalTheory.T.regularSemisimple_unitaryGroup_concrete} on regular
semisimple elements of unitary groups). The vanishing is the type-test vanishing of weighted orbital integrals
(\lref{GlobalTheory.T.weightedOrbitalIntegral_typeTest_vanish}, with weight $1$): if $\gamma$ commutes with a
self-adjoint idempotent $\epsilon$ with $(1-\epsilon)^*J\epsilon=0$ and $0<\mathrm{rank}\,\epsilon<n$, and $\tau_v$ has no
non-zero vector fixed by $U_{\mathrm{rank}\,\epsilon}(k_v)$, then the orbital integral at $\gamma$ of the type test function
$f_v$ vanishes. It uses that the constant terms of the type idempotent along $U_a$ vanish
(\lref{GlobalTheory.T.typeIdempotent_constantTerm_vanish}).

(b) Since $\gamma$ is regular semisimple, its centralisers at all places are commutative, so their Haar measures are also
right invariant, and the orbits of $\gamma$ are proper (\lref{GlobalTheory.Orb.isProperOrbit_GArch} and its finite
analogue). Since $\gamma$ is of non-field type its centraliser is anisotropic and $\mathrm{vol}(\mathcal T)<\infty$
(\lref{GlobalTheory.NonFieldType.covol_lt_top}). The factorisation of the adelic orbital integral, with $o_v=1$ for
almost all $v$, is \lref{GlobalTheory.Orb.hasOrbInt_tensor} (the $f_w$ are smooth with compact support and $(f_v)$ is a
finite family). Finally the class integral equals $\mathrm{vol}(\mathcal T)\,O_\gamma(f)$ by unfolding the class kernel
over $G_\gamma(F)\backslash G(F)$ (\lref{GlobalTheory.T.classTerm_eq_volume_mul_orbitalIntegral}).

(c) Choose Haar measures $dt_w$ on the archimedean centralisers and, at finite $v$, the Haar measure $dt_v$ giving volume
$1$ to the compact open subgroup $G_\gamma(F_v)\cap K_v$; their restricted product $dt$ exists
(\lref{GlobalTheory.Orb.existsUnique_isCentProd}) and so does a fundamental domain $\mathcal T$
(\lref{GlobalTheory.AutomorphicL2Unitary.isFundamentalDomain_stdDom}). Fix $e\colon E[\gamma]^+\cong\prod_{i<k}K_i$; since $k\ge2$ the
factor $i=0$ exists, and by (a) $a=\mathrm{rank}\,\epsilon_0$ satisfies $0<a<n$. By $\multncst{Rad}$ there is $v\in S$ at
which $\tau_v$ has no non-zero vector fixed by $U_a(k_v)$, and by (a), $o_v=O_\gamma(f_v)=0$. The product $\prod_vo_v$ has
only finitely many factors different from $1$ and one factor $0$, so it vanishes, and by (b) the class integral of $f$ at
$\gamma$ is $\mathrm{vol}(\mathcal T)\prod_wo_w\cdot0=0$.
\end{proof}

Only the radical hypothesis, the restricted product property, the smoothness and compact support of the $f_w$ and the
properties of the finite test family are used; the oddness of $n$ and the pseudo-coefficient property of the $f_w$ are not.

For the remaining classes the vanishing comes from the archimedean place. We first isolate an abstract uniqueness
principle for invariant functionals supported on a regular semisimple orbit. Let $H$ be a locally compact second
countable group with a closed embedding into some $M_N(\C)$ and a Haar measure $dh$ which is also right invariant, and
let $s\in H$. A \emph{weight} for $s$ is a smooth function $W\ge0$ of two matrix variables with $W(s,1)>0$, with an
invariance property under translation by the centraliser $H_s$ and a properness property; weights exist when $H$ is the
unitary group of a non-degenerate hermitian form and $s$ is regular semisimple with proper orbit
(\lref{MainStatements.WeightKernel.weight}, from the diagonal case
\lref{MainStatements.WeightKernel.weightDiag}). A \emph{kernel functional} is $\Lambda(\phi)=\int_\Omega\sum_{i\in I}
c_i(x)\,\phi(q_i(x))\,dm(x)$ for a measure space $(\Omega,m)$, a countable $I$ and measurable $c_i,q_i$ such that, for each
$x$ and compact $C$, only finitely many $i$ have $c_i(x)\ne0$ and $q_i(x)\in C$, and the integrand is integrable for every
$\phi\in C_c^\infty(H)$.

\begin{lemma}[Invariant kernel functionals]\label{pf:multnc:vanish}
Let $s$, $W$ and $\Lambda$ be as above, and assume that every $q_i(x)$ is conjugate to $s$ and that
$\Lambda(\phi(k\cdot k^{-1}))=\Lambda(\phi)$ for $k\in H$ and $\phi\in C_c^\infty(H)$. If $\phi\in C_c^\infty(H)$ has orbital
integral $O_s(\phi)=0$, then $\Lambda(\phi)=0$.
\formallabel{MainStatements.WeightKernel.vanish}
\end{lemma}

\begin{proof}
We use three statements of the development: an interchange of $\Lambda$ with integration over a compact set of $k$, for
smooth families $y\mapsto\Psi(y,k)$ with uniformly compact support (\lref{MainStatements.WeightKernel.kerSwap}); for
every $\psi\ge0$ in $C_c(H)$ with $\psi(s)>0$ there is $a\ne0$ with $\int_HW(z^{-1}sz,k)\psi(k^{-1}z^{-1}szk)\,dk=a$ for all
$z\in H$ (\lref{MainStatements.WeightKernel.identOne}); and if $O_s(\phi)=0$ then $\int_H\phi(kyk^{-1})W(kyk^{-1},k)\,dk=0$
for every conjugate $y$ of $s$ (\lref{MainStatements.WeightKernel.identTwo}). Let $\psi$ be a smooth bump function
with $\psi(s)>0$ and $a$ as in the second statement. By the properness of $W$ there is a compact $\mathcal K\subset H$ such
that $W(y,k)\phi(y)\psi(k^{-1}yk)=0$ and $\phi(kyk^{-1})W(kyk^{-1},k)\psi(y)=0$ for all conjugates $y$ of $s$ and all
$k\notin\mathcal K$. Since every $q_i(x)$ is a conjugate of $s$,
\begin{align*}
a\,\Lambda(\phi)&=\Lambda\Bigl(y\mapsto\phi(y)\int_{\mathcal K}W(y,k)\psi(k^{-1}yk)\,dk\Bigr)
=\int_{\mathcal K}\Lambda\bigl(y\mapsto\phi(y)W(y,k)\psi(k^{-1}yk)\bigr)\,dk\\
&=\int_{\mathcal K}\Lambda\bigl(y\mapsto\phi(kyk^{-1})W(kyk^{-1},k)\psi(y)\bigr)\,dk\\
&=\Lambda\Bigl(y\mapsto\psi(y)\int_{\mathcal K}\phi(kyk^{-1})W(kyk^{-1},k)\,dk\Bigr)=\Lambda(0)=0,
\end{align*}
using the interchange twice and the conjugation invariance of $\Lambda$ in the third step (with the function
$y\mapsto\phi(y)W(y,k)\psi(k^{-1}yk)$ conjugated by $k$). As $a\ne0$, $\Lambda(\phi)=0$.
\end{proof}

\begin{lemma}[Archimedean vanishing]\label{pf:multnc:archzero}
Assume the conclusion of \cref{pf:multnc:regint}. Let $\mu$ be a Haar measure on $G(\A)$ which is also right invariant, $f_w$
smooth compactly supported
functions on the $G_w$ and $(f_v)$ a finite family such that $f'_\infty\otimes f_f$ is a cusp function for every family
$(f'_w)$ of smooth compactly supported functions. Let $\gamma\in G(F)$ be regular semisimple and $w$ an archimedean place
with $O_\gamma(f_w)=0$ for some Haar measure on the centraliser of $\gamma$ in $G_w$. Then the class integral of
$f=f_\infty\otimes f_f$ at $\gamma$ is $0$.
\formallabel{MainStatements.WeightKernel.archZero}
\end{lemma}

\begin{proof}
Let $\mathcal F$ be a fundamental domain. The class kernel converges absolutely, having only finitely many non-zero terms
near each point ($f$ has compact support and $G(F)$ is discrete in $G(\A)$), and it is integrable on $\mathcal F$ by
\cref{pf:multnc:regint}, since a left $G(F)$-invariant function which is integrable on the quotient is integrable on every
fundamental domain. For the vanishing
(\lref{MainStatements.WeightKernel.PArchZero.integral_class_eq_zero}) write $f(g)=R(g)\,f_w(g_w)$, where $R$ is the
product of the other local components, and consider, for $\phi\in C_c^\infty(G_w)$,
\[
\Lambda(\phi)=\int_{\mathcal F}\sum_{\delta\in[\gamma]}R(x^{-1}\delta x)\,\phi((x^{-1}\delta x)_w)\,dx .
\]
This is a kernel functional on $H=G_w$ (with $\Omega=\mathcal F$, $I=[\gamma]$, $c_\delta(x)=R(x^{-1}\delta x)$ and
$q_\delta(x)=(x^{-1}\delta x)_w$): the finiteness condition holds as $R\,\phi$ has compact support, and the integrability
holds by \cref{pf:multnc:regint}, since $R\,\phi$ is again a cusp function by hypothesis. Each
$q_\delta(x)=(gx)_w^{-1}\gamma_w(gx)_w$ ($\delta=g^{-1}\gamma g$) is a conjugate of $\gamma_w$. $\Lambda$ is invariant
under conjugation: replacing $\phi$ by $\phi(k\cdot k^{-1})$ amounts to translating $x$ on the right by $k^{-1}\in G_w$,
since $R$ is invariant under conjugation by $G_w$, and the integral over $\mathcal F$ of a left $G(F)$-invariant function
is invariant under right translations, $\mu$ being also right invariant. The group $G_w$ is the unitary group of a
non-degenerate hermitian matrix, its Haar measures are also right invariant, and $\gamma_w$ is regular semisimple with
proper orbit (\lref{GlobalTheory.Orb.isProperOrbit_GArch}); so $\gamma_w$ has a weight, and \cref{pf:multnc:vanish}
gives $\Lambda(f_w)=0$, which is the claim.
\end{proof}

\begin{lemma}[The remaining reducible classes]\label{pf:multnc:hyperzero}
Assume the conclusion of \cref{pf:multnc:regint} and the structure statement for regular semisimple elements of unitary
groups (\lref{GlobalTheory.T.regularSemisimple_unitaryGroup_concrete}). For every set of test data and every regular
semisimple $\gamma\in G(F)$ whose characteristic polynomial is reducible and which is not of non-field type, the class
integral of $f$ at $\gamma$ is $0$.
\formallabel{GlobalTheory.PsCo.Plug.hyperX}
\formallabel{GlobalTheory.PsCo.Plug.hyperZero_of_intX}
\end{lemma}

\begin{proof}
Fix a good Witt frame of $V$ (\lref{GlobalTheory.TruncKer.exists_isGood}). Since the characteristic polynomial of
$\gamma$ is reducible and $\gamma$ is not of non-field type, $\gamma$ is $G(F)$-conjugate to an element $\gamma_0$ of a proper
standard Levi subgroup (\lref{GlobalTheory.T.exists_conj_mem_standardLevi_of_reducible_charpoly}; this is where the
structure statement enters); $\gamma_0$ is regular semisimple and, at every archimedean place $w$, its centraliser in $G_w$
is commutative and not compact (\lref{GlobalTheory.T.centralizer_regular_levi_commutative_noncompact}). Fix $w$. Since
$f_w$ is a pseudo-coefficient (\cref{pf:multnc:def-pc}), $O_{\gamma_0}(f_w)=0$ for every Haar measure on the centraliser. For
every family $(f'_w)$ of smooth compactly supported functions, $f'_\infty\otimes f_f$ is a cusp function, because the type
datum of test data is cuspidal (\lref{MainStatements.WeightKernel.cuspAny}). By \cref{pf:multnc:archzero}
the class integral of $f$ at $\gamma_0$ is $0$; as $[\gamma]=[\gamma_0]$, the class kernels of $f$ at $\gamma$ and at
$\gamma_0$ coincide.
\end{proof}

\subsubsection*{The elliptic expansion}

\begin{proposition}[The elliptic expansion]\label{pf:multnc:ellexp}
Let $E/F$ be a CM extension unramified at all finite places, $J\in M_n(E)$ hermitian non-degenerate, $L$ an
$\OO_E$-lattice, $\mu$ the restricted product of Haar measures $\mu_w,\mu_v$, $f_w$ smooth compactly supported functions on
the $G_w$ and $(f_v)$ a finite family. For every plus polynomial $P$ of degree $n$ let $dt_P$ be the restricted product of
Haar measures $dt_{P,w},dt_{P,v}$ on $T_P$, and define $\mathrm{t}^{\mathrm{orb}}(P)$ by \eqref{pf:multnc:eq-termorb}. Then
for every $P$: $\mathrm{vol}(T_P(F)\backslash T_P(\A))<\infty$, every $\multncPhi_v(P;f_v)$ and every $\multncPhi_w(P;f_w)$
is well defined, and $\multncPhi_v(P;f_v)=1$ for almost all $v$; $\mathrm{t}^{\mathrm{orb}}(P)\ne0$ for only finitely many
$P$; $K^{\mathrm{ell}}_f$ is integrable on $G(F)\backslash G(\A)$; and $\int K^{\mathrm{ell}}_f=\sum_P
\mathrm{t}^{\mathrm{orb}}(P)$.
\formallabel{GlobalTheory.TrEllSum.elliptic_expansion}
\end{proposition}

\begin{proof}
For a plus polynomial $P$ let $K_{f,P}(x)=\sum_{\delta}f(x^{-1}\delta x)$, the sum over $\delta\in G(F)$ with
characteristic polynomial $P$. We use the following statements of the development.
\begin{enumerate}[label=(\alph*)]
\item (\lref{GlobalTheory.T.integral_charpolyKernel_eq_stableOrbitalIntegrals}) For every $P$, $K_{f,P}$ is
integrable on $G(F)\backslash G(\A)$ with integral $\mathrm{t}^{\mathrm{orb}}(P)$.
\item (\lref{GlobalTheory.TrFinite.finite_charpoly_and_ratPts}) For compact $C\subset G(\A)$, the set of characteristic
polynomials of the $\delta\in G(F)$ having a conjugate $x^{-1}\delta x\in C$ is finite, and for each $x$ only finitely many
$\delta$ satisfy $x^{-1}\delta x\in C$.
\item (\lref{GlobalTheory.T.exists_torusClass_of_irreducible_charpoly_unitary}) Every $\delta\in G(F)$ with irreducible
characteristic polynomial has a plus polynomial as characteristic polynomial (and a class in $X_P(F,J)$).
\item (\lref{GlobalTheory.T.unitaryTorus_compact_quotient_finite_volume}) $T_P(F)\backslash T_P(\A)$ is compact and has
finite volume.
\item (\lref{GlobalTheory.T.local_stable_orbital_integrals_well_defined}) For continuous compactly supported $f_v$, resp.\
$f_w$, $\multncPhi_v(P;f_v)$, resp.\ $\multncPhi_w(P;f_w)$, is well defined.
\item (\lref{GlobalTheory.T.eventually_unitaryTorus_orbitalClass_trivial_local}) For almost all $v$: the trivial class
lies in $X_P(F_v,J)$, and if $\mu_v(K_v)=1$ and $dt_{P,v}$ gives volume $1$ to the stabiliser of the standard lattice of
$K_P$, the orbital integral of $\one_{K_v}$ at $\gamma_q$ is $1$ for the trivial class $q$ and $0$ for the other classes.
\end{enumerate}
The finiteness and well-definedness assertions follow from (d), (e), and from (f) together with $f_v=\one_{K_v}$ and the
normalisations of $\mu_v$ and $dt_{P,v}$ at almost all $v$, which give $\multncPhi_v(P;f_v)=1$ there. Let $C=\supp f$ and let
$\mathcal S$ be the finite set of plus polynomials which are characteristic polynomials of elements $\delta\in G(F)$ with
$x^{-1}\delta x\in C$ for some $x$ (b). For every $x$, the sum defining $K^{\mathrm{ell}}_f(x)$ is finite by (b), and by (c),
and since plus polynomials are irreducible, grouping it according to the characteristic polynomial gives
$K^{\mathrm{ell}}_f=\sum_{P\in\mathcal S}K_{f,P}$. If $P\notin\mathcal S$, then $K_{f,P}=0$ and $\mathrm{t}^{\mathrm{orb}}(P)=
\int K_{f,P}=0$ by (a); so $\mathrm{t}^{\mathrm{orb}}$ is supported in $\mathcal S$. By (a), $K^{\mathrm{ell}}_f$ is integrable
as a finite sum of integrable functions, and $\int K^{\mathrm{ell}}_f=\sum_{P\in\mathcal S}\mathrm{t}^{\mathrm{orb}}(P)=\sum_P
\mathrm{t}^{\mathrm{orb}}(P)$.
\end{proof}

% ===========================================================================
\subsection{The normalising constant of the row counts: statement}\label{pf:multnc:sec-nc}

\paragraph{Setting.} In \cref{pf:multnc:sec-nc,pf:multnc:sec-ncpf,pf:multnc:sec-ncin,pf:multnc:sec-ncproofs} $D$ is again a global datum, with
totally real field $F$, CM field $E$ (unramified over $F$ at all finite places) and ranks $k$, $m$, $r_S$ and $r=k+m-1$,
where $m=2(r_S+1)$ and $m+3\le k$; thus $r=k+2r_S+1$. Let $f=[F:\Q]$, which is also the number of archimedean places of
$F$ and of $E$, and $d_F$ the discriminant of $F$. Let $g_0,g_1$ be the two auxiliary primes of the host of $D$ and
$g_S=g_0g_1$. A finite place $v$ of $F$ is \emph{typed} if $v\mid g_0g_1$ and \emph{ordinary} otherwise. For a word $u$ of
the row construction, $T(u)$ is its level prime.

A \emph{Hida setting} $X$ for $D$ consists of a large amount of data, of which we use: the hermitian space $(E^n,J)$ with
$n=r$ and its lattice $L$ (identified with the lattice of $D$); the unitary group $G=\Uni(J)$ over $F$; a hyperbolic
system $e_1,\dots,e_{r_S},f_1,\dots,f_{r_S}$ ($\langle e_a,e_b\rangle=\langle f_a,f_b\rangle=0$, $\langle e_a,f_b\rangle=
\delta_{ab}$) and, for each word $u$, a vector $z_u$ orthogonal to it with $\langle z_u,z_u\rangle=-T(u)$; the
\emph{frame} $P_u=(e,f,z_u)$ of $2r_S+1$ vectors and its stabiliser $H_u\subset G$; Haar measures on $G(\A)$,
$G(\multncAf)$, $G(F_v)$, $G(F\otimes\R)$, on $H_u(\multncAf)$ and on the archimedean groups $H_u(F\otimes\R)$ and
$N_u(F\otimes\R)$, where $N_u$ is the unipotent group attached to $(e,z_u)$; and the open compact subgroup
$K_f=\prod_vK_v$ of $G(\multncAf)$, $K_v$ the stabiliser of $L\otimes\OO_v$.

\paragraph{Local factors.} For a finite place $v$ of $F$ let $q_v=\#\OO_F/v$, and $\eta_v=1$ if $v$ splits in $E$,
$\eta_v=-1$ otherwise. For integers $k\le N$ and a natural number $T$ put
\begin{equation}\label{pf:multnc:eq-aP}
a_v(k,N,T)=\prod_{i=k_0}^{N}\bigl(1-\eta_v^iq_v^{-i}\bigr),\qquad k_0=\begin{cases}k&\text{if }v\mid T,\\k+1&\text{if }
v\nmid T,\end{cases}
\end{equation}
and $\mathrm{al}_v=a_v(k,n,T(u))$ if $v$ is typed, $\mathrm{al}_v=a_v(k,n,T(u))\,a_v(k,n,T(u'))$ otherwise. For $i\ge2$ let
\begin{equation}\label{pf:multnc:eq-Lval}
L_i=\prod_{\mathfrak p}\bigl(1-\eta_{\mathfrak p}^iN\mathfrak p^{-i}\bigr)^{-1},
\end{equation}
the product over the maximal ideals of $\OO_F$, with $\eta_{\mathfrak p}=-1$ if $\mathfrak p\OO_E$ is prime and $1$
otherwise; thus $L_i=L(i,\eta^i)$ for the quadratic character $\eta$ of $E/F$.

\paragraph{The mass datum.} A \emph{mass datum} indexed by a finite set $\mathcal P$ is a record of real numbers. We use the
following fields and derived quantities:
\[
\mu_K=\mu_{K,o}\prod_{p\in\mathcal P}\mu_{K,p},\quad \mu_K^\natural=\mu_{K,o}\prod_p\mu^\natural_{K,p},\quad
\mu_S=\mu_{S,o}\prod_p\mu_{S,p},\quad\mu_{S'}=\mu_{S',o}\prod_p\mu_{S',p},
\]
\[
\mathrm{MM}=M_o\prod_pM_p,\qquad U=U_o\prod_pU_p,\qquad \mathrm{AA}=\mathrm{vol}_B\,C_{oo}\,(\mathrm{MM}/\mathrm{vol}_{ZZ})/\mu_K,
\]
together with the fields $\mu_H,\mu_{H'},\mu_N,\mu_{N'}$, $\mathrm{vol}_{ZZ}$, $\mathrm{al}$, $\mathrm{vol}_G$,
$\mathrm{vol}_{Ha},\mathrm{vol}_{Ha'},\mathrm{vol}_{Na},\mathrm{vol}_{Na'}$, $\mathrm{vol}_B$, $\mathrm{vol}_H$,
$\mathrm{vol}_N,\mathrm{vol}_{N'}$, $C_{oo}$, $\rho_0$ and $n_{0c}$. The relations we use are:
\begin{enumerate}[label=(R\arabic*),start=0]
\item positivity of all the fields above (and of several others);
\item $n_{0c}\rho_0=\mathrm{AA}$;
\item $\rho_0=\mathrm{vol}_{Ha}/(\mathrm{vol}_H\mu_H)$, $\mathrm{vol}_{Na}=\mathrm{vol}_N\mu_N$,
$\mathrm{vol}_{Na'}=\mathrm{vol}_{N'}\mu_{N'}$ (and a primed analogue of the first);
\item $\mu_S=\mu_N\mu_H$ and $\mu_{S'}=\mu_{N'}\mu_{H'}$ (and two analogues);
\item $M_o/U_o=\mu_{K,o}^2/(\mu_{S,o}\mu_{S',o})$;
\item $M_p/U_p=\mu_{K,p}\mu^\natural_{K,p}/(\mu_{S,p}\mu_{S',p})$ for all $p$ (and two further relations);
\item $\mathrm{vol}_B\mathrm{vol}_G=\mathrm{vol}_{Ha}\mathrm{vol}_{Ha'}$ and $\mathrm{vol}_{Na}=\mathrm{vol}_{Na'}=1$;
\item $\mathrm{vol}_{Ha'}=\mathrm{vol}_G$.
\end{enumerate}

\begin{definition}[The mass datum of a Hida setting]\label{pf:multnc:def-bcf}
Let $X$ be a Hida setting for $D$, $u,u'$ words, $\mathrm{Fd}_N,\mathrm{Fd}_{N'}$ subsets of $N_u(F\otimes\R)$,
$N_{u'}(F\otimes\R)$, $X_0$ a frame and $\psi$ an additive character of $\A_F/F$ (together with further parameters which
play no role below). The mass datum $\multncB=\multncB(X,u,u',\dots)$ is indexed by the typed places, and its relevant
fields are: $\mu_{K,p}=\mu_{G,p}(K_p)$ and $\mu_{K,o}=\mu_{G,f}(K_f)/\prod_p\mu_{G,p}(K_p)$, so that $\mu_K=\mu_{G,f}(K_f)$;
$\mu^\natural_{K,p}$, the volume of the subgroup $K^\natural_p\subset K_p$ generated by the stabiliser of the frame and the
principal congruence subgroup of level $g_p$; $\mu_{H'}$, the volume of $K_f\cap H_{u'}(\multncAf)$ for the measure of
$H_{u'}(\multncAf)$; $\mathrm{vol}_H$, the volume of $H_u(F\otimes\R)$; $\mathrm{vol}_N$ and $\mathrm{vol}_{N'}$, the volumes
of $\mathrm{Fd}_N$ and $\mathrm{Fd}_{N'}$; $U_p=\mathrm{lvu}_p$ and $U_o=\prod_{v\text{ ordinary}}\mathrm{lvu}_v$,
$\mathrm{vol}_{ZZ}=\prod_v\mathrm{ZZ}_v$, where $\mathrm{lvu}_v$ and $\mathrm{ZZ}_v$ are explicit local volumes of sets of
unipotent coordinates and of boxes of Gram entries for the self-dual measures attached to $\psi_v$;
$\mathrm{al}=\prod_{v\text{ typed}}\mathrm{al}_v\cdot\prod_{v\text{ ordinary}}\mathrm{al}_v$; the covolumes
$\mathrm{vol}_G$ of $G(F)$ in $G(\A)$, $\mathrm{vol}_{Ha}$, $\mathrm{vol}_{Ha'}$ of the rational points in $H_u(\A)$,
$H_{u'}(\A)$, and $\mathrm{vol}_B$ in the stabiliser of $X_0$; and
\[
C_{oo}=\kappa_N(D)\cdot C_{\mathrm{dis}}(X,u,u'),\qquad n_{0c}=\mathrm{AA}/\rho_0,
\]
where $C_{\mathrm{dis}}(X,u,u')$ is the positive constant of an archimedean disintegration statement (and $0$ if there is
none) and $\kappa_N(D)\ne0$ is an explicit constant which cancels below. Let
$\nu_g(D)=\sqrt{|N_{E/\Q}(g_S^m)|}^{\,k-1}$.
\formal{Formalism.L2.SG.MassDataConc.bCF}
\formal{Formalism.L2.SG.MassDataConc.kapN}
\formal{Formalism.L2.SG.MassDataConc.ngFac}
\end{definition}

\begin{definition}[The explicit constant]\label{pf:multnc:def-cexpl}
Let $v_E=\frac12|d_E|^{1/(2f)}$, let $\lambda_E(i)$ be the Euler product of \cref{prop:families:mass} (which equals $L_i$ for
$i\ge2$, since $E/F$ is unramified at the finite places), and
\begin{equation}\label{pf:multnc:eq-mr}
\mathrm{mr}(i)=\frac{\bigl((i-1)!\,(v_E/\pi)^i\bigr)^{f}}{\sqrt{(2v_E)^{f}}}\,\lambda_E(i),
\end{equation}
the right-hand side of the mass ratio of \cref{prop:families:mass}. Let $c_{\mathrm{mass}}(j)=c_1\prod_{i=2}^{j}\mathrm{mr}(i)$,
where $c_1>0$ is the mass constant of unimodular hermitian lattices of rank $1$ over $E$. For a flag $F_0$ and a finite set
$\mathcal R$ of lattices let $\mathrm{mass}(\mathcal R)=\sum_{\Lambda\in\mathcal R}\#\mathrm{Aut}_{F_0}(\Lambda)^{-1}$,
where $\mathrm{Aut}_{F_0}(\Lambda)$ is the group of automorphisms of $\Lambda$ fixing the flag. The explicit constant is
\begin{equation}\label{pf:multnc:eq-cexpl}
c_{\mathrm{expl}}=\frac{c_{\mathrm{mass}}(k-1)\cdot|d_F|^{-2r_S(r_S+1)(k-1)}\cdot\prod_{i=k}^{r}L_i}{\mathrm{mass}(\mathcal R)}.
\end{equation}
\formal{Formalism.L2.SG.Pf.Cnt.cExpl}
\formal{Formalism.L2.SG.Pf.Cnt.cMass}
\end{definition}

\begin{hypothesis}[Hypotheses of the second theorem]\label{pf:multnc:hyp-nc}
We consider the following statements, in the formulation of \cite{LeanSource2026}.
\begin{enumerate}[label=(N\arabic*)]
\item The reference covolume identity: for every Hida setting, all words $u,u'$ and fundamental sets
$\mathrm{Fd}_N,\mathrm{Fd}_{N'}$ as in (N4), the product of the volumes of $\mathrm{Fd}_N$ and of $\mathrm{Fd}_{N'}$, both for
the product over the archimedean places of the reference measures attached to $u$, is $|d_F|^{e_{DF}}$, where
$e_{DF}=-(8r_S^3+6r_S^2+6kr_S^2+r_S)/3$. It is proved in the development
(\lref{Formalism.L2.SG.Closed.CovolRef_outright}, which establishes exactly this identity).
\item The Siegel--Weil input: for every Hida setting, fundamental domains of $G(F)$ in $G(\A)$ have finite volume, and for
$2a<n$, a non-degenerate hermitian $T\in M_a(E)$ and a factorizable Schwartz function $\Phi$ on $a$-tuples of vectors, the
mean over a fundamental domain of the theta series $\sum_{\mathrm{gram}(Y)=T}\Phi(gY)$ equals $g_\infty(T)\prod_vg_v(T)$,
where $g_\infty,g_v$ are the continuous densities of the push-forwards of the self-dual measures under the Gram map
\cite{Weil1965}.
\item Three auxiliary facts: the total mass of the gauge measure of the compact unitary group $\Uni(k)$ is
$\mathrm{vol}\,\Uni(k)$, with a gauge density; a limit formula for the volumes of small balls in a model space; and every
totally positive element of the maximal totally real subfield of a CM field which is unramified at the finite places is
represented by every binary form $as\bar s+bt\bar t$ with $a,b$ totally positive. They enter the formal statement as
instance hypotheses and are proved in the development (\lref{GlobalTheory.GaugeVolumeUD.gaugeVolUD},
\lref{GlobalTheory.RankOneChart.rhsWLim} and \lref{GlobalTheory.hermRepUnram}, which establish these three
facts).
\item Hypotheses on the data: measure ties of the Hida setting, a normalised level subgroup, and $\mathrm{Fd}_N$,
$\mathrm{Fd}_{N'}$ are fundamental sets for the integral points of $N_u$, $N_{u'}$; measure compatibility for $\psi$ (for
$P_u$, $P_{u'}$ and every frame $X_0$ with non-degenerate Gram matrix, the measures of the setting disintegrate along the
orbit map $g\mapsto gP$ against the self-dual orbit measures, and the $N_u$-measures have covolume $1$); $E/F$ is CM, $J$ is
hermitian non-degenerate, and $\psi$ is a character of $\A_F/F$ which is standard at infinity; $L$ is an $\OO_E$-lattice;
the Gram matrices of $P_{u'}$ and of $X_0$ are non-degenerate; positivity of the local factors $\mathrm{lvu}_v$,
$\mathrm{ZZ}_v$, $\mathrm{al}_v$ and of further local factors, and convergence and positivity of their products; local
measure ties at $u$ and $u'$ and facts at the typed places; positivity of four quotient volumes; $s_d\ge0$, $C\ge0$ and
$T_h>0$ for three real parameters of the record; the congruence $T(u)\equiv T(u')\pmod{g_S^2}$; and $\mathcal R$ is a set of
row representatives for $F_0$: every row lattice for $F_0$ is the image of exactly one element of $\mathcal R$ under the
stabiliser of $F_0$.
\end{enumerate}
\formal{Formalism.L2.SG.Pf.NcReal.JX.CovolRef}
\formal{PI.VolStabS.MassDataS.SWXS}
\end{hypothesis}

\begin{theorem}[Explicit normalising constant]\label{pf:multnc:main-nc}
Assume (N1)--(N4). Then for the mass datum $\multncB$ of \cref{pf:multnc:def-bcf},
\begin{equation}\label{pf:multnc:eq-nctiec}
c_{\mathrm{expl}}=\frac{n_{0c}\cdot\nu_g(D)}{\mathrm{al}\cdot\kappa_N(D)}.
\end{equation}
\formallabel{Formalism.L2.SG.Pf.NcAllC.ncTieC_of_closed}
\end{theorem}

The theorem is used to show that all normalising constants of the row counts are positive and to identify the mean count
constant with the unscaled main constant on the row-counting side.

% ===========================================================================
\subsection{Proof of the explicit normalising constant}\label{pf:multnc:sec-ncpf}

\subsubsection*{Step 1: algebraic reduction}

\begin{lemma}\label{pf:multnc:ncalg}
Let $B$ be a mass datum satisfying (R0)--(R5) and $\mathrm{vol}_{Na}=\mathrm{vol}_{Na'}=1$. Then
\begin{equation}\label{pf:multnc:eq-n0c}
n_{0c}=\frac{\mathrm{vol}_B}{\mathrm{vol}_{Ha}}\,C_{oo}\,\mathrm{vol}_H\,\mathrm{vol}_N\,\mathrm{vol}_{N'}\,
\frac{U}{\mathrm{vol}_{ZZ}}\,\frac{\mu_K^\natural}{\mu_{H'}},
\end{equation}
and for all $c$, $\nu$ and $\kappa\ne0$: $c=n_{0c}\,\nu/(\mathrm{al}\,\kappa)$ if and only if $c\,\mathrm{al}\,\kappa$
equals $\nu$ times the right-hand side of \eqref{pf:multnc:eq-n0c}.
\formallabel{Formalism.L2.SG.Pf.NcCut.normalizing_constant_eq}
\formallabel{Formalism.L2.SG.Pf.NcCut.nc_iff}
\end{lemma}

\begin{proof}
From (R4) and (R5), $M_o=U_o\mu_{K,o}^2/(\mu_{S,o}\mu_{S',o})$ and $M_p=U_p\mu_{K,p}\mu^\natural_{K,p}/(\mu_{S,p}\mu_{S',p})$;
multiplying, $\mathrm{MM}=U\mu_K\mu_K^\natural/(\mu_S\mu_{S'})$. By (R1) and the definition of $\mathrm{AA}$,
$n_{0c}=\mathrm{vol}_BC_{oo}\mathrm{MM}/(\mathrm{vol}_{ZZ}\mu_K\rho_0)$. Substitute $\mathrm{MM}$, then $\mu_S=\mu_N\mu_H$ and
$\mu_{S'}=\mu_{N'}\mu_{H'}$ by (R3), $\rho_0=\mathrm{vol}_{Ha}/(\mathrm{vol}_H\mu_H)$ and $\mu_N=1/\mathrm{vol}_N$,
$\mu_{N'}=1/\mathrm{vol}_{N'}$ by (R2) and $\mathrm{vol}_{Na}=\mathrm{vol}_{Na'}=1$. All quantities involved are positive by
(R0), and \eqref{pf:multnc:eq-n0c} follows by cancelling. The equivalence follows since $\mathrm{al}>0$ and $\kappa\ne0$.
\end{proof}

\begin{lemma}[The record relations]\label{pf:multnc:record}
Under (N2) and (N4) the mass datum $\multncB$ satisfies (R0)--(R7); moreover $C_{oo}=\kappa_N(D)\,C_{\mathrm{dis}}(X,u,u')$
and $\kappa_N(D)\ne0$.
\formallabel{Formalism.L2.SG.Pf.MassDataX.pos_bCF}
\formallabel{Formalism.L2.SG.Pf.MassDataX.inMass_bCF}
\formallabel{Formalism.L2.SG.Pf.MassDataX.inMasses_bCF}
\formallabel{Formalism.L2.SG.Pf.MassDataX.inLevelS_bCF}
\formallabel{Formalism.L2.SG.Pf.MassDataX.inTyped_bCF}
\formallabel{Formalism.L2.SG.Pf.MassDataX.inVolumes_bCF}
\formallabel{Formalism.L2.SG.Pf.MassDataX.volHa'_eq}
\end{lemma}

\begin{proof}
The identity $C_{oo}=\kappa_N(D)\,C_{\mathrm{dis}}$ holds by the definition of $\multncB$, and $\kappa_N(D)$ is a power of a
positive real number times the inverse of a power of $g_S$, hence non-zero. (R1) holds by the definition
$n_{0c}=\mathrm{AA}/\rho_0$, because $\rho_0$ is the mass $\sum_c\#(\Gamma\cap cKc^{-1})^{-1}$ of a non-empty finite
double coset space $\Gamma\backslash H_u(\multncAf)/K$ with finite stabilisers, hence positive
(\lref{Formalism.L2.SG.Pf.MassDataX.massH_pos}, which proves this positivity). (R4) holds by the definition of $M_o$ in $\multncB$ and the positivity of $U_o$
(\lref{Formalism.L2.SG.MassDataConc.inOrdinary_bCF}, which records this deduction). The remaining relations are
consequences of statements proved in the development whose proofs are deferred: (R0) from the positivity of Haar measures
of open sets, of the covolumes and of the local factors (\lref{Formalism.L2.SG.Pf.MassDataPos.massDataPosC_of},
which assembles these positivity statements); (R2) and (R3) from the product structure of the measures of the stabilisers
of the frames (\lref{Formalism.L2.SG.Pf.MassDataX.inMassesAt_holds} and
\lref{level_s_at_instance}, which prove these identities for every mass datum of the
form of \cref{pf:multnc:def-bcf}); (R5) from the local self-dual measures at the typed places
(\lref{Formalism.L2.SG.Pf.MassDataX.inTypedC_loc}, which proves the typed relations from the local ties); and (R6),
(R7) from the Siegel--Weil input (N2) and the measure compatibility (\lref{volumes_s_at_instance}
and \lref{Formalism.L2.SG.Pf.MassDataX.covolH_eq_volG}, which give the covolume relations and
$\mathrm{vol}_{Ha'}=\mathrm{vol}_G$).
\end{proof}

\begin{definition}[The value identity]\label{pf:multnc:def-ncvalue}
The \emph{value identity} \multncst{V} for $\multncB$ is
\begin{equation}\label{pf:multnc:eq-ncvalue}
c_{\mathrm{expl}}\cdot\mathrm{al}=C_{\mathrm{dis}}\cdot\mathrm{vol}_H\,\mathrm{vol}_N\,\mathrm{vol}_{N'}\cdot
\frac{U}{\mathrm{vol}_{ZZ}}\cdot\frac{\mu_K^\natural}{\mu_{H'}}\cdot\nu_g(D).
\end{equation}
\formallabel{Formalism.L2.SG.Pf.NcAsm.NcValueC}
\end{definition}

\begin{proposition}\label{pf:multnc:ncvalue-suffices}
If $\multncB$ satisfies the relations of \cref{pf:multnc:record} and \multncst{V}, then \eqref{pf:multnc:eq-nctiec} holds.
\formallabel{Formalism.L2.SG.Pf.NcAsm.ncTieC_of'}
\formallabel{Formalism.L2.SG.Pf.NcAsm.ncTieC_of}
\end{proposition}

\begin{proof}
By (R6) and (R7), $\mathrm{vol}_B\mathrm{vol}_G=\mathrm{vol}_{Ha}\mathrm{vol}_G$, hence $\mathrm{vol}_B=\mathrm{vol}_{Ha}$
since $\mathrm{vol}_G>0$. By \cref{pf:multnc:ncalg} with $\kappa=\kappa_N(D)$ and $C_{oo}=\kappa\,C_{\mathrm{dis}}$, the claim
is equivalent to $c_{\mathrm{expl}}\,\mathrm{al}\,\kappa=\kappa\cdot(\text{right-hand side of \eqref{pf:multnc:eq-ncvalue}})$,
which is $\kappa$ times \eqref{pf:multnc:eq-ncvalue}.
\end{proof}


\subsubsection*{Step 2: from local volumes to global identities}

Write $X_e=\frac{(r_S+1)(8r_S^2+10r_S+3)}{3}$, $e_2=n(2r_S+1)-\frac{(2r_S+1)^2}2$, $|\cdot|_v$ for the normalised absolute
value of $F_v$, $\mathfrak d_v$ for a generator of the local different of $F$ at $v$, and $g_v=g_0$ if $v\mid g_0$,
$g_v=g_1$ otherwise. Let $\omega_v$ be the local orbit measure on $(2r_S+1)$-tuples of vectors in $F_v\otimes E^n$ attached
to the Gram matrix of $P_{u'}$ (the fibre over that Gram matrix of the self-dual measure for $\psi_v$), and
\[
\Omega=\prod_v\omega_v(K_v\cdot P_{u'}).
\]

\begin{definition}[Local and field identities]\label{pf:multnc:def-localids}
Assume that $\psi$ is a character of $\A_F/F$, standard at infinity. We consider:
\begin{enumerate}[label=(L\arabic*)]
\item (ordinary volumes) $\mathrm{lvu}_v/\mathrm{ZZ}_v=|\mathfrak d_v|_v^{-X_e}$ for every ordinary $v$;
\item (typed volumes) if $T(u)\equiv T(u')\pmod{g_S^2}$, then $\mathrm{lvu}_v/\mathrm{ZZ}_v=g_v^{(r_S+1)(2r_S+3)}$ for every
typed $v$;
\item (typed indices) $\mu^\natural_{K,p}/\mu_{K,p}=g_p^{-(r_S+1)(2k+2r_S+1)}/a_p(k,n,T(u))$ for every typed $p$;
\item (orbit volumes) $\omega_v(K_v\cdot P_{u'})=|\mathfrak d_v|_v^{e_2}\,a_v(k,n,T(u'))$ for every finite $v$;
\item (different) $\prod_v|\mathfrak d_v|_v=|d_F|^{-1}$;
\item (field dictionary) the discriminant and the number of archimedean places of the base field of $X$ are $d_F$ and $f$;
for $i\ge2$, $\prod_v(1-\eta_v^iq_v^{-i})^{-1}=L_i$; $\prod_v\mathrm{tf}_v=P_T$, where $\mathrm{tf}_v=1-\eta_v^kq_v^{-k}$ if
$v\mid T(u)$ and $1$ otherwise, and $P_T=\prod_{\mathfrak p\mid T(u)}(1-\eta_{\mathfrak p}^kN\mathfrak p^{-k})$;
$v_E=|d_F|^{1/f}/2$; $E$ has $f$ archimedean places; $n=r$; $|\mathfrak d_p|_p=1$ and $a_p(k,n,T(u'))=a_p(k,n,T(u))$ for
typed $p$; and $\prod_{p\text{ typed}}g_p=g_S^{f}$;
\item (row mass) if $T(u)\equiv T(u')\pmod{g_S^2}$ and $\mathcal R$ is a set of row representatives for $F_0$, then
$\mathrm{mass}(\mathcal R)=c_{\mathrm{mass}}(k)\,T(u)^{kf}\,P_T$.
\end{enumerate}
\formallabel{Formalism.L2.SG.Pf.NcIdx.IndexTypC}
\formallabel{Formalism.L2.SG.Pf.NcVal.OmV}
\formallabel{Formalism.L2.SG.Pf.NcVal.DiffProd}
\formallabel{Formalism.L2.SG.Pf.NcVal.FieldDict}
\formallabel{Formalism.L2.SG.Pf.Cnt.RowsMass}
\end{definition}

\begin{lemma}\label{pf:multnc:different}
Identity (L5) holds.
\formallabel{Formalism.L2.SG.Pf.NcVal.diffProd_holds}
\end{lemma}

\begin{proof}
Let $\mathfrak D$ be the different of $F/\Q$. For each finite $v$, $\mathfrak d_v$ generates $\mathfrak D\OO_v$, so
$|\mathfrak d_v|_v=q_v^{-\ord_v\mathfrak D}$, which is $1$ for almost all $v$. Hence
$\prod_v|\mathfrak d_v|_v=N(\mathfrak D)^{-1}=|d_F|^{-1}$, since the absolute norm of the different is the absolute value
of the discriminant \cite[III.2.9]{Neukirch1999}.
\end{proof}

\begin{lemma}\label{pf:multnc:localids-hold}
For the data of \cref{pf:multnc:main-nc}, identities (L1)--(L7) hold.
\end{lemma}

\begin{proof}
(L1) and (L2) are \lref{Formalism.L2.SG.Pf.NcVal.volOrdV_holds} and \lref{Formalism.L2.SG.Pf.NcVal.volTypV_holds},
which state exactly these local volume identities. (L3) is \lref{Formalism.L2.SG.Pf.AlTieIdx.indexTypC_holds}, a
computation of the index of $K^\natural_p$ in $K_p$ (deferred). (L4) is \cref{pf:multnc:omv}. (L5) is
\cref{pf:multnc:different}. (L6) is \lref{Formalism.L2.SG.Pf.AlTieLG.fieldDict_holds}
(deferred); it collects identifications of the fields of the Hida setting with those of $D$, the Euler product of
$L(i,\eta^i)$ (absolutely convergent for $i\ge2$), the identity $v_E=|d_F|^{1/f}/2$ (as $E/F$ is unramified at all finite
places, $|d_E|=d_F^2$) and three facts about the typed places. (L7) is \cref{pf:multnc:rowmass} below.
\end{proof}

Put
\begin{gather*}
A_T=\prod_va_v(k,n,T(u)),\qquad A'=\prod_va_v(k,n,T(u')),\\
A_{\mathrm{typ}}=\prod_{p\text{ typed}}a_p(k,n,T(u)),\qquad\Lambda=\prod_{i=k+1}^{r}L_i .
\end{gather*}

\begin{proposition}[Global identities]\label{pf:multnc:global}
Assume (L1)--(L6), the positivity and convergence hypotheses of (N4) and $T(u)\equiv T(u')\pmod{g_S^2}$. Then the products
above converge, $\Lambda,P_T,A_T,A',A_{\mathrm{typ}}>0$, and
\begin{align}
&\frac{U}{\mathrm{vol}_{ZZ}}=|d_F|^{X_e}\,g_S^{(r_S+1)(2r_S+3)f},\label{pf:multnc:eq-G1}\\
&\frac{\mu^\natural_K}{\mu_K}=\frac{g_S^{-(r_S+1)(2k+2r_S+1)f}}{A_{\mathrm{typ}}},\label{pf:multnc:eq-G2}\\
&\Omega=|d_F|^{-e_2}A',\label{pf:multnc:eq-G3}\\
&\mathrm{al}\cdot A_{\mathrm{typ}}=A_TA',\label{pf:multnc:eq-G4}\\
&A_T\,\Lambda=P_T,\label{pf:multnc:eq-G5}\\
&\nu_g(D)=g_S^{mf(k-1)}.\label{pf:multnc:eq-G6}
\end{align}
\end{proposition}

\begin{proof}
We use repeatedly that a product of convergent infinite products converges to the product of the values, and that a
product over all places can be split into the finite product over the typed places and the product over the ordinary places.

\eqref{pf:multnc:eq-G5}: since $n=r\ge k+1$, \eqref{pf:multnc:eq-aP} gives
$a_v(k,n,T)=\prod_{i=k+1}^{r}(1-\eta_v^iq_v^{-i})\cdot\mathrm{tf}_v$. For each $i\in[k+1,r]$ we have $i\ge2$, so by (L6)
$\prod_v(1-\eta_v^iq_v^{-i})=L_i^{-1}$, with $L_i>0$ (\lref{Formalism.HermitianLattices.Ext.L3.lValEta_pos}, positivity
of these Euler products); and $\prod_v\mathrm{tf}_v=P_T$ is a finite product of positive factors. Hence $A_T=\Lambda^{-1}P_T$;
similarly $A'=\Lambda^{-1}P'_T$ with $P'_T$ the analogous finite product for $T(u')$.

\eqref{pf:multnc:eq-G4}: by definition $\mathrm{al}=\prod_{p\text{ typed}}a_p(T(u))\prod_{v\text{ ordinary}}a_v(T(u))
a_v(T(u'))$. Splitting the convergent product $A_TA'=\prod_va_v(T(u))a_v(T(u'))$ along the typed places and using
$a_p(T(u'))=a_p(T(u))$ at typed $p$ (L6), $\prod_{v\text{ ordinary}}a_v(T(u))a_v(T(u'))=A_TA'/A_{\mathrm{typ}}^2$, whence
$\mathrm{al}=A_{\mathrm{typ}}\cdot A_TA'/A_{\mathrm{typ}}^2$.

\eqref{pf:multnc:eq-G3}: by (L4), $\Omega=\prod_v|\mathfrak d_v|_v^{e_2}\cdot\prod_va_v(T(u'))=(|d_F|^{-1})^{e_2}A'$ by (L5)
(raising the convergent product of (L5) to the real power $e_2$).

\eqref{pf:multnc:eq-G1}: by the definitions, $U=\prod_v\mathrm{lvu}_v$ and $\mathrm{vol}_{ZZ}=\prod_v\mathrm{ZZ}_v$, both
convergent with positive values (N4). By (L1), (L2) and $|\mathfrak d_v|_v=1$ at typed $v$ (L6),
$\mathrm{lvu}_v/\mathrm{ZZ}_v=|\mathfrak d_v|_v^{-X_e}\cdot c_v$ for every $v$, where $c_v=g_v^{(r_S+1)(2r_S+3)}$ if $v$ is
typed and $c_v=1$ otherwise. Taking products and using (L5) and $\prod_{p\text{ typed}}g_p=g_S^f$ gives
$U/\mathrm{vol}_{ZZ}=|d_F|^{X_e}g_S^{(r_S+1)(2r_S+3)f}$.

\eqref{pf:multnc:eq-G2}: $\mu_K^\natural/\mu_K=\prod_p\mu^\natural_{K,p}/\mu_{K,p}$ (the factor $\mu_{K,o}>0$ cancels); by (L3)
this is $\prod_pg_p^{-(r_S+1)(2k+2r_S+1)}/A_{\mathrm{typ}}=g_S^{-(r_S+1)(2k+2r_S+1)f}/A_{\mathrm{typ}}$.

\eqref{pf:multnc:eq-G6}: $g_S^m\in\Q$, so $N_{E/\Q}(g_S^m)=g_S^{m[E:\Q]}=g_S^{2mf}$ and
$\sqrt{|N_{E/\Q}(g_S^m)|}^{\,k-1}=g_S^{mf(k-1)}$.
\end{proof}

\subsubsection*{Step 3: the final computation}

\begin{proposition}\label{pf:multnc:final}
Let $a'>0$ and assume the conclusions of \cref{pf:multnc:global}, (L7), (R0) for $\multncB$, the third fact of (N3), and
\begin{align}
&\mu_K\,a'=\mu_{H'}\,\Omega,\label{pf:multnc:eq-S}\\
&C_{\mathrm{dis}}\cdot\mathrm{vol}_H\,\mathrm{vol}_N\,\mathrm{vol}_{N'}=a'\,2^{kf}\Bigl(\frac{\pi^k}{(k-1)!}\Bigr)^{f}
T(u)^{-kf}\,|d_F|^{e_{DF}}.\label{pf:multnc:eq-R}
\end{align}
Then the value identity \multncst{V} holds.
\formallabel{Formalism.L2.SG.Pf.NcArithG.ncValueC_of_global}
\end{proposition}

\begin{proof}
Write $d=|d_F|$, $s=r_S$, $g=g_S$, $T=T(u)$. We use: $c_{\mathrm{mass}}(k)=c_{\mathrm{mass}}(k-1)\,\mathrm{mr}(k)$ for $k\ge2$
(immediate from the definition); $\lambda_E(k)=L_k>0$ (\lref{Formalism.HermitianLattices.Ext.L3.lValEta_eq_lValEtaF},
which identifies the two Euler products for CM fields unramified at the finite places, applicable by
\lref{Formalism.L2.SG.Pf.Cnt.field_dictionary_holds}, which supplies these hypotheses for $D$); $c_{\mathrm{mass}}(k-1)>0$
(\lref{Formalism.L2.SG.Pf.Asm.cMass_spec}, positivity of the mass constants); and, from $v_E=d^{1/f}/2$,
$\sqrt{(2v_E)^f}=d^{1/2}$ and $v_E^{kf}=d^k/2^{kf}$. Hence by \eqref{pf:multnc:eq-mr}
\[
\mathrm{mr}(k)=\frac{((k-1)!)^f\,d^k}{2^{kf}\pi^{kf}\,d^{1/2}}\,L_k .
\]
\emph{Left-hand side of \eqref{pf:multnc:eq-ncvalue}.} By \eqref{pf:multnc:eq-cexpl}, (L7), \eqref{pf:multnc:eq-G4} and
\eqref{pf:multnc:eq-G5}, writing $\prod_{i=k}^rL_i=L_k\Lambda$,
\[
c_{\mathrm{expl}}\,\mathrm{al}=\frac{c_{\mathrm{mass}}(k-1)\,d^{-2s(s+1)(k-1)}L_k\Lambda}{c_{\mathrm{mass}}(k-1)\,\mathrm{mr}(k)\,
T^{kf}P_T}\cdot\frac{A_TA'}{A_{\mathrm{typ}}}
=d^{E_1}\,2^{kf}\Bigl(\frac{\pi^k}{(k-1)!}\Bigr)^{f}T^{-kf}\,\frac{A'}{A_{\mathrm{typ}}},
\]
with $E_1=-2s(s+1)(k-1)+\frac12-k$.

\emph{Right-hand side.} By \eqref{pf:multnc:eq-S}, $\mu_K/\mu_{H'}=\Omega/a'$, so
$\mu_K^\natural/\mu_{H'}=(\mu_K^\natural/\mu_K)(\Omega/a')$. Using \eqref{pf:multnc:eq-R}, \eqref{pf:multnc:eq-G1},
\eqref{pf:multnc:eq-G2}, \eqref{pf:multnc:eq-G3} and \eqref{pf:multnc:eq-G6}, the right-hand side of
\eqref{pf:multnc:eq-ncvalue} equals
\[
2^{kf}\Bigl(\frac{\pi^k}{(k-1)!}\Bigr)^{f}T^{-kf}\,d^{\,e_{DF}+X_e-e_2}\,
g^{(s+1)(2s+3)f+mf(k-1)-(s+1)(2k+2s+1)f}\,\frac{A'}{A_{\mathrm{typ}}}.
\]
Since $m=2(s+1)$, the exponent of $g$ is $(s+1)f\bigl((2s+3)+2(k-1)-(2k+2s+1)\bigr)=0$. Since $r=k+2s+1$,
$e_2=(k+2s+1)(2s+1)-\frac{(2s+1)^2}2=2ks+k+2s^2+2s+\frac12$, so
\[
E_1+e_2-X_e=\bigl(-2s^2k+4s^2+4s+1\bigr)-\frac{8s^3+18s^2+13s+3}{3}=-\frac{8s^3+6s^2+6ks^2+s}{3}=e_{DF}.
\]
Hence $e_{DF}+X_e-e_2=E_1$, and the two sides agree.
\end{proof}

\begin{corollary}\label{pf:multnc:arith}
Let $a'>0$. Assume \eqref{pf:multnc:eq-S}, \eqref{pf:multnc:eq-R}, (L1)--(L7), the positivity and convergence hypotheses of
(N4), (R0), the congruence $T(u)\equiv T(u')\pmod{g_S^2}$ and that $\mathcal R$ is a set of row representatives. Then
\multncst{V} holds.
\formallabel{Formalism.L2.SG.Pf.NcGlob.normalizing_arithmetic_holds}
\end{corollary}

\begin{proof}
Combine \cref{pf:multnc:global} and \cref{pf:multnc:final}.
\end{proof}

\subsubsection*{Step 4: assembly}

The identities \eqref{pf:multnc:eq-S} and \eqref{pf:multnc:eq-R} come from the following two statements.

\begin{definition}[The splitting identity]\label{pf:multnc:def-splitom}
The statement \multncst{Split} is: let $E/F$ be a CM extension, $J\in M_n(E)$ hermitian non-degenerate, $\psi$ a character of
$\A_F/F$ standard at infinity, $P$ a frame of $2r+1$ vectors with non-degenerate Gram matrix and $2(2r+1)<n$, and fix
normalised measures on $G=\Uni(J)$ and on the stabiliser $H_P$, and an orbit measure $\nu_H$ on tuples which is a fibre
family for the Gram map over the Gram matrix of $P$ and through which the measure of $G(\A)$ disintegrates against that of
$H_P(\A)$ along $g\mapsto gP$; let $\Lambda$ be an $\OO_E$-lattice with $0<\mu_{G,f}(K_f)<\infty$, and assume the two parts
of the Siegel--Weil input (N2). Then there is $a'>0$ such that (i) the measure of $G(F\otimes\R)$ disintegrates as $a'$
times the integral, over the archimedean orbit of $P$ with the archimedean self-dual fibre measures, of the measure of
$H_P(F\otimes\R)$ (we say that \emph{$a'$ is an archimedean disintegration constant for $P$}), and (ii)
$\mu_{G,f}(K_f)\,a'=\mu_{H_P,f}(K_f\cap H_P(\multncAf))\,\prod_v\omega_v(K_v\cdot P)$.
\formallabel{Formalism.L2.SG.Pf.NcSplitS.SplitOm}
\end{definition}

\begin{definition}[The real bracket identity]\label{pf:multnc:def-realbracket}
The statement \multncst{RB} is: for every Hida setting $X$, words $u,u'$, fundamental sets $\mathrm{Fd}_N,\mathrm{Fd}_{N'}$
and $\psi$ standard at infinity satisfying the hypotheses (N4) on measures, and every archimedean disintegration constant
$a'>0$ for $P_{u'}$, identity \eqref{pf:multnc:eq-R} holds.
\formallabel{Formalism.L2.SG.Pf.NcReal.RealBracket}
\end{definition}

\begin{theorem}\label{pf:multnc:splitom}
The splitting identity \multncst{Split} holds.
\formallabel{Formalism.L2.SG.Pf.NcS.SplitPf.splitOm_of_intWitt}
\end{theorem}

\begin{proof}
By \cref{pf:multnc:splitwitt}, \multncst{Split} follows from the local Witt theorem, which holds by
\lref{Formalism.L2.SG.Pf.NcS.WittLocPf.local_witt_theorem} (proved via a Witt splitting,
\lref{Formalism.L2.SG.Pf.NcS.LT.WSp.witt_split}; compare \cite{Jacobowitz1962}), and the integral Witt theorem almost
everywhere, which holds by \lref{Formalism.L2.SG.Pf.NcS.IntWittPf.integral_witt_almost_everywhere}.
\end{proof}

\begin{theorem}\label{pf:multnc:realbracket}
The reference covolume identity (N1), together with the first two facts of (N3), implies \multncst{RB}.
\formallabel{Formalism.L2.SG.Closed.RealBracket_holds}
\end{theorem}

\begin{proof}
By \cref{pf:multnc:jacgauge} (with the second fact of (N3)) \multncst{Jac} holds, and by \cref{pf:multnc:aprime} (with the
first fact of (N3)) the archimedean disintegration constants are as stated there. Gauge densities exist for all signatures
(\lref{Formalism.MainNotions.S.Gauge.gauge_density_exists}). Hence \cref{pf:multnc:rb-jac} applies.
\end{proof}

\begin{proposition}\label{pf:multnc:nctieall}
Assume \multncst{Split}, \multncst{RB}, (N2), (N4) and the third fact of (N3). Then \eqref{pf:multnc:eq-nctiec} holds.
\formallabel{Formalism.L2.SG.Pf.NcAllS.ncTieAllS_of}
\formallabel{Formalism.L2.SG.Pf.NcAllS.ncTie0All_of}
\formallabel{Formalism.L2.SG.Pf.NcLeaves.ncTieC_of_leaves}
\end{proposition}

\begin{proof}
By \cref{pf:multnc:record} the relations (R0)--(R7) hold. The measure compatibility for $u'$ provides an orbit measure for
$P_{u'}$ with the fibre and disintegration properties required in \multncst{Split}. We have $n=r=k+2r_S+1$ and
$k\ge m+3=2r_S+5$, hence $2(2r_S+1)<n$. By (R0), $\mu_{K,o}>0$ and $\mu_{K,p}>0$, hence $0<\mu_{G,f}(K_f)<\infty$. So
\multncst{Split} gives an archimedean disintegration constant $a'>0$ for $P_{u'}$ with $\mu_{G,f}(K_f)a'=\mu_{H'}\Omega$; since
$\mu_K=\mu_{K,o}\prod_p\mu_{K,p}=\mu_{G,f}(K_f)$ by definition, this is \eqref{pf:multnc:eq-S}. By \multncst{RB} with this
$a'$ we get \eqref{pf:multnc:eq-R}. The identities (L1)--(L7) hold by \cref{pf:multnc:localids-hold}. By
\cref{pf:multnc:arith}, \multncst{V} holds, and by \cref{pf:multnc:ncvalue-suffices}, \eqref{pf:multnc:eq-nctiec} holds.
\end{proof}

\begin{proof}[Proof of \cref{pf:multnc:main-nc}]
By \cref{pf:multnc:realbracket}, (N1) and (N3) give \multncst{RB}; \multncst{Split} holds by \cref{pf:multnc:splitom};
conclude by \cref{pf:multnc:nctieall}.
\end{proof}

% ===========================================================================
\subsection{The row-lattice mass formula}\label{pf:multnc:sec-ncin}

We prove the row-lattice mass formula (L7) from a mass formula for glued lattices. The proofs of (L4), of
\cref{pf:multnc:splitom} and of \cref{pf:multnc:realbracket} are completed in \cref{pf:multnc:sec-ncproofs}.

\begin{proposition}\label{pf:multnc:rowmass}
Identity (L7) holds.
\formallabel{PI.UnitaryGroupMassFormula.CntFormula.RowsMassW.row_lattice_mass}
\end{proposition}

\begin{proof}
Let $T=T(u)$. To each row lattice $\Lambda$ for the flag $F_0$ the development attaches a lattice $M(\Lambda)$ in a hermitian
space $V_1$ which contains a vector $z$ with $\langle z,z\rangle=-T$ and whose orthogonal complement $z^\perp$ is positive
definite of dimension $k$ (\lref{PI.UnitaryGroupMassFormula.CntFormula.RowsMassW.anisotropic_complement_definite},
which proves this definiteness and computes the dimension). The lattice $M(\Lambda)$ is a \emph{glued} lattice: a unimodular
lattice in $V_1$ containing $z$ as a primitive vector. Two facts relate row lattices and glued lattices
(\lref{PI.UnitaryGroupMassFormula.CntFormula.RowsMassW.row_class_holds}): $\Lambda_1,\Lambda_2$ are conjugate under the
stabiliser of $F_0$ if and only if $M(\Lambda_1),M(\Lambda_2)$ are conjugate under the stabiliser of $z$; and every glued
lattice is of the form $M(\Lambda)$ (\cref{pf:multnc:rowsurj}). Since
$\mathcal R$ is a set of row representatives, it follows that $\Lambda\mapsto M(\Lambda)$ is injective on $\mathcal R$ and
that its image is a set of representatives of the glued lattices modulo the stabiliser of $z$; here $M(\Lambda)$ is a glued
lattice for every row lattice $\Lambda$ (\lref{PI.UnitaryGroupMassFormula.CntFormula.RowsMassW.row_lattice_mass_holds},
which proves this). Moreover $\#\mathrm{Aut}_{F_0}(\Lambda)=\#\mathrm{Stab}_z(M(\Lambda))$
(\lref{PI.UnitaryGroupMassFormula.CntFormula.RowsMassW.row_weight_formula}, which identifies the two automorphism
groups). Glued lattices exist, since the image of the standard unimodular lattice is one
(\lref{PI.UnitaryGroupMassFormula.CntFormula.RowsMassW.GluedLatticeMass.glued_exists}, existence of glued lattices
when $\dim z^\perp\ge2$ and a unimodular lattice exists). The mass formula for glued lattices
(\lref{PI.UnitaryGroupMassFormula.CntFormula.RowsMassW.GluedLatticeMass.glued_lattice_mass_formula}) states that, for a set
of representatives $\mathcal M$ of the glued lattices modulo the stabiliser of $z$ with $\dim z^\perp=k\ge1$,
\[
\sum_{M\in\mathcal M}\frac1{\#\mathrm{Stab}_z(M)}=c_{\mathrm{mass}}(k)\,T^{kf}\prod_{\mathfrak p\mid T}
\bigl(1-\eta_{\mathfrak p}^kN\mathfrak p^{-k}\bigr).
\]
Hence $\mathrm{mass}(\mathcal R)=\sum_{\Lambda\in\mathcal R}\#\mathrm{Stab}_z(M(\Lambda))^{-1}=c_{\mathrm{mass}}(k)\,T^{kf}P_T$.
\end{proof}

The glued lattice mass formula (about $270$ Lean lines, deferred) is deduced from the mass constants $c_{\mathrm{mass}}$, that
is, from the mass formula for unimodular hermitian lattices in the form of the mass ratio of \cref{prop:families:mass},
applied in the positive definite space $z^\perp$; it uses that two unimodular lattices in the same non-degenerate hermitian space lie in the same
genus (\lref{PI.UnitaryGroupMassFormula.CntFormula.UG.unimodular_one_genus}, which proves this over CM fields unramified
at the finite places).

% ===========================================================================
\subsection{Proofs of the analytic inputs of the normalising constant}\label{pf:multnc:sec-ncproofs}

We prove the orbit volume identity (L4), the splitting identity \multncst{Split}, the real bracket identity
\multncst{RB} and the surjectivity statement used for (L7). Throughout, $E/F$ is a CM extension, $J\in M_n(E)$ hermitian and
non-degenerate, $G=\Uni(J)$, and $\psi$ a character of $\A_F/F$ which is standard at infinity; we say that the
\emph{standard hypotheses} hold. For an $F$-algebra $R$ (a completion $F_v$, $F\otimes\R$ or $\multncAf$) and $a\ge1$ let
$\mathrm{gram}(Y)=(\langle Y_i,Y_j\rangle)_{i,j}$ be the Gram matrix of an $a$-tuple $Y$ of vectors in $R\otimes_FE^n$, and
let $dY$ and $dT$ be the self-dual measures for $\psi_R$ on $a$-tuples and on hermitian $a\times a$ matrices.

\subsubsection*{Fibre measures and local orbit volumes}

A \emph{fibre family} for the Gram map is a family $(\omega_T)_T$ of measures on $a$-tuples, indexed by the hermitian
$a\times a$ matrices $T$, such that: $\omega_T$ is concentrated on the linearly independent tuples $Y$ with
$\mathrm{gram}(Y)=T$; $\int_{\mathcal U}\phi\,dY=\int\!\int\phi\,d\omega_T\,dT$ for every measurable $\phi\ge0$, where
$\mathcal U$ is the open set of linearly independent tuples; and for every continuous compactly supported $\phi$ with support
in $\mathcal U$, $\phi$ is $\omega_T$-integrable and $T\mapsto\int\phi\,d\omega_T$ is continuous.

\begin{lemma}[Fibre measures]\label{pf:multnc:omunique}
Under the standard hypotheses, for every finite place $v$ and every $a$, a fibre family for the Gram map on $a$-tuples
over $F_v$, with respect to the self-dual measures, exists and is unique.
\formallabel{PI.OmFibreUniq.W.omUnique_of}
\formallabel{PI.OmFibreUniq.W.omUnique_of_sdEx}
\formallabel{omUnique}
\end{lemma}

\begin{proof}
Since $\psi_v$ is non-trivial (\lref{PI.OmFibreUniq.W.Sd.additive_character_nontrivial}), the self-dual measures $dY$
and $dT$ exist and are Haar measures (\lref{PI.OmFibreUniq.W.SdInst.sdEx_of_psiNT}). \emph{Existence}
(\lref{PI.OmFibreUniq.W.gram_map_fibre_existence}): for Haar measures on tuples and on hermitian matrices, the
push-forward of $dY|_{\mathcal U}$ under the Gram map has a continuous density with respect to $dT$
(\lref{PI.OmFibreUniq.W.densAll_gen}), and a general disintegration statement for maps with continuous densities on
totally disconnected spaces produces the fibre family (\lref{PI.VolStabSW.DiscreteGroups.fibOfDens_holds}).
\emph{Uniqueness} (\lref{PI.OmFibreUniq.W.gram_map_fibre_uniqueness}, an application of the general statement
\lref{PI.OmFibreUniq.W.fib_unique}): let $(\omega_T)$ and $(\omega'_T)$ be fibre families. For a compact open
$C\subset\mathcal U$ the functions $T\mapsto\omega_T(C)$ and $T\mapsto\omega'_T(C)$ are continuous (the indicator of $C$ is
continuous with compact support); applying the disintegration identity to $\one_C\cdot(\one_A\circ\mathrm{gram})$ for
measurable sets $A$ of matrices shows that they agree $dT$-almost everywhere, hence everywhere, since $dT$ is a Haar measure
and so charges every non-empty open set. As $\mathcal U$ is open in a finite-dimensional space over the local field $F_v$,
the compact open subsets of $\mathcal U$ form a family stable under finite intersections which generates the Borel sets of
$\mathcal U$ and covers it by countably many members; both measures are concentrated on $\mathcal U$, hence
$\omega_T=\omega'_T$ for all $T$.
\end{proof}

By \cref{pf:multnc:omunique}, the local orbit measure $\omega_v$ in (L4) is well defined.

\begin{proposition}[Local orbit volumes]\label{pf:multnc:omv}
For every Hida setting $X$, every word $u'$ and every $\psi$, identity (L4) holds: if $\psi$ is a character of $\A_F/F$
standard at infinity, then $\omega_v(K_v\cdot P_{u'})=|\mathfrak d_v|_v^{e_2}\,a_v(k,n,T(u'))$ for every finite place $v$.
\formallabel{Formalism.L2.SG.Pf.NcX.OmVSt.omV_of}
\formallabel{Formalism.L2.SG.Pf.NcX.OmVSt.omVZ_holdsX}
\formallabel{Formalism.L2.SG.Pf.NcVal.OmVSt.OmVZ.omV_holdsX}
\end{proposition}

\begin{proof}
Let $a=2r_S+1$; since $n=r=k+2r_S+1$, we have $a\le n$, $n-a=k$ and $na-a^2/2=e_2$. Fix a finite place $v$ and let $Y_v$ be
the tuple obtained from $P_{u'}$ by multiplying its vectors by the local scalars of the development. We use:
\begin{enumerate}[label=(\alph*)]
\item the standard hypotheses hold, $L$ is a lattice, $L\otimes\OO_v$ is self-dual, and $E/F$ is unramified at $v$
(\lref{Formalism.L2.SG.Pf.NcX.OmVSt.omega_standard_holds}, \lref{Formalism.L2.SG.Pf.NcX.OmVSt.omega_lattice_holds},
\lref{Formalism.L2.SG.Pf.NcVal.OmVSt.OmX.omVSD_holdsX}, \lref{Formalism.L2.SG.Pf.NcX.OmVSt.omega_unramified_holds});
\item $\mathrm{gram}(Y_v)=\mathrm{gram}(P_{u'})$ (\lref{Formalism.L2.SG.Pf.NcX.OmVSt.omVZGram}), $Y_v$ is primitive in
$L\otimes\OO_v$ (\lref{Formalism.L2.SG.Pf.NcVal.OmVSt.OmX.omVPrim_holdsX}), and $\omega_v(K_v\cdot P_{u'})=
\omega_v(K_v\cdot Y_v)$ (\lref{Formalism.L2.SG.Pf.NcVal.OmVSt.OmVc.omVScale_holdsX});
\item if $v\nmid T(u')$, the determinant of $\mathrm{gram}(P_{u'})$ is a unit at $v$
(\lref{Formalism.L2.SG.Pf.NcX.OmVSt.omVZDet});
\item (local unitary orbit volumes; \lref{PI.OmLocU.Done.omLocU_closed}) under the standard hypotheses, for a
self-dual lattice at a place $v$ unramified in $E$ and a primitive $a$-tuple $Y$ ($a\le n$) whose Gram matrix has unit
determinant, $\omega_v(K_v\cdot Y)=|\mathfrak d_v|_v^{na-a^2/2}\prod_{j=n-a+1}^{n}(1-\eta_v^jq_v^{-j})$;
\item (places dividing the level; \lref{Formalism.L2.SG.Pf.NcVal.OmVSt.OmVZ.omVBadT_holdsX}) if $v\mid T(u')$, then
$\omega_v(K_v\cdot Y_v)=|\mathfrak d_v|_v^{na-a^2/2}\prod_{j=n-a}^{n}(1-\eta_v^jq_v^{-j})$.
\end{enumerate}
If $v\nmid T(u')$, (a)--(d) give $\omega_v(K_v\cdot P_{u'})=|\mathfrak d_v|_v^{e_2}\prod_{j=k+1}^n(1-\eta_v^jq_v^{-j})$, and
the product is $a_v(k,n,T(u'))$ by \eqref{pf:multnc:eq-aP}, since $v\nmid T(u')$. If $v\mid T(u')$, (b) and (e) give the
product from $j=k$ to $n$, which is again $a_v(k,n,T(u'))$.
\end{proof}

The two cited local volume computations carry the weight of (L4): (d) is a count of primitive representations over the
residue field combined with Hensel's lemma (about $1\,500$ Lean lines), and (e) a similar computation at the places
dividing $T(u')$, where the Gram matrix is not unimodular (about $2\,900$ lines).

\subsubsection*{The archimedean constant}

Let $X$ be a Hida setting, $u,u'$ words and $\mathrm{Fd}_N,\mathrm{Fd}_{N'}$, $\psi$ as in (N4). Write $T_{u'}$ for the Gram
matrix of $P_{u'}$, $\mathrm{orb}(g)=g\cdot P_{u'}$ for the orbit map on $G(F\otimes\R)$, $\mu_{G,\infty}$ and $\mu_{H,\infty}$
for the measures of $G(F\otimes\R)$ and $H_{u'}(F\otimes\R)$ of the setting, and let
\[
\kappa=c_{\mathrm{gauge}}\cdot\mathrm{vol}(\Uni(k))^f\,T(u')^{-kf},\qquad a_0=\kappa/\mathrm{vol}(H_{u'}(F\otimes\R)),
\]
where $c_{\mathrm{gauge}}=\prod_wc_{\mathrm{gauge}}(\eps_w)$ is the product of the gauge constants of the archimedean
places. We use the following statements of the development.
\begin{enumerate}[label=(A\arabic*)]
\item $H_{u'}(F\otimes\R)$ is compact (\lref{Formalism.L2.SG.Pf.MassDataX.hinfCompact_holds}); hence its measure is
positive and finite and is the one used in (RB) (\lref{PI.AprimeInfW.htotInf_of}, \lref{PI.AprimeInfW.hinfFin_of}).
\item Archimedean self-dual measures and fibre families of the Gram map exist (\lref{PI.AprimeInfW.exInf_holds}).
\item (Weil's formula; \lref{PI.AprimeInfW.weilInf_of}) For measurable $f\ge0$ on $G(F\otimes\R)$ and every map $s$
with $s(y)\cdot P_{u'}=y$ on the orbit of $P_{u'}$,
\[
\int f\,d\mu_{G,\infty}=\frac1{\mathrm{vol}(H_{u'}(F\otimes\R))}\int\!\int f(s(y)h)\,d\mu_{H,\infty}(h)\,
d(\mathrm{orb}_*\mu_{G,\infty})(y).
\]
\item (Push-forward; \lref{PI.AprimeInfW.pushInf_of_pushLFin_of_orbUniq_of_constTest}) For every archimedean fibre
family, $\mathrm{orb}_*\mu_{G,\infty}=\kappa\,\omega_{\infty,T_{u'}}$.
\item There is a measurable $C\subset G(F\otimes\R)$ with $0<\mu_{G,\infty}(C)<\infty$ (\lref{PI.AprimeInfW.ginfTest_holds}).
\end{enumerate}
The push-forward statement (A4) is the archimedean orbit computation. Both measures in (A4) are invariant under
$G(F\otimes\R)$, locally finite and concentrated on the orbit, and such measures are proportional
(\lref{PI.AprimeInfW.pushLFin_of_compact}, \lref{PI.AprimeInfW.orbUniq_of_compact}); the constant is identified as
$\kappa$ by evaluating both sides on one test function (\lref{PI.AprimeInfW.constTest_of_polar_compact}), through
polar coordinates near $P_{u'}$ and a Jacobian identity in charts (\lref{PI.AprimeInfW.jacOneC_of_charts}), which rests
on the normalisation of the measure of $G(F\otimes\R)$ by the gauge densities and the gauge volume of $\Uni(k)$
(\lref{PI.AprimeInfW.grpStd_full}) and on the form of the self-dual measures on tuples in these charts
(\lref{PI.AprimeInfW.tupStd_of_sm}).

\begin{lemma}[Value of the archimedean constant]\label{pf:multnc:ainf}
Assume (A1)--(A4). Then $a_0$ is an archimedean disintegration constant for $P_{u'}$. If moreover (A5) holds, then every
archimedean disintegration constant $a'>0$ for $P_{u'}$ equals $a_0$.
\formallabel{PI.AprimeInfW.ainfValue_of}
\formallabel{PI.AprimeInfW.aprimeInf_of}
\formallabel{PI.AprimeInfW.aprimeInf_of_steps}
\end{lemma}

\begin{proof}
By (A2) self-dual measures and fibre families exist. Let $(\omega_{\infty,T})$ be a fibre family, $f\ge0$ measurable and $s$
a section of the orbit map. By (A3), (A4) and (A1),
\begin{align*}
\int f\,d\mu_{G,\infty}&=\frac{\kappa}{\mathrm{vol}(H_{u'}(F\otimes\R))}\int\!\int f(s(y)h)\,d\mu_{H,\infty}(h)\,
d\omega_{\infty,T_{u'}}(y)\\
&=a_0\int\!\int f(s(y)h)\,d\mu_{H,\infty}(h)\,d\omega_{\infty,T_{u'}}(y),
\end{align*}
so $a_0$ is an archimedean disintegration constant. Let $a'>0$ be another one and $C$ as in (A5). Applying both identities
to $f=\one_C$, with the same fibre family and section, gives $\mu_{G,\infty}(C)=a'I=a_0I$ with
$I=\int\!\int\one_C(s(y)h)\,d\mu_{H,\infty}\,d\omega_{\infty,T_{u'}}$. Since $0<\mu_{G,\infty}(C)<\infty$ and $a',a_0>0$,
$0<I<\infty$, and $a'=a_0$.
\end{proof}

\begin{theorem}[The archimedean constant]\label{pf:multnc:aprime}
Assume the first fact of (N3) (the gauge volume of $\Uni(k)$). Let $X,u,u',\mathrm{Fd}_N,\mathrm{Fd}_{N'},\psi$ satisfy the
hypotheses of (N4), and assume gauge densities at all archimedean places. Then every archimedean disintegration constant
$a'>0$ for $P_{u'}$ equals $c_{\mathrm{gauge}}\cdot\mathrm{vol}(\Uni(k))^f\,T(u')^{-kf}/\mathrm{vol}(H_{u'}(F\otimes\R))$.
\formallabel{PI.AprimeInfAsm.aprime_infinite_real}
\formallabel{Formalism.L2.SG.Closed.AprimeInf_holds}
\end{theorem}

\begin{proof}
The first fact of (N3) is supplied by \lref{GlobalTheory.GaugeVolumeUD.gaugeVolUD}; it enters (A4) through
\lref{PI.AprimeInfW.grpStd_full}. Statements (A1)--(A5) hold as cited above, and \cref{pf:multnc:ainf} gives the claim.
\end{proof}

\subsubsection*{The Jacobian identity for gauge measures}

The real bracket is computed place by place in a model space. For $k\ge1$ and $r\ge0$ let $H_{k,r}=\C^{2r+k+1}$ with the
hermitian form of the development, $\mathcal G=\Uni(H_{k,r})$ with its \emph{gauge measure} $\mu$ (the Haar measure
normalised by the gauge density), $W_b\subset H_{k,r}$ the $(k+1)$-dimensional complement of the hyperbolic part spanned by
$e_1,\dots,e_r,f_1,\dots,f_r$, and $\Uni(W_b)\subset\mathcal G$ with its gauge measure $\mu_b$. For $z\in W_b$ with
$T(z)=-\langle z,z\rangle>0$ let $\mathcal W_z\subset\Uni(W_b)$ be the stabiliser of $z$ (a compact group), $\mathcal N$ the
unipotent radical of the stabiliser of the isotropic flag $(e_a)$, with the Haar measure $\nu_\kappa$ which is $\kappa$
times the measure normalised by its Lie algebra, and $\mathcal S_z=\mathcal N\mathcal W_z$ with the product measure
$\mu_{\mathcal S_z}$ of $\nu_\kappa$ and a Haar measure $\mu_{\mathcal W_z}$. Let $\mathrm{bc}(c,w)$ ($c\in\C^r$,
$w\in\Uni(W_b)$) be the big cell representatives, $\zeta\mapsto\sigma(\zeta)$ a section of the double cosets
$\mathcal W_z\backslash\Uni(W_b)/\mathcal W_{z'}$, and $\mathrm{wt}(c,\zeta)=\mathrm{wt}_1(c)\,\mathrm{wt}_2(\zeta)$ the weight
of the development.

\begin{definition}[The Jacobian identity]\label{pf:multnc:def-jac}
The statement \multncst{Jac} is: for $k\ge1$, $r\ge0$, $z,z'\in W_b$ with $T=T(z)>0$, $T'=T(z')>0$, Haar measures
$\mu_{\mathcal W_z},\mu_{\mathcal W_{z'}}$ of total masses $a,a'$, $\kappa>0$, a section $\sigma$ and measurable
$\phi\ge0$ on $\mathcal G$,
\begin{gather*}
\int\phi\,d\mu=c_{\mathrm{Jac}}\int\mathrm{wt}(c,\zeta)\int_{\mathcal S_z}\int_{\mathcal S_{z'}}\phi\bigl(s\,
\mathrm{bc}(c,\sigma(\zeta))\,s'^{-1}\bigr)\,d\mu_{\mathcal S_{z'}}(s')\,d\mu_{\mathcal S_z}(s)\,dc\,d\zeta,\\
c_{\mathrm{Jac}}=\frac{2^{r+1}\mathrm{vol}(\Uni(k))^2}{(TT')^k\,\mathrm{vol}(\Uni(k-1))\,a\,a'\,\kappa^2}.
\end{gather*}
\formallabel{Formalism.L2.SG.Pf.JacGauge}
\end{definition}

\begin{proposition}[The Jacobian identity]\label{pf:multnc:jacgauge}
Assume the second fact of (N3) (the limit formula for small balls). Then \multncst{Jac} holds.
\formallabel{Formalism.L2.SG.Pf.Jac.jacGauge_of}
\formallabel{Formalism.L2.SG.Pf.jacGauge_of}
\formallabel{Formalism.L2.SG.Pf.jacGauge}
\formallabel{Formalism.L2.SG.Closed.JacGauge_holds}
\end{proposition}

\begin{proof}
The second fact of (N3) is supplied by \lref{GlobalTheory.RankOneChart.rhsWLim}. We use three integration formulas.
\begin{enumerate}[label=(\alph*)]
\item (Bruhat coordinates) For gauge measures $\mu$ on $\mathcal G$ and $\mu_b$ on $\Uni(W_b)$,
\[
\int\phi\,d\mu=2^r\int\mathrm{wt}_1(c)\int_{\mathcal N}\int_{\Uni(W_b)}\int_{\mathcal N}\phi\bigl(n\,\mathrm{bc}(c,1)\,m\,
n'^{-1}\bigr)\,d\nu_1(n')\,d\mu_b(m)\,d\nu_1(n)\,dc .
\]
It follows from the sandwich estimate (\lref{Formalism.L2.SG.Pf.BH.sandwich'}: for small $\rho$, the point with Bruhat
coordinates $(n,c,m,n')$ lies in the gauge ball of radius $\rho$ only if an explicit quadratic expression in the coordinates
is $<((1+\epsilon)\rho)^2$, and does if it is $<((1-\epsilon)\rho)^2$) and the volume limit of gauge balls
(\lref{Formalism.L2.SG.Pf.BH.VL.volLim}): the right-hand side defines a Haar measure on $\mathcal G$
(\lref{Formalism.L2.SG.Pf.GE.rhsHaar}), hence a multiple of $\mu$, and the multiple is found by comparing the two
measures on small balls (\lref{Formalism.L2.SG.Pf.bruhatHaar_of_sandwich_volLim}).
\item (Double cosets of $\Uni(W_b)$) With $c_W=2\,\mathrm{vol}(\Uni(k))^2/((TT')^k\mathrm{vol}(\Uni(k-1))aa')$,
\[
\int\psi\,d\mu_b=c_W\int\mathrm{wt}_2(\zeta)\int_{\mathcal W_z}\int_{\mathcal W_{z'}}\psi\bigl(h\,\sigma(\zeta)\,h'^{-1}\bigr)\,
d\mu_{\mathcal W_{z'}}(h')\,d\mu_{\mathcal W_z}(h)\,d\zeta ;
\]
it follows in the same way from the limit formula for small balls
(\lref{Formalism.L2.SG.Pf.wbHaar_of_rhsWLim}).
\item Gauge measures on $\mathcal G$ and on $\Uni(W_b)$ exist (\lref{Formalism.L2.SG.Pf.gaugeExists_of_sandwich_rhsWLim},
using \lref{Formalism.L2.SG.Pf.GE.existsGaugeNormalizedStabilizer}).
\end{enumerate}
Take $\mu_b$ from (c). Substitute (b), applied to $m\mapsto\phi(n\,\mathrm{bc}(c,1)\,m\,n'^{-1})$, into (a). Regrouping the
integrations over $\mathcal N$ and $\mathcal W_z$, resp.\ $\mathcal N$ and $\mathcal W_{z'}$, into integrations over
$\mathcal S_z$, resp.\ $\mathcal S_{z'}$ (\lref{Formalism.L2.SG.Pf.Jac.lhs_c}), and noting that
$\mu_{\mathcal S_z}\otimes\mu_{\mathcal S_{z'}}$ carries the factor $\kappa^2$ relative to $\nu_1$
(\lref{Formalism.L2.SG.Pf.Jac.rhs_eq}), one obtains the identity of \multncst{Jac} with the weight
$\mathrm{wt}_1(c)\mathrm{wt}_2(\zeta)=\mathrm{wt}(c,\zeta)$ (\lref{Formalism.FourierDecayAlt.P.weight_factorization})
and the constant $2^rc_W/\kappa^2=c_{\mathrm{Jac}}$, both being
$2^{r+1}\mathrm{vol}(\Uni(k))^2/((TT')^k\mathrm{vol}(\Uni(k-1))aa'\kappa^2)$.
\end{proof}

\subsubsection*{The real bracket}

\begin{lemma}[Place-by-place splitting]\label{pf:multnc:placesplit}
Assume \multncst{Jac} and (N1). Let $X,u,u',\mathrm{Fd}_N,\mathrm{Fd}_{N'},\psi$ satisfy (N4), with gauge densities at all
archimedean places. Then there are, for every archimedean place $w$, Haar measures $\mu_{N,w},\mu_{N',w},\mu_{H,w},
\mu_{H',w}$ on the local groups and positive numbers $a_w,b_w,\kappa_w,\kappa'_w$, and real numbers $c_N,c_{N'}$, such that:
the measures of $N_u,N_{u'},H_u,H_{u'}$ over $F\otimes\R$ are the products of the $\mu_{N,w}$, $\mu_{N',w}$, $\mu_{H,w}$,
$\mu_{H',w}$; $\mu_{H,w}$ and $\mu_{H',w}$ have total masses $a_w$ and $b_w$; at every $w$ the local disintegration of the
gauge measure along the local double cosets holds with the constant
\begin{equation}\label{pf:multnc:eq-Cw}
C_w=\frac{c_{\mathrm{gauge}}(\eps_w)\,\mathrm{vol}(\Uni(k))^2}{(T(u)T(u'))^k\,\mathrm{vol}(\Uni(k-1))\,a_wb_w\,\kappa_w
\kappa'_w};
\end{equation}
$\mathrm{vol}_N=(\prod_w\kappa_w)c_N$, $\mathrm{vol}_{N'}=(\prod_w\kappa'_w)c_{N'}$; and $c_Nc_{N'}=|d_F|^{e_{DF}}$.
\formallabel{Formalism.L2.SG.Pf.NcReal.JX.placeSplit_of}
\end{lemma}

\begin{proof}
The measures of $N_u$ and $N_{u'}$ are products of local Haar measures up to global constants
(\lref{Formalism.L2.SG.RealObj.InfCut.ProductExactness.muN_eq}), and the local factors are multiples
$d_w\cdot\mathrm{ref}_{N,w}$, $d'_w\cdot\mathrm{ref}_{N,w}$ of the reference measures attached to $u$
(\lref{Formalism.L2.SG.Pf.NcReal.JX.muNv_eq}); similarly for $H_u$ and $H_{u'}$
(\lref{Formalism.L2.SG.RealObj.InfCut.ProductExactness.muH_eq}). Fix an archimedean place $w_0$ and put the global
constants at $w_0$: $\kappa_w$ is $d_w$ times the global constant of $N_u$ if $w=w_0$ and $d_w$ otherwise, similarly
$\kappa'_w$, $\mu_{N,w}=\kappa_w\mathrm{ref}_{N,w}$, $\mu_{N',w}=\kappa'_w\mathrm{ref}_{N,w}$, and $\mu_{H,w},\mu_{H',w}$
the local factors of $H_u,H_{u'}$ with the global constants put at $w_0$; let $a_w,b_w$ be their total masses (positive and
finite, \lref{Formalism.L2.SG.Pf.NcReal.JX.muHv_mass}). Then the product decompositions hold, and
$\mathrm{vol}_N=(\prod_w\kappa_w)c_N$, $\mathrm{vol}_{N'}=(\prod_w\kappa'_w)c_{N'}$ with $c_N$, $c_{N'}$ the volumes of
$\mathrm{Fd}_N$ and $\mathrm{Fd}_{N'}$ for the product of the reference measures attached to $u$; by (N1),
$c_Nc_{N'}=|d_F|^{e_{DF}}$. The local disintegration at $w$ with the constant \eqref{pf:multnc:eq-Cw} follows from
\multncst{Jac} by transporting it to the local frame at $w$ (\lref{Formalism.L2.SG.Pf.NcReal.JX.placeStep}): the matrix
which maps the standard frame of the model space to $(e,f,w_f,T(u)^{-1/2}z_u)$, $w_f$ an orthonormal basis of the
complement, identifies the local groups and measures with those of $\mathcal G$, and $z_{u'}$ with a vector of $W_b$.
\end{proof}

\begin{theorem}[The real bracket identity]\label{pf:multnc:rb-jac}
Assume \multncst{Jac}, (N1), the conclusion of \cref{pf:multnc:aprime} and gauge densities for all signatures. Then
\multncst{RB} holds.
\formallabel{Formalism.L2.SG.Pf.NcReal.realBracket_of}
\formallabel{Formalism.L2.SG.Pf.NcReal.JX.realBracket_of_jac}
\end{theorem}

\begin{proof}
Let $X,u,u',\mathrm{Fd}_N,\mathrm{Fd}_{N'},\psi$ satisfy (N4) and let $a'>0$ be an archimedean disintegration constant for
$P_{u'}$. Take the data of \cref{pf:multnc:placesplit}. By the product exactness of local disintegrations
(\lref{Formalism.L2.SG.RealObj.InfCut.ProductExactness.ProdExact_holds}: if the measures of $N_u,N_{u'},H_u,H_{u'}$ are
products of local measures and the local disintegrations hold with positive constants $C_w$, then the archimedean
disintegration statement of the setting holds with the constant $\prod_wC_w$), and since the constant of that statement is
unique (\lref{Formalism.L2.SG.Pf.NcReal.cooOf_eq}, by evaluation on a compact set of positive finite measure),
$C_{\mathrm{dis}}(X,u,u')=\prod_wC_w$. Write $V_k=\mathrm{vol}(\Uni(k))$, $T=T(u)$, $T'=T(u')$. We have
$\mathrm{vol}_H=\prod_wa_w$ and $\mathrm{vol}(H_{u'}(F\otimes\R))=\prod_wb_w$, and by \eqref{pf:multnc:eq-Cw} and
$c_{\mathrm{gauge}}=\prod_wc_{\mathrm{gauge}}(\eps_w)$,
\[
\prod_wC_w\,a_wb_w\,\kappa_w\kappa'_w=c_{\mathrm{gauge}}\Bigl(\frac{V_k^2}{(TT')^kV_{k-1}}\Bigr)^f .
\]
Since $V_j=\prod_{i<j}(2\pi)^{i+1}/i!$, we have $V_k/V_{k-1}=2^k\pi^k/(k-1)!$, so
\[
\frac{V_k^2}{(TT')^kV_{k-1}}=V_k\,2^k\,\frac{\pi^k}{(k-1)!}\,T^{-k}T'^{-k}.
\]
By \cref{pf:multnc:aprime}, $a'=c_{\mathrm{gauge}}V_k^fT'^{-kf}/\prod_wb_w$. Therefore
\begin{align*}
C_{\mathrm{dis}}\,\mathrm{vol}_H\,\mathrm{vol}_N\,\mathrm{vol}_{N'}&=\prod_wC_w\cdot\prod_wa_w\cdot\Bigl(\prod_w\kappa_w\Bigr)c_N
\cdot\Bigl(\prod_w\kappa'_w\Bigr)c_{N'}\\
&=\frac{c_{\mathrm{gauge}}}{\prod_wb_w}\Bigl(\frac{V_k^2}{(TT')^kV_{k-1}}\Bigr)^f|d_F|^{e_{DF}}\\
&=\frac{c_{\mathrm{gauge}}V_k^fT'^{-kf}}{\prod_wb_w}\,2^{kf}\Bigl(\frac{\pi^k}{(k-1)!}\Bigr)^fT^{-kf}|d_F|^{e_{DF}}\\
&=a'\,2^{kf}\Bigl(\frac{\pi^k}{(k-1)!}\Bigr)^fT^{-kf}|d_F|^{e_{DF}},
\end{align*}
which is \eqref{pf:multnc:eq-R}.
\end{proof}

\subsubsection*{The splitting identity}

Write $G(\A)=G(F\otimes\R)\times G(\multncAf)$, $K_f=\prod_vK_v$, and let $P$ be a frame as in \multncst{Split}, with Gram
matrix $T_P$, orbit measure $\nu_H$ on adelic tuples and stabiliser $H_P$. Let $\omega_{\infty,T}$ be an archimedean fibre
family. We use the following statements of the development.
\begin{enumerate}[label=(W\arabic*)]
\item (\lref{Formalism.L2.SG.Pf.NcS.AdTupPf.admissible_tuple_holds}) Measure-theoretic facts on the decomposition
$G(\A)=G(F\otimes\R)\times G(\multncAf)$, compatible with the action on tuples; $K_f$ is measurable; and, if
$0<\mu_{G,f}(K_f)<\infty$, then $0<\mu_{H,f}(K_f\cap H_P(\multncAf))<\infty$ (\lref{Formalism.L2.SG.Pf.NcS.AdTupPf.kf_compact_open}).
\item (\lref{Formalism.L2.SG.Pf.NcS.OmDecompPf.omega_decomposition_holds}) There is a measure $\omega_f$ on finite
adelic tuples with $\int A(y_\infty)\one_S(y_f)\,d\nu_H(y)=\omega_f(S)\int A\,d\omega_{\infty,T_P}$ for measurable $A\ge0$ and
$S$.
\item (\lref{Formalism.L2.SG.Pf.NcS.SplitSk.SplitSk_ident}) For measurable $f\ge0$ on $G(F\otimes\R)$ and a section
$s$ of the archimedean orbit map of $P$,
\[
\Bigl(\int f\,d\mu_{G,\infty}\Bigr)\mu_{G,f}(K_f)=\mu_{H,f}(K_f\cap H_P(\multncAf))\;\omega_f(K_f\cdot P)\int\!\int
f(s(y)h)\,d\mu_{H,\infty}(h)\,d\omega_{\infty,T_P}(y);
\]
this is the disintegration of the measure of $G(\A)$ along $g\mapsto gP$, applied to $f\otimes\one_{K_f}$, with the finite
part computed by (W2).
\item (\lref{Formalism.L2.SG.Pf.NcS.KfOrbPf.finite_orbit_holds}) $K_f\cdot P$ is measurable, and $x\in K_f\cdot P$ if and
only if $x_v\in K_v\cdot P$ for every $v$.
\item (\lref{Formalism.L2.SG.Pf.NcS.BoxVolPf.box_volume_holds}) If $\omega_{\infty,T_P}\ne0$ and $\omega_f$ is as in (W2),
then $\omega_f$ is concentrated on the tuples $x$ with $\mathrm{gram}(x_v)=T_P$ for all $v$, and for every family of
measurable Schwartz--Bruhat sets $C_v$ equal to the standard lattice box for almost all $v$,
$\omega_f(\{x:x_v\in C_v\ \forall v\})=\prod_v\omega_v(C_v\cap\{\mathrm{gram}=T_P\})$.
\end{enumerate}

\begin{lemma}[Finite volume of the orbit]\label{pf:multnc:omfinvol}
Assume the \emph{local Witt theorem} (for every finite $v$, every $a$-tuple over $F_v$ whose Gram matrix is that of a rational
tuple $P$ with non-degenerate Gram matrix is $G(F_v)$-conjugate to $P$) and the \emph{integral Witt theorem almost
everywhere} (for every lattice $L$ and almost all $v$, every tuple of the standard lattice box at $v$ with the Gram matrix of
$P$ is $K_v$-conjugate to $P$). Then, if $\omega_{\infty,T_P}\ne0$ and $\omega_f$ is as in (W2),
$\omega_f(K_f\cdot P)=\prod_v\omega_v(K_v\cdot P)$.
\formallabel{Formalism.L2.SG.Pf.NcS.FinVolSk.omFinVol_of}
\formallabel{Formalism.L2.SG.Pf.NcS.SplitPf.omega_finite_volume_of_witt}
\end{lemma}

\begin{proof}
By the two Witt theorems (\lref{Formalism.L2.SG.Pf.NcS.LocSBPf.locSB_of}) there are measurable Schwartz--Bruhat sets
$C_v$ with $C_v\cap\{\mathrm{gram}=T_P\}=K_v\cdot P$ for all $v$ and $C_v$ equal to the standard lattice box for almost all
$v$: by the local Witt theorem $K_v\cdot P$ is the intersection of a compact open set with the fibre, and by the integral
Witt theorem the standard box can be taken for almost all $v$. By (W5), $\omega_f$-almost every $x$ has
$\mathrm{gram}(x_v)=T_P$ for all $v$; for such $x$, by (W4), $x\in K_f\cdot P$ if and only if $x_v\in C_v$ for all $v$. Hence
$\omega_f(K_f\cdot P)=\omega_f(\{x:x_v\in C_v\ \forall v\})=\prod_v\omega_v(C_v\cap\{\mathrm{gram}=T_P\})=
\prod_v\omega_v(K_v\cdot P)$ by (W5).
\end{proof}

\begin{theorem}[The splitting identity]\label{pf:multnc:splitwitt}
The local Witt theorem and the integral Witt theorem almost everywhere imply \multncst{Split}.
\formallabel{Formalism.L2.SG.Pf.NcS.SplitSk.splitOm_of}
\formallabel{Formalism.L2.SG.Pf.NcS.SplitPf.omega_splitting_of_witt}
\end{theorem}

\begin{proof}
Let the data of \multncst{Split} be given, with $0<\mu_{G,f}(K_f)<\infty$, and put $\Pi=\prod_v\omega_v(K_v\cdot P)$ (a real
number $\ge0$). For every choice of archimedean self-dual measures and fibre family, (W2) provides $\omega_f$, and by (W3),
the core identity
\[
\Bigl(\int f\,d\mu_{G,\infty}\Bigr)\mu_{G,f}(K_f)=\mu_{H,f}(K_f\cap H_P(\multncAf))\,\omega_f(K_f\cdot P)\int\!\int f(s(y)h)\,
d\mu_{H,\infty}(h)\,d\omega_{\infty,T_P}(y)
\]
holds. For $f=1$ the left-hand side is non-zero ($\mu_{G,\infty}$ is a Haar measure and $\mu_{G,f}(K_f)>0$), hence
$\omega_{\infty,T_P}\ne0$, and by \cref{pf:multnc:omfinvol}, $\omega_f(K_f\cdot P)=\Pi$; so the core identity holds with
$\omega_f(K_f\cdot P)$ replaced by $\Pi$, and, for $f=1$ again, $\Pi>0$. Put
\[
a'=\frac{\mu_{H,f}(K_f\cap H_P(\multncAf))\,\Pi}{\mu_{G,f}(K_f)},
\]
which is positive and finite by (W1). Dividing the core identity by $\mu_{G,f}(K_f)$ gives (i) of \multncst{Split}, for every
choice of measures and sections, and (ii) is the definition of $a'$.
\end{proof}

\subsubsection*{Glued lattices come from row lattices}

We use the notation of \cref{pf:multnc:rowmass}: $F_0$ is the flag $(e_1,\dots,e_{r_S};z)$ of the row construction in the
hermitian space $V$ of $D$, $\mathsf X$ the span of the $e_a$, $V_1=\mathsf X^\perp/\mathsf X$ with the induced form,
$\bar z$ the image of $z$, and $M(\Lambda)$ the image of $\Lambda\cap\mathsf X^\perp$ in $V_1$. A \emph{row lattice} for
$F_0$ is an $\OO_E$-lattice $\Lambda\subset V$ in the genus of the base lattice $L_0$, unimodular, which is a lift of the flag
($z\in\Lambda$, primitive in the associated subquotient), has the prescribed local form at the primes dividing $g_S$, and
satisfies $\Lambda\cap\mathsf X=\mathsf X_S:=\sum_a\mathfrak d_E^{-(r_S-a)}e_a$, where $\mathfrak d_E$ is the different of $F$
extended to $E$ (\formal{Formalism.MainNotions.S.IsRowLat}).

\begin{proposition}[Row surjectivity]\label{pf:multnc:rowsurj}
Every glued lattice $M\subset V_1$ (a unimodular lattice containing $\bar z$ as a primitive vector) is $M(\Lambda)$ for some
row lattice $\Lambda$ for $F_0$.
\formallabel{PI.UnitaryGroupMassFormula.CntFormula.RowsMassW.rowSurj_of}
\formallabel{PI.UnitaryGroupMassFormula.CntFormula.RowsMassW.rowSurj_of_genus}
\formallabel{PI.UnitaryGroupMassFormula.CntFormula.RowsMassW.row_surjective}
\end{proposition}

\begin{proof}
Choose partner vectors $f_1,\dots,f_{r_S}$ with $\langle e_a,f_b\rangle=\delta_{ab}$, $\langle f_a,f_b\rangle=0$ and
$\langle f_a,z\rangle=0$ (\lref{PI.UnitaryGroupMassFormula.CntFormula.RowsMassW.partner_vector_exists}). For $v\in V$
let $\mathrm{pr}(v)=v-\sum_a\langle e_a,v\rangle f_a\in\mathsf X^\perp$ and $\mathfrak a_a=\mathfrak d_E^{-(r_S-a)}$, and let
$\Lambda_M$ be the set of $v\in V$ such that, for all $a\le r_S$, $\overline{\langle e_a,v\rangle}\in\mathfrak a_a^{-1}$ and
$\langle f_a,\mathrm{pr}(v)\rangle\in\mathfrak a_a$, and such that the image of $\mathrm{pr}(v)$ in $V_1$ lies in $M$; it is
an $\OO_E$-submodule. The development proves: $\Lambda_M$ is a lattice
(\lref{PI.UnitaryGroupMassFormula.CntFormula.RowsMassW.lattice_integral}); it is unimodular
(\lref{PI.UnitaryGroupMassFormula.CntFormula.RowsMassW.lamUnimod}); $z\in\Lambda_M$, $\Lambda_M\cap\mathsf X=\mathsf X_S$,
$z$ is primitive in the associated subquotient, and $M(\Lambda_M)=M$
(\lref{PI.UnitaryGroupMassFormula.CntFormula.RowsMassW.lattice_flag_exists}); and $\mathsf X_S$ has the prescribed local
form at every prime (\lref{PI.UnitaryGroupMassFormula.CntFormula.RowsMassW.local_genus_exists}). Finally $\Lambda_M$ and
$L_0$ are unimodular lattices in the same hermitian space over a CM field unramified at all finite places, so they lie in the
same genus (\lref{PI.UnitaryGroupMassFormula.CntFormula.UG.unimodular_one_genus}; locally, unimodular hermitian lattices
of the same rank over an unramified quadratic extension are isometric \cite{Jacobowitz1962}). Hence $\Lambda_M$ is a row
lattice with $M(\Lambda_M)=M$.
\end{proof}

% ===========================================================================
\subsection{Remarks on the formal statements}\label{pf:multnc:sec-remarks}

\begin{remark}\label{pf:multnc:rem-formal}
\begin{enumerate}[label=(\alph*)]
\item The first theorem has the ten published statements (H1)--(H10) as hypotheses. Two of them, (H1) and (H7), are not
used in the proof. Hypothesis (H10) is Green's parametrisation, which is a theorem (\cref{lit:top:green}).
\item In the proof of the elliptic trace identity, one hypothesis of the intermediate statement
(\lref{GlobalTheory.PsCo.Plug.ellTrace_of_middleX}), namely a statement on regular semisimple classes, is not used. In
the proof of the spectral expansion, two hypotheses of the intermediate statement
(\lref{GlobalTheory.PsCo.SpecX.specTrace}) are the true proposition.
\item The existence of pseudo-coefficients $\multncst{PC}$ differs from the classical existence theorem
(\cref{pf:multnc:lit-ds}(c), \cite{ClozelDelorme1984}): it is stated only for odd $r$ and for Harish-Chandra parameters
satisfying the strong gap condition $\lambda_i-\lambda_{i+1}\ge100r^4$, which are satisfied in the application
(\cref{pf:multnc:params}); on the other hand the trace vanishes on every irreducible unitary representation other than
$\pi(\eps,\lambda)$, not only on the tempered ones, and the orbital integrals at elliptic elements are prescribed. The
elliptic expansion uses only that $\mu$ is a restricted product, that the $f_w$ are smooth with compact support and that
$(f_v)$ is a finite family. In \cref{pf:multnc:nonfield} the oddness of the rank and the pseudo-coefficient property of the
$f_w$ are not used.
\item The second theorem has three instance hypotheses (the three facts of (N3)); in the development they are discharged by
global instances, that is, by the theorems cited in (N3), and do not appear as explicit arguments.
\item In the reference covolume identity (N1), the fundamental sets attached to $u$ and to $u'$ are both measured with the
reference measures attached to $u$. The positivity hypotheses of (N4) contain two conjuncts which are the true proposition.
The three real parameters $s_d$, $C$, $T_h$ of the record and several of its fields (those not listed in
\cref{pf:multnc:def-bcf}) play no role in the identity \eqref{pf:multnc:eq-nctiec}.
\item The constant $C_{\mathrm{dis}}(X,u,u')$ is defined by a choice among the positive constants for which an archimedean
disintegration statement holds (and as $0$ if there is none); the development proves that such a constant exists
(\lref{Formalism.L2.SG.Pf.MassDataPos.cooOf_pos}, positivity of $C_{\mathrm{dis}}$), so $C_{\mathrm{dis}}>0$.
\end{enumerate}
\end{remark}

\begin{remark}[Deferred proofs]\label{pf:multnc:rem-deferred}
The following results are used above with their statements, but their formal proofs, which belong to the part of the
development treated in this section, are not converted here: \cref{pf:multnc:cuspmult}, \cref{pf:multnc:specfin},
\cref{pf:multnc:typevec}, \cref{pf:multnc:localiso}, the facts of \cref{pf:multnc:realtori,pf:multnc:transport,pf:multnc:aux},
the statements (S2)--(S4) and (S7)--(S12), the facts (a)--(d) in the proof of \cref{pf:multnc:spectrace}, (i)--(iv) in the
proof of \cref{pf:multnc:separation}, (a)--(c) in the proof of \cref{pf:multnc:elltrace}, (L3), (L6), the Witt theorems, the
glued lattice mass formula with the facts on row lattices cited in \cref{pf:multnc:rowmass}, and the ingredients cited in
\cref{pf:multnc:sec-kpoly,pf:multnc:sec-classint,pf:multnc:sec-ncproofs}: (K1)--(K10), (B1)--(B8), the reduction-theoretic
and local lemmas (a)--(f) of \cref{pf:multnc:regfin}, the class bound, the local statements used in
\cref{pf:multnc:nonfield,pf:multnc:vanish,pf:multnc:hyperzero,pf:multnc:ellexp}, the local volume computations (a)--(e) of
\cref{pf:multnc:omv}, (A1)--(A5), the integration formulas (a)--(c) in the proof of \cref{pf:multnc:jacgauge}, the local
statements of \cref{pf:multnc:placesplit,pf:multnc:rb-jac}, (W1)--(W5) and the lattice statements of
\cref{pf:multnc:rowsurj}. The record relations cited in the proof of \cref{pf:multnc:record}, (L1), (L2) and the row weight
formula are proved elsewhere in the development.
\end{remark}
